\documentclass[11pt]{article}
\usepackage[T1]{fontenc}
\usepackage{lmodern,microtype}
\usepackage[margin=1in]{geometry}
\usepackage{amsmath,amssymb,amsthm,mathtools}
\usepackage{mathrsfs}
\usepackage{booktabs,tabularx,array,enumitem,longtable}
\usepackage[dvipsnames]{xcolor}
\usepackage[colorlinks=true,linkcolor=MidnightBlue,citecolor=MidnightBlue,urlcolor=MidnightBlue]{hyperref}
\usepackage{bookmark}
\setlength{\emergencystretch}{2em}
\setcounter{tocdepth}{1}
\setlist{topsep=4pt,itemsep=3pt,parsep=0pt}
\newtheorem{theorem}{Theorem}[section]
\newtheorem{proposition}[theorem]{Proposition}
\newtheorem{lemma}[theorem]{Lemma}
\newtheorem{corollary}[theorem]{Corollary}
\theoremstyle{definition}
\newtheorem{definition}[theorem]{Definition}
\theoremstyle{remark}
\newtheorem{remark}[theorem]{Remark}
\numberwithin{equation}{section}
\newcommand{\R}{\mathbb R}
\newcommand{\C}{\mathbb C}
\newcommand{\dd}{\,\mathrm d}
\newcommand{\ee}{\mathrm e}
\newcommand{\ii}{\mathrm i}
\newcommand{\calJ}{\mathcal J}
\newcommand{\abs}[1]{\left|#1\right|}
\newcommand{\norm}[1]{\left\lVert#1\right\rVert}
\DeclareMathOperator{\atanh}{atanh}
\DeclareMathOperator{\Log}{Log}
\newcolumntype{Y}{>{\raggedright\arraybackslash}X}
\renewcommand{\arraystretch}{1.12}
\newcommand{\Tr}{\operatorname{Tr}}
\newcommand{\Dom}{\operatorname{Dom}}
\title{Green Potentials and Arithmetic Observations\\in the Newman Flow}
\author{Drafted for Edward Baker}
\date{12 November 2026\\\small Analytical synthesis through GNP-RS-B10, Investigation 09}
\hypersetup{pdftitle={Green Potentials and Arithmetic Observations in the Newman Flow},pdfauthor={Drafted for Edward Baker}}
\begin{document}
\maketitle
\begin{abstract}
We assemble the analytical foundations of the Green-potential approach
to the Newman family, beginning with its complete theta source and the
geometric and operator realizations that motivated it. The source
normalization, heat clock and logarithmic product lead to an exact
positive strip potential. At time zero this potential is a signed
critical-line logarithmic average minus an explicit convergent prime
baseline. We derive acquisition estimates through multiple real zeros,
complete arithmetic and radial error payments, and a contour-deflated
method that separates singular root logarithms from a regular factor.
Scoped obstructions explain why source accuracy and positive
quadratic estimates alone need not produce the required signed bound.
We connect these estimates to Riemann--Siegel expansions in the
critical strip and to their exact Gaussian evolution in Newman time.
Off-line zeros are cancellations of the complete expression; they do
not force degeneracy of the classical phase saddles. The connection
supplies controlled approximants. On the transition scale
\(\lambda=c/\log(X/(4\pi))\), \(1/2\le c\le1\), we derive the
uniform outer prime baseline, the complete-field logarithmic transfer,
and the coherent moment law. Finite Euler references give zero dyadic
density of centers failing the fixed Green margin at \(c=1\), although
each fixed reference has unbounded residuals at arbitrarily high centers.
A further consequence is convergence to zero of every finite power
mean of the Green potential at \(c=1\), obtained from density and
an exponential-moment bound. At \(c=1\), the block-constant prime
cutoff \(P_U=\lfloor \log^2(U/(4\pi))\rfloor\) yields quantitative
residual and Green mean estimates on translated intervals of length
\(U^{17/81+\rho}\le\ell\le U\), for fixed \(0<\rho\le47/162\).
The proof uses exact-cutoff Mellin projection and a cumulative zeta
mean-square theorem. The reference phase estimate is uniform in
the regulator and in \(c\). The every-center residual estimate
remains open.
A separate theorem obstructs the specified fixed-reference signed
tangent bound. We retain the collision and localized-correlation developments as
parallel inputs, and distinguish their quantitative consumers. Numerical
certificates appear only as supporting evidence with their original
finite scope. The final section states the outstanding analytical tasks.
This is a synthesis of inherited proofs and explicitly conditional
methods, with no new claim of zero exclusion or independent review.
\end{abstract}
\clearpage
\tableofcontents
\clearpage
\section{The question and the analytical structure}\label{sec:intro}

The Green approach asks for a signed upper estimate on an observable
whose zero representation is positive. Its usefulness rests on keeping
these two descriptions of the same quantity together. Positivity of a
source, regularity of its heat flow and accuracy of an approximation
are valuable inputs, but none by itself estimates that signed value.

This manuscript continues the Green-potential investigation and integrates
the relevant foundations from both of its parent manuscripts
\cite{GNPBackground,GNPOperator,OGParent}. It is organized around the chain
\[
 \text{complete theta source}\ \longrightarrow\ H_\lambda
 \ \longrightarrow\ \log|H_\lambda|\ \longrightarrow\ J_\lambda
 \ \longrightarrow\ \text{a quantitative defect comparison}.
\]
The spatial variable is \(z\), Newman time is \(\lambda\), and forward
heat time is \(\tau=\lambda_0-\lambda\). The half-line cosine convention
is fixed in \eqref{eq:H}; in particular \(H_0(z)=\xi((1-\ii z)/2)/8\).
The original source variable and the observed spatial variable are
distinct. A heat operation on one need not induce the same operation on
the other.

Sections~\ref{sec:source}--\ref{sec:log-collision} derive the source,
geometric and operator dictionaries and explain why collision singularities
belong to root and logarithmic coordinates while the source orbit stays
regular. Sections~\ref{sec:green}--\ref{sec:prime} identify the complete
Green potential and its exact time-zero arithmetic baseline. The boundary
logarithm, finite-source and deflation arguments in
Sections~\ref{sec:boundary-log}--\ref{sec:deflation} give distinct ways of
acquiring that observable. Section~\ref{sec:RS} explains how classical and
heat-evolved Riemann--Siegel formulas enter those estimates.
Sections~\ref{sec:transition}--\ref{sec:euler-reference} then develop
the child's results through E06: the uniform outer reduction, the
limits of quadratic acquisition, and the Euler-reference comparison.
Theorem~\ref{thm:b10-averaged} extends the residual mean to a growing
block-constant prime cutoff on shorter translated intervals at \(c=1\).
Corollary~\ref{cor:b10-green} gives the corresponding Green mean and
fixed-threshold exceptional proportion. Each deduction keeps the
physical time fixed inside the spatial integral at each center.

The conditioning example in Section~\ref{sec:conditioning} concerns a
specified class containing noncanonical translated sources. The
short-mollifier theorem in Section~\ref{sec:quadratic} concerns one
polynomial shared by a whole interval of centers. Their hypotheses and
quantifiers are part of the conclusions. Neither is a general obstruction
to theta-specific signed arithmetic estimates. Section~\ref{sec:parallel}
keeps localized correlations and collision constraints as parallel
investigations, identifying exactly which inputs they could supply.

There are four different levels of assertion throughout this paper:
an exact representation, an estimate of acquisition error, a conditional
algorithm with unpaid data requirements, and an upper value small enough
for an exclusion. The first three do not imply the fourth. The one-defect
threshold \eqref{eq:margin} provides a numerical scale against which an
estimate can be judged; a reformulation of the all-height target is
identified once at its first use. The premise screen in
Section~\ref{sec:evidence} applies the repository's collision controls
\cite{OGControls}; the source/evidence map accompanying this manuscript
points to the result and avoidance indexes and the exact inherited files.

The author-facing purpose is a coherent entry point for further Green
research. The child has selected a transition-scale consumer and proved
both acquisition estimates and scoped obstructions. Its outstanding
issue is the arithmetic value at exceptional centers, rather than
selection of another representation. The final section distinguishes
that issue from the parent's time-zero question and from the additional
uniformity needed for other defect depths or time bands. This synthesis
does not assert a new theorem of canonical smallness.

\section{The complete theta source and its heat evolution}\label{sec:source}

\subsection{Reflection, the endpoint and the canonical normalization}

Let
\[
 \vartheta(t)=\sum_{n\in\mathbb Z}\ee^{-\pi n^2t},\qquad
 p(u)=\frac{\ee^u}{2}\{\vartheta(\ee^{4u})-1\}
     =\ee^u\sum_{n\ge1}\ee^{-\pi n^2\ee^{4u}}.
\]
Jacobi's transformation
\(\vartheta(t)=t^{-1/2}\vartheta(1/t)\)
\cite{DLMFThetaModular} gives
\begin{equation}\label{eq:primitive-reflection}
 p(-u)-p(u)=\sinh u,\qquad p'(0)=-\frac12.
\end{equation}
Define
\begin{equation}\label{eq:source}
 \Phi(u)=\frac{p''(u)-p(u)}8
 =\ee^u\sum_{n\ge1}(2r_n(u)^2-3r_n(u))\ee^{-r_n(u)},
 \qquad r_n(u)=\pi n^2\ee^{4u}.
\end{equation}
Applying \(\partial_u^2-1\) to \eqref{eq:primitive-reflection}
proves that \(\Phi\) is even. The complete series is smooth, and
each summand is positive when \(u\ge0\), since \(r_n\ge\pi>3/2\).
Thus \(\Phi>0\) on \(\R\). Reflection is a property of the complete
sum; an evenly reflected finite prefix need not have matching odd
derivatives at zero.

With \(W_\lambda(u)=\ee^{\lambda u^2}\Phi(u)\), put
\begin{equation}\label{eq:H}
 H_\lambda(z)=\frac12\int_\R W_\lambda(u)\ee^{\ii zu}\dd u
            =\int_0^\infty W_\lambda(u)\cos(zu)\dd u.
\end{equation}
There is no division by \(H_\lambda(0)\).
For \(J^{\rm prim}(z)=\int_0^\infty p(u)\cos(zu)\dd u\),
two integrations by parts, using the endpoint in
\eqref{eq:primitive-reflection}, give
\begin{equation}\label{eq:primitive-H}
 8H_0(z)=\frac12-(1+z^2)J^{\rm prim}(z).
\end{equation}
For completeness, splitting the theta Mellin integral at \(t=1\)
and using Jacobi reflection gives
\[
 \xi(s)=\frac12+\frac{s(s-1)}2
 \int_1^\infty \frac{\vartheta(t)-1}{2}
       \{t^{s/2-1}+t^{(1-s)/2-1}\}\dd t.
\]
The integral is entire in \(s\). Substituting \(t=\ee^{4u}\)
and then \(s=(1-\ii z)/2\) turns the integral into
\(8J^{\rm prim}(z)\), while \(4s(s-1)=-(1+z^2)\).
Equation~\eqref{eq:primitive-H} therefore proves
\begin{equation}\label{eq:xi}
 H_0(z)=\frac18\xi(s(z)),\quad
 s(z)=\frac12-\frac{\ii z}2,\quad
 \xi(s)=\frac12s(s-1)\pi^{-s/2}\Gamma(s/2)\zeta(s).
\end{equation}
The reflected argument gives the same field
\cite{DLMFGreenPrime}. The constant \(1/2\) in
\eqref{eq:primitive-H} is indispensable.

\subsection{Complete domination, heat evolution and the product}

Set \(\rho(u)=\pi\ee^{4|u|}\), and recursively define
\[
 P_0(r)=2r^2-3r,\qquad
 P_{j+1}(r)=(1-4r)P_j(r)+4rP_j'(r).
\]
Write \(a(P)\) for the sum of the absolute coefficients of a polynomial,
and put
\[
 C_d(N)=\sum_{n>N}n^{2d}\ee^{-\pi(n^2-1)},\qquad
 B_j=a(P_j)C_{j+2}(0).
\]
Each \(C_d(N)\) is finite and tends to zero with \(N\).

\begin{lemma}[Complete source bound]\label{lem:source}
For each fixed integer \(j\ge0\),
\[
 |\Phi^{(j)}(u)|\le B_j\ee^{|u|}\rho(u)^{j+2}\ee^{-\rho(u)}.
\]
On \(u\ge0\), the differentiated arithmetic tail satisfies the same
bound with \(B_j\) replaced by \(a(P_j)C_{j+2}(N)\).
\end{lemma}
\begin{proof}
Differentiating a summand gives
\(\ee^uP_j(n^2\rho)\ee^{-n^2\rho}\).
Since \(\deg P_j=j+2\) and \(\rho\ge\pi>1\),
\[
 |P_j(n^2\rho)|\le a(P_j)n^{2j+4}\rho^{j+2},\qquad
 \ee^{-n^2\rho}\le\ee^{-\rho}\ee^{-\pi(n^2-1)}.
\]
Summation proves the half-line estimate and differentiated local
convergence. Complete smooth reflection gives the full-line estimate.
\end{proof}

For fixed \(L,b\ge0\) and integers \(k,d\ge0\), the integral
\[
 \int_\R |u|^k\ee^{Lu^2+(b+1)|u|}
             \rho(u)^d\ee^{-\rho(u)}\dd u
\]
is finite: substitution \(t=\pi\ee^{4|u|}\) leaves a product of
powers of \(t,\log(t/\pi)\), and
\(\exp\{(L/16)\log^2(t/\pi)-t\}\).
Consequently \(H\) is jointly entire in \((\lambda,z)\), and all
fixed derivatives on compact parameter sets pass under the integral.
In particular,
\begin{equation}\label{eq:heat}
 \partial_\lambda H_\lambda=-\partial_z^2H_\lambda.
\end{equation}
The opposite clock \(h(\tau,z)=H_{\lambda_0-\tau}(z)\) satisfies
\(h_\tau=h_{zz}\). No uniform assertion in a growing derivative
order follows from these fixed-order bounds.

We retain the canonical strip input
\begin{equation}\label{eq:zero-strip}
 |\Im z_j(\lambda)|\le1\qquad(0\le\lambda\le1/2)
\end{equation}
from the parent development \cite{OGParent,GreenStripProof}.
This is strip control, not an all-real-zero assumption.
The source bound also gives, uniformly in this compact time range,
\(\log\max_{|z|\le r}|H_\lambda(z)|=O(r\log(r+2))\).
One can obtain this bound directly by splitting the source integral
at a constant multiple of \(\log(r+2)\); beyond that point the
double exponential dominates \(r|u|+u^2/2+O(|u|)\).
Positivity gives a common lower bound for \(H_\lambda(0)\).
Jensen's formula on a disk of radius \(2r\) therefore gives a
full origin-centered count
\begin{equation}\label{eq:count}
 N_\lambda(r)\le Cr\log(r+2),\qquad r\ge1.
\end{equation}
The even Hadamard product is consequently
\begin{equation}\label{eq:product}
 H_\lambda(z)=H_\lambda(0)
 \prod_{\zeta\in\mathcal R_\lambda}(1-z^2/\zeta^2).
\end{equation}
Here \(\mathcal R_\lambda\) contains one representative of each
opposite pair, with multiplicity. The paired product converges
locally uniformly because \(\sum|\zeta|^{-2}<\infty\). The order
bound permits at most a linear exponential factor; evenness removes
its linear term, and evaluation at zero fixes the constant.
Thus no unspecified exponential factor remains.


\section{Geometric origins and the calibrated heat realization}
\label{sec:geometry}

The preceding source can be described by geometric traces, scalar
potentials and a prepared operator orbit. These descriptions explain the
origin of the investigation. We derive the identities needed for the
Green thread while retaining their preparation and observation assumptions
\cite{GNPBackground,OperatorProof,OperatorAssessment}.

\subsection{The cylinder and the multiplicity-cancelling trace}
Let \(D\) be the Dirac operator on the unit two-sphere. Its eigenvalues
are \(\pm n\), each of multiplicity \(2n\), \(n\ge1\)
\cite{CamporesiHiguchi1996}. On the cylinder
\(\R_v\times S^2_{1/2}\), with the radial Clifford matrix anticommuting
with the angular operator, the geometric square is
\begin{equation}\label{eq:cylinder-square}
 L_0=-\partial_v^2+4A,\qquad A=D^2.
\end{equation}
The scalar curvature is \(8\), so the Lichnerowicz curvature term is
\(2\); it is already contained in this geometric square. The
\(A=n^2\) shell has total rank \(4n\). Its radial Green kernel is
\[
 g_n(v,w)=\frac{\ee^{-2n|v-w|}}{4n}.
\]
Indeed it solves \((-\partial_v^2+4n^2)g_n=\delta_w\): the derivative
jump at \(w\) is \(-1\). Thus the coincident kernel induces the angular
weight \(W=(4\sqrt A)^{-1}\), and
\begin{equation}\label{eq:geometric-trace}
 Q(\beta)=\Tr(W\ee^{-\beta A})=\sum_{n\ge1}\ee^{-\beta n^2}.
\end{equation}
All traces at \(\beta>0\) converge absolutely. The factor \(4n\) from
the shell cancels the diagonal Green factor \(1/(4n)\).

Supply the scan \(\beta(u)=\pi\ee^{4u}\). Then
\(p(u)=\ee^u Q(\beta(u))\), and differentiating twice gives
\begin{equation}\label{eq:trace-source}
 \Phi(u)=\ee^u Q(\beta)
 \Tr\bigl[\rho_\beta(2\beta^2A^2-3\beta A)\bigr],\qquad
 \rho_\beta=\frac{W\ee^{-\beta A}}{Q(\beta)}.
\end{equation}
The differentiation is justified by Gaussian angular summability on
compact \(u\) sets. This proves the positive-half-line source formula
and, together with theta reflection, the full source. The scan is a
specified relation between variables: translation of the radial Green
probe does not generate it, since \(g_n(v,v)\) is constant in \(v\).

\subsection{A scalar potential and its transported source}
Write \(\Theta(R)=1+2Q(\pi R^2)\). Jacobi reflection says
\(\Theta(R)=R^{-1}\Theta(1/R)\). The scalar potential
\(V(R)=1+R^{-1}-\Theta(R)\) gives the even function
\begin{equation}\label{eq:scalar-G0}
 \mathscr G_0(u)=\tfrac12\ee^u V(\ee^{2u})
 =\tfrac12\ee^{-u}-p(u)\quad(u\ge0).
\end{equation}
Let \(L=1-\partial_u^2\). The reflection endpoint \(p'(0)=-1/2\)
makes the derivative of \eqref{eq:scalar-G0} zero at the origin, so its
even extension has no distributional delta there. Hence
\(L\mathscr G_0=8\Phi\). More generally define
\begin{equation}\label{eq:scalar-G-lambda}
 \mathscr G_\lambda=8L^{-1}W_\lambda
 =4\int_\R \ee^{-|u-w|}W_\lambda(w)\dd w.
\end{equation}
The positive convolution and the complete source bounds justify every
fixed derivative and yield
\(\mathscr G_\lambda(u)\sim c_\lambda\ee^{-|u|}\), where
\(c_\lambda=8H_\lambda(\ii)>0\) and \(c_0=1/2\).
For example, as \(u\to+\infty\), the part \(w\le u\) equals
\(4\ee^{-u}\int_{-\infty}^u\ee^w W_\lambda(w)\dd w\); the complementary
tail is negligible by double-exponential decay. Evenness proves the
other end.

Fourier transformation of \(L\mathscr G_\lambda=8W_\lambda\) gives
\begin{equation}\label{eq:scalar-fourier}
 16H_\lambda(z)=(1+z^2)\int_\R\mathscr G_\lambda(u)\ee^{\ii zu}\dd u,
 \qquad |\Im z|<1.
\end{equation}
The raw integral is only asserted in this strip. At \(z=\pm\ii\) one
must use the entire source definition of \(H_\lambda\), rather than
multiply a vanishing polynomial by a divergent integral.

If \(X\) denotes multiplication by \(u\), source transport induces
\begin{equation}\label{eq:scalar-transport}
 \partial_\lambda\mathscr G_\lambda=\mathcal B\mathscr G_\lambda,
 \quad \mathcal B=L^{-1}X^2L,
 \quad \mathcal Bf=u^2f+L^{-1}(4uf'+2f).
\end{equation}
The commutator identity follows by expanding \(L(u^2f)\).
On Schwartz functions \(\mathcal B^j=L^{-1}X^{2j}L\); the displayed
prepared family supplies its own domains through \eqref{eq:scalar-G-lambda}.
Multiplying \(\mathscr G_0\) by \(\ee^{\lambda u^2}\) is a different
operation and is not integrable for \(\lambda>0\). This distinction is
essential when interpreting a geometric potential as a Newman source.

\subsection{A fixed operator with the correct heat clock}
The calibrated operator is
\begin{equation}\label{eq:calibrated-K}
 K=L_0-4I=-\partial_v^2+4(A-I).
\end{equation}
Under the unitary radial Fourier transform its shell multiplier is
\(m_n(p)=p^2+4(n^2-1)\), with maximal domain
\[
 \Dom K=\left\{\psi:\sum_n\int_\R m_n(p)^2
             \|\widehat\psi_n(p)\|^2\dd p<\infty\right\}.
\]
Real nonnegative multiplication proves self-adjointness and positivity;
the spectrum is purely absolutely continuous and equals \([0,\infty)\).
The value zero is not an eigenvalue. The semigroup \(S_\tau=\ee^{-\tau K}\)
has norm one, and for \(j\ge1,\tau>0\),
\(\|K^jS_\tau\|\le(j/(\ee\tau))^j\), by maximizing
\(r^j\ee^{-\tau r}\) over \(r\ge0\).

Fix a unit vector \(\chi\) in the first angular shell. Define
\(\widehat{E_\chi f}(p)=f(p)\chi\) and
\(R_\chi\psi(x)=\sqrt{\pi/2}\langle\chi,\psi(x)\rangle\).
The first is an isometry and the second is bounded between the relevant
\(L^2\) spaces. On their natural domains,
\begin{equation}\label{eq:operator-intertwine}
 KE_\chi=E_\chi X^2,\qquad R_\chi K=-\partial_x^2R_\chi.
\end{equation}
The complete prepared amplitude \(W_{\lambda_0}\) therefore satisfies
\begin{equation}\label{eq:operator-response}
 R_\chi S_\tau E_\chi W_{\lambda_0}=H_{\lambda_0-\tau}=h(\tau,\cdot).
\end{equation}
To verify the constant, inverse Fourier transformation gives
\(\sqrt{\pi/2}/\sqrt{2\pi}=1/2\), exactly the coefficient in
\eqref{eq:H}. Evenness removes the Fourier-sign convention. The source
bounds place the prepared orbit in every fixed graph domain. Its complex
frequency continuation and pointwise derivatives follow from weighted
source integrals, not bounded point evaluation on arbitrary \(L^2\).

Using \(L_0\) instead would produce \(\ee^{-4\tau}h\). The calibration
is material. In fact a fixed bounded linear observation of
\(\ee^{-\tau L_0}\psi\), with fixed \(\psi\), decays exponentially in
\(L^2\), whereas
\[
 \|H_{\lambda_0-\tau}\|_2
 \sim\Phi(0)(\pi/2)^{3/4}\tau^{-1/4}.
\]
This follows from Plancherel and scaling
\(\int\ee^{-2\tau u^2}|W_{\lambda_0}(u)|^2\dd u\).
It excludes this particular uncalibrated bounded-observation realization
for all \(\tau\), not all alternative preparations or observations.
Likewise reverse heat is unbounded on ambient states, while the
canonical source itself is entire in \(\lambda\).

The representation preserves its source-accuracy burden exactly:
\begin{equation}\label{eq:operator-error}
 \|R_\chi S_\tau E_\chi(f-g)\|_2
 =\sqrt{\pi/2}\,\|\ee^{-\tau u^2}(f-g)\|_2.
\end{equation}
No arithmetic cancellation estimate is supplied merely by changing
coordinates to this operator realization.

\subsection{Energy remains regular at a collision}
For the prepared orbit put
\(N=\int|W_{\lambda_0-\tau}|^2=(2/\pi)\int h^2\) and
\(T=\int u^2|W_{\lambda_0-\tau}|^2=(2/\pi)\int h_x^2\).
Differentiation and Plancherel give
\begin{equation}\label{eq:energy-laws}
 N'=-2T,\qquad T'=-2\|K S_\tau E_\chi W_{\lambda_0}\|_2^2,
 \qquad (T/N)'=-2\operatorname{Var}_{|W|^2/N}(u^2).
\end{equation}
The linear source probability \(W/\int W\) and the spectral probability
\(|W|^2/N\) are distinct; the latter produces the factor two above.
Locally, \((h^2)_\tau=(h^2)_{xx}-2h_x^2\). At a double zero with
\(h=h_x=0\), \(h_{xx}=a\ne0\), one has
\((h^2)_\tau=0\) and \((h^2)_{\tau\tau}=2a^2\).
The regular energy laws are compatible with a change in the real zero
configuration and impose no collision barrier.

\subsection{The other parent realizations and their scope}
For completeness, the shared parent also supplied two exact observation
models \cite{TNDDirect,TNDAux}. They are not new inputs to a Green sign
estimate. In the direct model prepare
\(f_\lambda=\sqrt{W_\lambda/Z_\lambda}\), \(Z_\lambda=2H_\lambda(0)\),
and supply the measured quadrature. With
\(U(t)=\ee^{-\ii\omega tL_0}\), \(Q_t=U(t)XU(t)^*\), the operations
\[
 K_j(t,s)=U(t)\ee^{-\gamma(t-s)X^2/2}
       \frac{(\sqrt{\gamma(t-s)}X)^j}{\sqrt{j!}}U(s)^*
\]
are complete since the positive sum \(\sum_jK_j^*K_j\) is the identity
by the scalar exponential series. The no-count amplitude multiplies
\(f_{1/2}\) by \(\ee^{-\gamma t u^2/2}\), giving
\(\lambda=1/2-\gamma t\), \(0\le t\le1/\gamma\). The effects
\((I\pm\cos(kQ_t))/2\) then obey
\[
 q_++q_-=H_\lambda(0)/H_{1/2}(0),\qquad
 H_\lambda(k)=H_{1/2}(0)(q_+-q_-).
\]
These are joint probabilities including failed conditioning; arbitrary
complex \(k\) is not claimed to define a bounded effect.

In the auxiliary model the supplied amplitude is instead
\(W_\lambda(v)B\), where \(BB^*=W\ee^{-\pi A}/Q(\pi)\) and
\(\|B\|_{\rm HS}=1\). Apply \(\ee^{-\tau L_0}\) on the physical side
and pair the Fourier amplitude with \(B\). Commuting radial and angular
heat factors and summing the complete angular trace gives
\begin{equation}\label{eq:auxiliary-response}
 \mathcal O_{\lambda,\tau}(z)
 =\ee^{-\tau z^2}\frac{Q(\pi+4\tau)}{Q(\pi)}H_\lambda(z).
\end{equation}
The multiplier never vanishes, so this response preserves zeros and
multiplicities. Recovery on real \(z=k\) amplifies errors by
\(\ee^{\tau k^2}Q(\pi)/Q(\pi+4\tau)\), unbounded in \(k\) for
\(\tau>0\). Proper time in this model is not Newman time.
The graph-domain and weighted-integral proofs for both models remain
in the parent sources; their loaders, quadrature access and calibrated
detector laws are supplied assumptions. Their exact observation formulas
do not generate a missing source-specific inequality.

\section{Logarithmic fields and collision coordinates}\label{sec:log-collision}

\subsection{The logarithmic field and its denominator}

Let \(\nu_\lambda=\sum_j\delta_{z_j(\lambda)}\), with multiplicities,
and use \(\partial_z=(\partial_x-\ii\partial_y)/2\).
Local factorization \(H=(z-z_j)^m h\), with \(h\ne0\), proves
\begin{equation}\label{eq:distributions}
 \Delta\log|H_\lambda|=2\pi\nu_\lambda,\qquad
 R_\lambda=\frac{H_\lambda'}{H_\lambda}
           =2\partial_z\log|H_\lambda|,\qquad
 \partial_{\bar z}R_\lambda=\pi\nu_\lambda .
\end{equation}
These are distributional identities; \(\log|z|\) and \(1/z\) are
locally integrable, with the stated delta masses obtained by integrating
around a small circle. A multiplicity-\(m\) zero gives residue \(m\).
Differentiating \eqref{eq:product} away from its zeros gives
\[
 R_\lambda(z)=\sum_{\zeta\in\mathcal R_\lambda}
                    \frac{2z}{z^2-\zeta^2}.
\]
Pairing is part of the convergence prescription.

On zero-free space-time patches, logarithmic differentiation of
\eqref{eq:heat} yields
\begin{equation}\label{eq:burgers}
 \partial_\lambda R=-R_{zz}-2RR_z.
\end{equation}
This local equation is not assigned an additional unproved global
meaning at its poles. Likewise, solving only the first or last
distributional equation in \eqref{eq:distributions} leaves a harmonic
or holomorphic summand. The complete field and its product fix it.

Source approximation alone does not control division. If
\(|H-\widehat H|\le\varepsilon\),
\(|H'-\widehat H'|\le\varepsilon_1\), and
\(b=|\widehat H|-\varepsilon>0\), then
\begin{equation}\label{eq:ratio-error}
 \left|R-\frac{\widehat H'}{\widehat H}\right|
 \le\frac{\varepsilon_1}{b}
       +\frac{|\widehat H'|\varepsilon}{b|\widehat H|}.
\end{equation}
This follows by adding and subtracting
\(\widehat H'/H\). The positive denominator margin is a separate
requirement. It fails on a real interval containing a zero, which
motivates the integrated logarithmic method developed later.


\subsection{Local collision geometry}\label{sec:double}
The following local derivation is inherited from the operator parent\cite{OperatorCollision,GNPOperator}.

\paragraph{Event assumptions and a finite heat expansion}
Fix an interior point \((\tau_*,x_*)\) such that
\begin{equation}\label{eq:double}
 h=h_x=0,\qquad a=h_{xx}\ne0,
 \qquad b=h_{xxx},\quad c=h_{xxxx},
\end{equation}
with derivatives evaluated at that point unless stated otherwise.
This is a hypothetical event of the canonical family; its existence is
not assumed elsewhere. Since \(h(\tau,0)>0\), its position is nonzero.
Evenness gives a reflected event at \(-x_*\); the analysis is local
near either one.

Set \(s=\tau-\tau_*\), \(y=x-x_*\).
The heat equation replaces each time derivative by two spatial
derivatives. Taylor expansion therefore gives
\begin{equation}\label{eq:heat-four-jet}
 h(\tau_*+s,x_*+y)
 =\frac a2(y^2+2s)+\frac b6(y^3+6sy)
  +\frac c{24}(y^4+12sy^2+12s^2)+\mathcal R(s,y).
\end{equation}
For fixed compact sets in \((S,Y)\),
\(\mathcal R(\epsilon^2S,\epsilon Y)=O(\epsilon^5)\).
This notation assigns weight two to \(s\) and weight one to \(y\).
To justify it, take a sufficiently high finite Taylor expansion in
\((s,y)\). Terms of weight below two vanish by \eqref{eq:double};
the terms of weights two, three and four are exactly those displayed.
Every omitted monomial has weight at least five, and the ordinary
Taylor remainder has the same bound after rescaling. The source bounds
justify the needed finite derivatives. The constants depend on the event
and the chosen compact set.

\paragraph{The regular critical curve}
\begin{proposition}[Critical-value crossing]\label{prop:critical}
There is a unique analytic local curve \(x_c(\tau)\) with
\(x_c(\tau_*)=x_*\) and \(h_x(\tau,x_c(\tau))=0\).
For \(v(\tau)=h(\tau,x_c(\tau))\),
\begin{equation}\label{eq:critical-exact}
 x_c'=-\frac{h_{xxx}}{h_{xx}},\qquad v'=h_{xx}
 \quad\hbox{along the curve}.
\end{equation}
Its expansions are
\begin{align}
 x_c(\tau_*+s)&=x_*-\frac ba s+O(s^2),\label{eq:critical-place}\\
 v(\tau_*+s)&=as+\frac12\left(c-\frac{b^2}{a}\right)s^2+O(s^3).
 \label{eq:critical-value}
\end{align}
\end{proposition}
\begin{proof}
Apply the analytic implicit-function theorem to \(h_x\), whose spatial
derivative at the event is \(a\ne0\). Differentiating \(h_x=0\) along
the curve gives \(h_{xxx}+h_{xx}x_c'=0\). Differentiating \(v\) gives
\(v'=h_\tau+h_xx_c'=h_{xx}\). A second differentiation at the event
yields \(v''=h_{xxxx}+h_{xxx}x_c'=c-b^2/a\), proving the expansions.
\end{proof}

Thus the extremum value crosses zero transversely, with finite velocity
\(a\). If \(a>0\), it is a local minimum which rises through zero; if
\(a<0\), it is a local maximum which falls through zero. This description
remains regular at the event.

\paragraph{The two-root cluster}
\begin{theorem}[Square-root splitting and drift]\label{thm:splitting}
In a sufficiently small fixed complex disk about \(x_*\), the field
has exactly two zeros counted with multiplicity for all sufficiently
small real \(|s|\). They satisfy
\begin{equation}\label{eq:roots}
 z_\pm(\tau_*+s)=x_*-\frac{2b}{3a}s+
 \begin{cases}
  \displaystyle\pm\sqrt{-2s}+O(|s|^{3/2}),&s<0,\\[2pt]
  \displaystyle\pm\ii\sqrt{2s}+O(s^{3/2}),&s>0.
 \end{cases}
\end{equation}
They are real and simple for \(s<0\), and a nonreal conjugate pair for
\(s>0\). Their midpoint is analytic in \(s\).
\end{theorem}
\begin{proof}
At \(s=0\) the zero has exact order two. Choose a disk containing no
other zero and having a nonzero boundary. Uniform analytic convergence
on that boundary preserves the zero count for small complex \(s\).
Put \(s=-t^2/2\), \(y=tY\). The quotient
\[
 t^{-2}h(\tau_*-t^2/2,x_*+tY)
\]
extends holomorphically to \(t=0\), with value
\(a(Y^2-1)/2\). The analytic implicit-function theorem at its simple
roots \(Y=\pm1\) supplies two analytic branches. Substitution in
\eqref{eq:heat-four-jet} gives
\begin{equation}\label{eq:t-roots}
 y_\pm(t)=\pm t+\frac{b}{3a}t^2+O(t^3).
\end{equation}
For real \(t\), the branches are real by conjugation and uniqueness.
For \(s>0\), take \(t=\ii\sqrt{2s}\); the leading imaginary parts are
nonzero and the roots are conjugate. Replacing \(t\) by \(-t\)
interchanges the two roots, so their sum and product are even analytic
functions of \(t\), hence analytic in \(s\). This also proves the
midpoint assertion and \eqref{eq:roots}.
\end{proof}

In increasing imaginary time, the real roots approach, meet and leave
the real axis. The direction is reversed in Newman \(\lambda\).
The total local complex multiplicity stays two. This is the standard
local forward-parabolic pattern \cite{Angenent1991}; a global zero count
on the whole line is not inferred from the local statement.



\paragraph{A regular squared gap and the surrounding factor}
The symmetric functions constructed above define a monic quadratic
factor. Local analytic factorization gives
\begin{equation}\label{eq:local-factor}
 h(\tau,x)=(x-z_-(\tau))(x-z_+(\tau))U(\tau,x),
\end{equation}
where \(U\) is analytic and nonzero near the event. One may obtain the
quadratic's coefficients from the contour integrals of
\(x^k h_x/h\) around the fixed disk, which are analytic in \(\tau\);
division removes precisely the two enclosed roots. The quotient is a
holomorphic unit, including at \(s=0\). This is the same local
factorization used in the earlier heat-multiplicity analysis
\cite{HeatMultiplicity}. No product over all zeros is required.

For \(s<0\), write \(r_\pm=z_\pm\), \(\delta=r_+-r_->0\), and
\(q=\delta^2\). Let \(g=U_x/U\).

\begin{proposition}[Pair motion and squared gap]\label{prop:gap}
On the real-pair side,
\begin{align}
 r_\pm'&=-\frac{2}{r_\pm-r_\mp}-2g(\tau,r_\pm),\label{eq:pair-motion}\\
 \delta'&=-\frac4\delta-2\bigl(g(\tau,r_+)-g(\tau,r_-)\bigr),\label{eq:delta}\\
 q'&=-8-4q\int_0^1g_x(\tau,r_-+t\delta)\dd t.\label{eq:q-exact}
\end{align}
The squared gap extends analytically through the event and
\begin{equation}\label{eq:q-expansion}
 q(\tau_*+s)=-8s+16B_*s^2+O(s^3),\qquad
 B_*=\frac{c}{6a}-\frac{b^2}{9a^2}.
\end{equation}
The midpoint has velocity \(-2b/(3a)\) at the event.
\end{proposition}
\begin{proof}
At a simple zero \(r\), implicit differentiation gives
\(r'=-h_{xx}/h_x\). Differentiating \eqref{eq:local-factor} twice
at \(r_\pm\) yields
\[
 \frac{h_{xx}}{h_x}(\tau,r_\pm)
 =\frac{2}{r_\pm-r_\mp}+2g(\tau,r_\pm).
\]
Subtracting proves \eqref{eq:delta}; multiplying by \(2\delta\) and
using the fundamental theorem of calculus gives \eqref{eq:q-exact}.
Analyticity of the symmetric functions makes \(q\) analytic.
At \(s=0\),
\[
 U(\tau_*,x_*+y)=\frac a2+\frac b6y+\frac c{24}y^2+O(y^3),
\]
so \(g_*=b/(3a)\) and \(g_{x,*}=B_*\).
The integral in \eqref{eq:q-exact} tends to \(B_*\); its symmetric
expansion is \(B_*+O(s)\), since the midpoint is \(O(s)\) from \(x_*\)
and \(\delta^2=O(s)\). Thus \(q'(0)=-8\) and \(q''(0)=32B_*\),
proving \eqref{eq:q-expansion}. Adding the two root equations cancels
their singular terms and gives the midpoint velocity \(-2g_*\).
\end{proof}

The critical point and root midpoint generally have different velocities:
\(-b/a\) and \(-2b/(3a)\). Individual zero velocities diverge as
\(|s|^{-1/2}\), whereas both collective coordinates remain regular.
The term \(-8\) in \eqref{eq:q-exact} is the universal terminal closing
rate of a separated collision. The regular factor affects drift
and subsequent terms; this is not a uniform estimate on its size.

Rodgers and Tao's zero dynamics and gap law have the corresponding
decomposition with the opposite Newman-time sign
\cite[Section~4, Theorem~4.16 and Lemma~4.18(i)]{RodgersTao2020}.
Their printed theorem uses principal-value sums on the paper's
contradiction interval \(\Lambda<t\leq0\). We do not extend that theorem
to a different interval or to arbitrary nonreal zero configurations.
The local derivation above requires only the analytic factor \(U\).


\paragraph{Amplitude, drift and relative curvature}
The coefficients in \eqref{eq:critical-place}--\eqref{eq:critical-value}
separate naturally into the amplitude \(a\) and two quantities unchanged
by a nonzero constant rescaling of the field:
\begin{equation}\label{eq:beta-gamma}
 \beta=\frac ba,\qquad \gamma=\frac ca-\frac{b^2}{a^2}.
\end{equation}
At the event, \(x_c'=-\beta\), \(v'=a\), and \(v''=a\gamma\).
At the reflected event, \(a,\gamma\) are unchanged and \(\beta\)
changes sign. The critical curve is the curve of stationary points
\(h_x=0\), not a curve of zeros of \(h\).

\begin{proposition}[Curvature transport and an event coordinate]
\label{prop:curvature-transport}
Along a regular critical curve, set
\(A(\tau)=h_{xx}(\tau,x_c(\tau))\ne0\) and \(r_c=-v/A\).
Then
\begin{align}
 A'&=h_{xxxx}-\frac{h_{xxx}^2}{h_{xx}},\qquad
 \frac{d}{d\tau}\log|A|
 =\frac{h_{xxxx}}{h_{xx}}-\frac{h_{xxx}^2}{h_{xx}^2},
 \label{eq:A-transport}\\
 r_c'&=-1-r_c\frac{A'}{A}.\label{eq:rc-transport}
\end{align}
In particular \(r_c(\tau_*+s)=-s+\tfrac12\gamma s^2+O(s^3)\)
and \(r_c\sim q/8\sim\tau_*-\tau\) on the real-pair side.
\end{proposition}
\begin{proof}
Differentiate \(A\) along the curve and use
\(x_c'=-h_{xxx}/h_{xx}\) and \(h_\tau=h_{xx}\).
The quotient rule and \(v'=A\) give \eqref{eq:rc-transport}.
Finally \(v=as+a\gamma s^2/2+O(s^3)\) and
\(A=a+a\gamma s+O(s^2)\), proving the expansion and, with
\eqref{eq:q-expansion}, the comparison.
\end{proof}

Thus \(\gamma\) is the relative rate of change of the extremum's
curvature. The regular event coordinates describe the same crossing;
they do not impose independent restrictions on it.

At the collision slice use the local factor \(U\) in
\eqref{eq:local-factor}, and write
\begin{equation}\label{eq:ell-kappa}
 \ell=\frac{U_x}{U}(\tau_*,x_*),\qquad
 \kappa=\partial_x\!\left(\frac{U_x}{U}\right)(\tau_*,x_*).
\end{equation}
The derivatives of \(h=(x-x_*)^2U\) give
\(U_*=a/2\), \(U_{x,*}=b/6\), \(U_{xx,*}=c/12\). Therefore
\begin{equation}\label{eq:coefficient-dictionary}
 \boxed{\ \beta=3\ell,\qquad c/a=6(\kappa+\ell^2),\qquad
 \gamma=6\kappa-3\ell^2,\qquad B_*=\kappa.\ }
\end{equation}
The same exterior curvature controls the second-order gap correction
and, together with the drift, the second-order critical value.

\paragraph{A conditional real-zero benchmark}
\begin{proposition}[Exterior inverse-square curvature]
\label{prop:LP-coefficients}
Suppose additionally that the collision slice belongs to the real
Laguerre--P\'olya class. If its product has Gaussian factor
\(\ee^{-\alpha z^2}\), \(\alpha\ge0\), and \(\{\rho_j\}\)
are its remaining zeros, with multiplicities, after the two at \(x_*\)
are removed, then
\begin{equation}\label{eq:LP-kappa}
 \kappa=-2\alpha-\sum_j\frac{1}{(x_*-\rho_j)^2}.
\end{equation}
For an even slice with \(x_*\ne0\),
\begin{equation}\label{eq:LP-gamma}
 \kappa\le-\frac{1}{2x_*^2},\qquad
 \gamma\le-\frac{3}{x_*^2}-\frac{\beta^2}{3}<0.
\end{equation}
\end{proposition}
\begin{proof}
The defining canonical product is
\(Cz^m\ee^{-\alpha z^2+d z}
\prod_j(1-z/\rho_j)\ee^{z/\rho_j}\), with real parameters and zeros,
and \(\sum_{\rho_j\ne0}|\rho_j|^{-2}<\infty\)
\cite[Definition~1.1]{Csordas2014PositiveKernels}.
After removing the pair at \(x_*\), take the second logarithmic
derivative. The linear exponential and regularizing linear terms
disappear, giving \eqref{eq:LP-kappa}. The sum converges absolutely
locally off the remaining zeros; only finitely many are nearby and
the distant terms are bounded by a constant times \(|\rho_j|^{-2}\).
Evenness supplies a double zero at \(-x_*\), whose contribution is
\(-2/(2x_*)^2\). All other terms are nonpositive. Equation
\eqref{eq:coefficient-dictionary} proves \eqref{eq:LP-gamma}.
\end{proof}

This is a conditional benchmark for an all-real slice, not an
unconditional theta estimate. The weaker assertion \(\gamma\le0\)
also follows from the derivative Laguerre inequality \(b^2-ac\ge0\).
For \(a>0\), \eqref{eq:LP-gamma} says that a rising minimum
decelerates; its transverse crossing remains possible. The polynomial
\((x^2-r^2)^2\) has \(\beta=3/r\), \(\gamma=-6/r^2\) at \(r\)
and attains the displayed reflected-pair bound. This is a check of
the sharp constant, not a canonical-source example.

The sign changes when exterior zeros need not be real. An exterior
conjugate pair \(d\pm\ii\eta\) contributes
\begin{equation}\label{eq:complex-exterior-curvature}
 -\frac{2((x_*-d)^2-\eta^2)}{((x_*-d)^2+\eta^2)^2}
\end{equation}
to the logarithmic curvature, which may have either sign. A collision
at an all-real boundary and a later collision among a nonreal
background therefore require different estimates. The local theory
neither assumes nor proves that the global Newman threshold is
attained at one finite collision.

Finally, a bounded exterior correction by itself cannot oppose the
terminal closing rate. If the integral in \eqref{eq:q-exact} has
absolute value at most \(C>0\), then
\(q'\le-8+4Cq\le-4\) whenever \(0<q\le1/C\).
If this bound and cluster isolation persist, the pair closes in
at most \(q(\tau_0)/4\) further imaginary time. Preventing entry
into that regime requires information in addition to boundedness
of the regular factor.

\subsection{Finite multiplicity}
\label{sec:multiplicity}


Suppose a real event has exact spatial order \(m\geq2\):
\[
 \partial_x^jh(\tau_*,x_*)=0\quad(0\leq j<m),\qquad
 a_m=\partial_x^mh(\tau_*,x_*)\ne0.
\]
For fixed \(m\), define
\begin{equation}\label{eq:heat-poly}
 p_m(S,Y)=
 \sum_{j=0}^{\lfloor m/2\rfloor}
 \frac{m!}{j!(m-2j)!}S^jY^{m-2j}.
\end{equation}

\begin{theorem}[Finite-multiplicity heat geometry]\label{thm:higher}
Uniformly for \((S,Y)\) in compact sets,
\begin{equation}\label{eq:higher-rescale}
 h(\tau_*+\epsilon^2S,x_*+\epsilon Y)
 =\frac{a_m}{m!}\epsilon^m p_m(S,Y)+O(\epsilon^{m+1}).
\end{equation}
In a sufficiently small fixed complex neighborhood of \(x_*\), the
cluster consists of \(m\) simple real zeros immediately before the event.
Immediately after it, there are no real zeros if \(m\) is even, and
exactly one simple real zero if \(m\) is odd. The remaining zeros form
nonreal conjugate pairs; the total complex multiplicity stays \(m\).
\end{theorem}
\begin{proof}
The coefficient of \(s^jy^k\) in the local Taylor series is
\(\partial_x^{2j+k}h/(j!k!)\), by the heat equation. All terms of
parabolic weight \(2j+k<m\) vanish, and the weight-\(m\) terms give
\eqref{eq:heat-poly}. A finite Taylor remainder, as in
Section~\ref{sec:double}, proves \eqref{eq:higher-rescale}.

For completeness, the real-root property of \(q_m(Y)=p_m(-1,Y)\)
can be proved directly. The displayed coefficients give
\[
 q_0=1,\quad q_1=Y,\quad
 q_{m+1}=Yq_m-2m q_{m-1}.
\]
These monic polynomials have simple real interlacing zeros by induction.
At successive zeros of \(q_m\), the values of \(q_{m+1}=-2m q_{m-1}\)
alternate in sign, giving \(m-1\) intervening roots. At the largest
zero of \(q_m\), \(q_{m-1}>0\) and \(q_{m+1}<0\), whereas
\(q_{m+1}(Y)>0\) for large positive \(Y\), giving a further root.
At the smallest zero, the sign of \(q_{m-1}\) is \((-1)^{m-1}\);
the sign of \(q_{m+1}\) there is \((-1)^m\), opposite its sign at
large negative \(Y\). This supplies the last root. The degree is
\(m+1\), so all are simple and the interlacing is strict.
The induction begins with \(q_1=Y\) and \(q_2=Y^2-2\).
Parity shows that precisely the odd polynomials have a zero at zero.

Set \(s=-\epsilon^2\) before the event. Dividing the field by
\(\epsilon^m\) gives an analytic family in \((\epsilon,Y)\), including
\(\epsilon=0\), with limiting simple roots those of \(q_m\).
The implicit-function theorem continues every root as a real simple
root for small real positive \(\epsilon\).
After the event, \(s=\epsilon^2\), and
\[
 p_m(1,\ii Y)=\ii^m p_m(-1,Y).
\]
The nonzero limiting roots are therefore purely imaginary and stay
nonreal under a sufficiently small perturbation. If \(m\) is odd,
the limiting root at zero is simple and its unique continued branch
is real by conjugation symmetry. For even \(m\) there is no such root.
Finally, choose a fixed disk containing only the order-\(m\) zero at
the event. Its nonzero boundary and analytic convergence preserve the
total zero count \(m\). The branches just constructed exhaust it.
\end{proof}

This recovers the local forward-parabolic multiplicity pattern in
\cite{Angenent1991,HeatMultiplicity}. It asserts neither that all
canonical events have multiplicity two nor that one can choose the neighborhoods
or remainder constants uniformly in \(m\) or in height.

\subsection{What the collision coefficients require from the source}
At a double event define the complete signed moments
\[
 C_j=\int_0^\infty u^jW_{\lambda_*}(u)\cos(x_*u)\dd u,
 \qquad S_j=\int_0^\infty u^jW_{\lambda_*}(u)\sin(x_*u)\dd u.
\]
Direct differentiation of \eqref{eq:H} gives
\begin{equation}\label{eq:collision-source-moments}
 C_0=S_1=0,\quad a=-C_2,\quad b=S_3,\quad c=C_4,
 \qquad \gamma=-\frac{C_2C_4+S_3^2}{C_2^2}.
\end{equation}
The moments are signed despite \(W_{\lambda_*}>0\). Separate absolute
bounds do not control these ratios without a lower bound on \(|C_2|\).
This is the concrete meaning of a source-specific collision input:
one needs coupled signed information or a paid regular-factor bound,
not positivity of the underlying source alone.

An exact alternative keeps the primitive endpoint. Put
\(\mathfrak P_\lambda(x)=\int_0^\infty p(u)\ee^{\lambda u^2}\cos(xu)\dd u\).
For every real \(\lambda,x\), integration by parts twice gives
\begin{equation}\label{eq:all-time-primitive}
 8H_\lambda=\tfrac12+(2\lambda-1-x^2)\mathfrak P_\lambda
       +4\lambda x\mathfrak P_\lambda'-4\lambda^2\mathfrak P_\lambda''.
\end{equation}
Indeed \(8\int\Phi f=\tfrac12 f(0)+p(0)f'(0)+\int p(f''-f)\),
and for \(f=\ee^{\lambda u^2}\cos(xu)\) the endpoint derivative is zero.
The coefficients of \(f''-f\) are
\((2\lambda-1-x^2+4\lambda^2u^2)\cos(xu)-4\lambda xu\sin(xu)\),
times \(\ee^{\lambda u^2}\). This proves the identity. Differentiating
it finitely many times gives all event jets, with no omitted boundary
constant. It remains a representation of the same oscillatory source.

The energy scales make the distinction between field regularity and
singular observation particularly explicit. If
\(\epsilon=\tau_*-\tau\downarrow0\) and \(I_\tau=(r_-,r_+)\), then
\begin{equation}\label{eq:local-lobe-scales}
 \int_{I_\tau}|h|\sim\frac{4\sqrt2}{3}|a|\epsilon^{3/2},\qquad
 \int_{I_\tau}h^2\sim\frac{16\sqrt2}{15}a^2\epsilon^{5/2},\qquad
 \int_{I_\tau}h_x^2\sim\frac{4\sqrt2}{3}a^2\epsilon^{3/2}.
\end{equation}
For the proof use \(x=x_*+\sqrt{2\epsilon}\,t\) in
\eqref{eq:heat-four-jet}; the limiting profiles are
\(h/\epsilon=a(t^2-1)\) and \(h_x/\sqrt{2\epsilon}=at\), with endpoints
\(t=\pm1\). The integrals of \(1-t^2,(1-t^2)^2,t^2\) are
\(4/3,16/15,2/3\). Uniform remainders on a common compact interval
justify dominated convergence. Thus the local kinetic-to-mass ratio is
\(5/(4\epsilon)\), even though both quantities vanish.
Differentiation on the moving interval gives the exact balances
\[
 \left(\int_{I_\tau}h^2\right)'=-2\int_{I_\tau}h_x^2,\qquad
 \left(\int_{I_\tau}h_x^2\right)'
 =[h_xh_{xx}]_{r_-}^{r_+}-2\int_{I_\tau}h_{xx}^2.
\]
Boundary velocities \(r_\pm'=-h_{xx}/h_x\) produce the latter boundary
term. A local normalized quotient is not the global operator energy.

Finally, for a source perturbation \(h_\eta=h+\eta g+O(\eta^2)\)
at a double event, the Jacobian of \((h,h_x)\) with respect to
\((\tau,x)\) is \(\left(\begin{smallmatrix}a&0\\b&a\end{smallmatrix}\right)\).
It is invertible, so the event continues locally and
\[
 \tau_*'(0)=-g/a,\qquad x_*'(0)=bg/a^2-g_x/a.
\]
This is the inherited joint-event response law \cite{SourceResponse}.
It identifies another denominator that must be paid; it supplies no
uniform stability at arbitrarily flat events.


\section{The complete strip Green potential}\label{sec:green}

\subsection{The finite-factor identity}

Fix
\begin{equation}\label{eq:kernels}
 a=\frac\pi8,\qquad K(t)=\frac1{16\cosh(at)},\qquad
 G(t,y)=\atanh\frac{\sin(ay)}{\cosh(at)}.
\end{equation}
Each copy of \(K\) has mass \(1/2\). Let \(\mu_X\) be the probability
measure on the two lines \(y=0,8\) that has density \(K(x-X)\) on
each. For an admissible function \(f\), write
\[
 \calJ[f](X)=\int\log|f|\dd\mu_X-\log|f(X+4\ii)|,
 \qquad J_\lambda(X)=\calJ[H_\lambda](X).
\]

The Green function of \(0<\Im z<B\), normalized by
\(-\Delta G_B=2\pi\delta\), is
\begin{equation}\label{eq:strip-green}
 G_B(z,\zeta)=\frac12\log
 \frac{\cosh(\pi(\Re z-\Re\zeta)/B)-\cos(\pi(\Im z+\Im\zeta)/B)}
      {\cosh(\pi(\Re z-\Re\zeta)/B)-\cos(\pi(\Im z-\Im\zeta)/B)}.
\end{equation}
Indeed \(z\mapsto\ee^{\pi z/B}\) maps the strip to the upper
half-plane; the logarithm of the reflected-distance ratio gives
\eqref{eq:strip-green}. Its singularity is \(-\log|z-\zeta|\).
At \(B=8\), observation \(X+4\ii\), its value is
\(G(X-\Re\zeta,\Im\zeta)\). Boundary normal derivatives divided
by \(2\pi\) give exactly the two densities \(K\).

For an interior zero \(\zeta\), the singularities of
\(\log|z-\zeta|+G_B(z,\zeta)\) cancel. This harmonic function has the
linear factor's boundary logarithm. Truncating the strip vertically
proves its Poisson representation: the side harmonic measures decay
exponentially while these logarithms grow only logarithmically.
It follows that \(\calJ[z-\zeta]\) is the positive Green value.
An exterior zero gives zero. A boundary zero is handled by its
locally integrable logarithm, or by a limit from the exterior; it
gives no half-weight atom. Finite multiplicities cause no change.

Apply this identity to every complete opposite pair in a finite
truncation of \eqref{eq:product}. Scalar logarithms cancel because
the boundary masses add to one. The only remaining issue is passage
to the complete product.

\subsection{Absolute logarithmic convergence}

Choose
\[
 m_0=\inf_\lambda H_\lambda(0)>0,\qquad
 M_1\ge\sup_{\lambda,|z|\le1}|H_\lambda'(z)|,\quad M_1>0,
\]
and set
\[
 r_*=\min(1,m_0/(2M_1)),\quad
 C_0=\sup_\lambda|\log H_\lambda(0)|,\quad
 C_*=C_0+20C(1+|\log r_*|).
\]
The segment integral of \(H'\) shows that no zero lies in
\(|z|\le r_*\). All constants are uniform for \(0\le\lambda\le1/2\).

\begin{lemma}[Unit-interval logarithmic bound]\label{lem:unit-log}
For \(v\in\R\), \(0\le b\le8\), and \(Q=|v|+12\),
\begin{equation}\label{eq:unit-log}
 \int_v^{v+1}|\log|H_\lambda(x+\ii b)||\dd x
 \le C_*Q\log^2(3Q).
\end{equation}
Moreover the series of paired factor logarithms is absolutely
convergent in local \(L^1\).
\end{lemma}
\begin{proof}
On the interval, \(|z|\le Q\). For a representative \(|\zeta|\le2Q\),
split the logarithm into those of \(z-\zeta\), \(z+\zeta\), and
\(-2\log|\zeta|\). Each positive distance logarithm is at most
\(\log(3Q)\). Its negative part has integral at most two, since
\(|x+\ii b-\zeta|\ge|x-\Re\zeta|\) and
\(\int_{-1}^1-\log|t|\dd t=2\).
The factor therefore costs at most
\(4\log(3Q)+4+2|\log r_*|\).
There are at most \(CQ\log(2Q+2)\) such representatives, by
\eqref{eq:count} and opposite pairing.

For \(|\zeta|>2Q\),
\[
 \left|\log\left|1-\frac{z^2}{\zeta^2}\right|\right|
 \le\frac43\frac{Q^2}{|\zeta|^2}.
\]
Partial summation with half the full count bounds their sum by
\[
 \frac23CQ\{\log(2Q+2)+1\}.
\]
These bounds, together with \(C_0\), give \eqref{eq:unit-log}.
They bound the sum of absolute factor logarithms, which proves the
stronger local \(L^1\) assertion.
\end{proof}

\begin{theorem}[Complete positive potential]\label{thm:green}
For every real \(X\) and \(0\le\lambda\le1/2\), both boundary
logarithms are absolutely integrable against \(K(x-X)\), and
\begin{equation}\label{eq:green-identity}
 J_\lambda(X)=
 \int_\R K(x-X)\{\log|H_\lambda(x)|+\log|H_\lambda(x+8\ii)|\}\dd x
 -\log|H_\lambda(X+4\ii)|
 =\sum_{j:y_j>0}G(X-x_j,y_j).
\end{equation}
The upper-zero sum is absolutely convergent, with all multiplicities.
There is no omitted harmonic term. The potential is nonnegative and
unchanged by multiplication of the field by a nonzero scalar.
\end{theorem}
\begin{proof}
The polynomial-logarithmic bound \eqref{eq:unit-log} is summable
against the exponentially decaying Poisson weights. Thus the finite
paired-product identity passes to the boundary integrals in \(L^1\).
At \(X+4\ii\), the product converges away from every zero by
\eqref{eq:zero-strip}. This proves the complete identity. Its finite
Green sums are nonnegative and increase to the displayed finite
value, hence converge absolutely. The product has already fixed all
harmonic data. Scalar invariance follows directly from the total
boundary mass one.
\end{proof}

\subsection{Two different tails and the exclusion margin}

For \(R\ge0\), let \(n=\lfloor R\rfloor\), \(q_0=\ee^{-a}\), and
\(S=|X|+13+n\). The explicit two-boundary estimate is
\begin{align}\label{eq:boundary-tail}
 &\int_{|x-X|>R}K(x-X)
  \{|\log|H_\lambda(x)||+|\log|H_\lambda(x+8\ii)||\}\dd x
 \nonumber\\
 &\hspace{1cm}\le
 \frac{C_*}{2}\ee^{-an}S\log^2(3S)
 \frac{1+4q_0+q_0^2}{(1-q_0)^4}.
\end{align}
To derive it, cover both sides by unit shells and use
\(K(t)\le\ee^{-a|t|}/8\) and Lemma~\ref{lem:unit-log} on both
boundaries. The remaining sum is bounded by
\(S\log^2(3S)\sum_{j\ge0}q_0^j(1+j)^3\).
The generating function of this cubic moment is the rational factor
in \eqref{eq:boundary-tail}.

For the sharper zero tail we use the additional canonical unit-count
input from the parent analysis \cite{SmoothLocalizationProof}:
after a fixed enlargement of \(A\),
\[
 \#\{j:u\le x_j(\lambda)\le u+1\}\le A\log(2+|u|)
\]
uniformly for \(0\le\lambda\le1/2\).
This inherited input includes the endpoint \(\lambda=0\); it is
separate from the global Jensen count proved above. Since
\begin{equation}\label{eq:G-decay}
 0<G(t,y)\le K_Gy\ee^{-a|t|},\qquad
 K_G=\frac{2a}{\cos^2a},\quad 0<y\le1,
\end{equation}
unit shells give
\begin{equation}\label{eq:zero-tail}
 \sum_{\substack{y_j>0\\|x_j-X|>R}}G(X-x_j,y_j)
 \le2AK_G\ee^{-an}
 \left\{\frac{\log(3+|X|)+\log(n+1)}{1-q_0}
       +\frac{q_0}{(n+1)(1-q_0)^2}\right\}.
\end{equation}
Here \(\atanh u\le u/(1-u^2)\), \(\sin(ay)\le ay\), and
\(\operatorname{sech}(at)\le2\ee^{-a|t|}\) prove
\eqref{eq:G-decay}. For the shell sum use
\(\log(3+|X|+k)\le\log(3+|X|)+\log(k+1)\) and
\(\log(n+j+1)\le\log(n+1)+j/(n+1)\).
The boundary and zero tails pay different truncations. In particular
\(J_\lambda(X)\ll_A\log(3+|X|)\), but that coarse bound is insufficient
for the following consumer.

\begin{corollary}[A single nonreal-zero margin]\label{cor:margin}
If \(d\le y_j\le1\) and \(|x_j-X|\le r\), then
\begin{equation}\label{eq:margin}
 J_\lambda(X)\ge m(d,r):=\frac{\sin(a)d}{\cosh(ar)}.
\end{equation}
An independently acquired strict upper bound below \(m(d,r)\)
excludes this specified zero class.
\end{corollary}
\begin{proof}
Use positivity of every term, \(\atanh u\ge u\), and the concavity
bound \(\sin(ay)\ge y\sin a\) for \(0\le y\le1\).
\end{proof}

An all-real slice has \(J=0\); any upper zero makes \(J>0\) at every
center. Accordingly universal time-zero vanishing is an RH-equivalent
reformulation, not an additional estimate. Uniform tails on compact
center intervals and continuity of finite root multisets show that
\(J\) is continuous through collisions. A pair approaching a double
real zero contributes a quantity linear in its upper height; when
that height has square-root time scale, so does the Green signal.
Continuity alone does not imply Lipschitz continuity or monotonicity.


\section{The exact arithmetic baseline at time zero}\label{sec:prime}

The Green identity is geometric. On the zeta slice, the safe boundary
and center can be removed exactly in favor of a convergent prime sum.
This step preserves the signed boundary information.

\begin{proposition}[Prime decomposition]\label{prop:prime}
Let
\[
 \mathcal L(X)=\int_\R K(x-X)\log|\zeta(1/2-\ii x/2)|\dd x,
\]
and
\begin{equation}\label{eq:P}
 P(X)=\sum_{n\ge2}\frac{\Lambda_{\rm vM}(n)}{\log n}
       \frac{n^{-5/2}}{1+n^{-4}}
       \cos\left(\frac X2\log n\right).
\end{equation}
Then the logarithmic integral is absolutely convergent, the prime-power
series converges absolutely and uniformly on \(\R\), and
\begin{equation}\label{eq:prime}
 J_0(X)=\mathcal L(X)-P(X)+G(X,1).
\end{equation}
For an integer \(Q\ge2\), its truncation satisfies
\begin{equation}\label{eq:prime-tail}
 |P(X)-P_Q(X)|\le\beta_Q:=\frac23Q^{-3/2}.
\end{equation}
\end{proposition}
\begin{proof}
Under \(s(z)=1/2-\ii z/2\), the heights \(0,4,8\) correspond to
real parts \(1/2,5/2,9/2\). In the xi factorization
\eqref{eq:xi}, the logarithmic moduli of \(s\), \(\pi^{-s/2}\),
\(\Gamma(s/2)\), and the constant are harmonic on a neighborhood
of the closed strip. The zero of \(s\) is at \(-\ii\), outside
the strip. Gamma has neither zeros nor poles there, and Stirling
growth makes its logarithmic modulus at most linear plus logarithmic
in horizontal height \cite{DLMFGreenPrime}. Exponential Poisson
decay therefore removes vertical-side terms and these harmonic
contributions have zero \(\calJ\).

The factor \(s-1\) has its zero at \(z=\ii\) and contributes
\emph{plus} \(G(X,1)\). This accounts for the pole of zeta; the
zeta logarithm is not treated as harmonic across \(s=1\).
Subtracting the elementary and gamma logarithms from
Lemma~\ref{lem:unit-log} also proves weighted absolute integrability
of the critical-line logarithm.

For \(\sigma>1\), the Euler product gives
\[
 \log|\zeta(\sigma-\ii x/2)|
 =\sum_{n\ge2}\frac{\Lambda_{\rm vM}(n)}{\log n}
       n^{-\sigma}\cos((x/2)\log n).
\]
The substitution \(u=\ee^{2at}\) and the beta integral yield
\[
 \widehat K(\omega)
 =\frac1{16a}\int_0^\infty
       \frac{u^{-1/2-\ii\omega/(2a)}}{1+u}\dd u
 =\frac12\operatorname{sech}(4\omega).
\]
At \(\omega=(\log n)/2\), this is \((n^2+n^{-2})^{-1}\).
The safe-boundary minus center coefficient is consequently
\[
 \frac{n^{-9/2}}{n^2+n^{-2}}-n^{-5/2}
 =-\frac{n^{-5/2}}{1+n^{-4}}.
\]
Absolute convergence justifies the interchanges and proves
\eqref{eq:prime}. Finally
\(\Lambda_{\rm vM}(p^k)/\log(p^k)=1/k\), and the coefficient
vanishes off prime powers. Thus
\[
 |P-P_Q|\le\sum_{n>Q}n^{-5/2}
          \le\int_Q^\infty t^{-5/2}\dd t=\frac23Q^{-3/2}.
\]
The integer cutoff is part of this integral-test bound.
\end{proof}

A useful pole-cancelled field is
\begin{equation}\label{eq:F}
 q(s)=\frac{s-1}{s},\qquad
 F(z)=q(s(z))\zeta(s(z))=D(z)H_0(z),\qquad
 D(z)=\frac{16\pi^{s(z)/2}}{s(z)^2\Gamma(s(z)/2)}.
\end{equation}
The multiplier \(D\) is holomorphic and zero-free on a neighborhood
of the closed strip, with harmonic logarithm of admissible growth.
Therefore \(\calJ[F]=J_0\). On the real boundary \(|q|=1\), so
\begin{equation}\label{eq:F-prime}
 J_0(X)=\int_\R K(x-X)\log|F(x)|\dd x-P(X)+G(X,1).
\end{equation}
The plus pole correction occurs once in this \emph{single-boundary}
prime representation. It is not added again to the already
pole-cancelled \emph{two-boundary} functional \(\calJ[F]\).

Positivity proves the lower comparison
\(\mathcal L(X)\ge P(X)-G(X,1)\).
The exclusion problem instead asks for an upper comparison
\[
 \mathcal L(X)<P(X)-G(X,1)+m(d,r).
\]
No such favorable bound is supplied by Proposition~\ref{prop:prime}.
No Euler product is assigned to nonzero Newman time.


\subsection{A complete pole-cancelled formula}
For \(\Re s>0\), the fractional-part formula gives
\begin{equation}\label{eq:fractional}
 q(s)\zeta(s)=1-(s-1)\int_1^\infty\{v\}v^{-s-1}\dd v.
\end{equation}
This follows for $\Re s>1$ by integrating the integer-counting function, and extends to $\Re s>0$ because the fractional-part integral converges locally uniformly there. The cancellation at $s=1$ is removable. In particular, it provides polynomial growth bounds without assuming any zero-free region on the critical line.

\section{One-sided boundary-logarithm approximation through zeros}\label{sec:boundary-log}

We now acquire the signed logarithm in the exact prime identity.
The inherited local theorem \cite{OGBoundaryLog,OGAssemblyAssessment}
controls how much an absolute field error can raise its integral, even at
multiple real zeros. Its starting data are a nonzero anchor away
from the real interval and an upper bound on a larger disk.

\subsection{A one-dimensional sublevel theorem}

Fix the nested radii and constants
\begin{equation}\label{eq:disk-constants}
 R_0=\frac92,\quad \rho_0=\frac{13}3,\quad r_0=\frac{25}6,\quad
 B=R_0+r_0=\frac{26}3,\quad
 \kappa=\log\frac{27}{26},\quad
 \frac{\rho_0+r_0}{\rho_0-r_0}=51 .
\end{equation}

\begin{theorem}[Line sublevels and a logarithmic upper envelope]\label{thm:local-log}
Let \(f\) be holomorphic on a neighborhood of the closed disk
\(\overline{D(0,R_0)}\), with \(|f|\le1\) there and
\(|f(0)|\ge m>0\). Put
\begin{equation}\label{eq:n-C}
 A=\log(1/m),\qquad
 n=\max\{2,\lceil A/\kappa\rceil\},\qquad
 C_f=2Bn\ee^{51A/n}.
\end{equation}
Let \(\Gamma\) be a horizontal segment of length at most one inside
\(\overline{D(0,r_0)}\). Then, for every \(t>0\),
\begin{equation}\label{eq:sublevel}
 |\{z\in\Gamma:|f(z)|<t\}|\le C_ft^{1/n}.
\end{equation}
If \(g\) is measurable on \(\Gamma\) and
\(|f-g|\le\varepsilon\), with \(\varepsilon>0\), then
\begin{equation}\label{eq:local-envelope}
 0\le\int_\Gamma
       [\log(|g|+\varepsilon)-\log|f|]\dd\ell
 \le\frac{C_fn^2}{n-1}(2\varepsilon)^{1/n}.
\end{equation}
No lower bound for \(|f|\) on \(\Gamma\), simplicity or separation
of its zeros, or nonvanishing or holomorphy of \(g\), is assumed.
\end{theorem}

\begin{proof}
Jensen's formula \cite{Tao2020LocalZeros} bounds the number \(k\)
of zeros inside the open \(\rho_0\)-disk, including multiplicity,
by \(k\kappa\le A\). Remove exactly these zeros with disk Blaschke
factors:
\[
 b_j(z)=\frac{R_0(z-a_j)}{R_0^2-\overline{a_j}z},\qquad
 h(z)=f(z)\Big/\prod_{j=1}^k b_j(z).
\]
Fill the removable singularities. Since every factor has boundary
modulus one, the maximum principle gives \(|h|\le1\) on the outer
disk. It is zero-free on the open \(\rho_0\)-disk and
\(|h(0)|\ge m\). Harnack's inequality for the nonnegative harmonic
function \(-\log|h|\) gives
\(|h(z)|\ge\ee^{-51A}\) when \(|z|\le r_0\)
\cite{Tao2016Harnack}. If \(h\) has a zero on the \(\rho_0\)-circle,
apply Harnack on smaller disks and pass to the limit.
The elementary denominator bound for each Blaschke factor gives
\begin{equation}\label{eq:factor-lower}
 |f(z)|\ge\ee^{-51A}
       \prod_{j=1}^k\frac{|z-a_j|}{B}\qquad(|z|\le r_0).
\end{equation}
Every distance ratio is less than one.

Let \(q=\ee^{51A}t\). If \(q<1\) and \(k=0\), the sublevel set
is empty. Otherwise, when \(|f(z)|<t\), at least one distance
is below \(Bq^{1/k}\). The set is covered by \(k\) disks of this
radius. A disk of radius \(b\) meets a line in length at most \(2b\),
so the total length is at most
\[
 2Bkq^{1/k}\le2Bnq^{1/n}=C_ft^{1/n}.
\]
For \(q\ge1\), this bound exceeds the segment length and is
automatic. This proves \eqref{eq:sublevel} in one dimension.
An area estimate would not suffice for this trace argument.

The triangle inequality gives, off the finitely many segment zeros,
\[
 0\le\log(|g|+\varepsilon)-\log|f|
 \le\log(1+2\varepsilon/|f|).
\]
Write \(h_*=2\varepsilon\). Tonelli and the layer-cake identity give
\[
 \int_\Gamma\log(1+h_*/|f|)\dd\ell
 =\int_0^\infty
 \frac{h_*}{t(t+h_*)}|\{|f|<t\}\cap\Gamma|\dd t.
\]
On \(0<t<h_*\), bound the kernel by \(1/t\); on \(t>h_*\),
bound it by \(h_*/t^2\). Equation~\eqref{eq:sublevel} then gives
the two bounds \(C_fn h_*^{1/n}\) and
\(C_fn h_*^{1/n}/(n-1)\). Their sum is
\eqref{eq:local-envelope}. The same sublevel integral proves local
integrability of \(\log^-|f|\). The choice \(n\ge2\) makes the
second integral finite and avoids the endpoint logarithm at \(n=1\).
\end{proof}

An immediate consequence, for \(M\ge1\), is
\begin{equation}\label{eq:absolute-log-local}
 \int_\Gamma|\log|Mf||\dd\ell\le\log M+C_fn.
\end{equation}
Indeed \(|\log|Mf||\le\log M+\log^-|f|\), and integrating
\eqref{eq:sublevel} over \(0<t<1\) with measure \(\dd t/t\)
bounds the latter integral by \(C_fn\).

\subsection{The canonical anchor and explicit height dependence}

For \(v\in\R\), center these disks at \(c_v=v+4\ii\), and take
the segment \(I_v=[v-1/2,v+1/2]\). Its maximum distance from
the center is \(\sqrt{65}/2<r_0\).
Throughout the outer disk, \(1/4\le\Re s\le19/4\).
The pole-cancelled fractional-part identity
\eqref{eq:fractional} gives
\[
 |F(z)|\le2|v|+25.
\]
In detail \(|s-1|\le |v|/2+6\), while the integral in that
identity is bounded by \(1/\Re s\le4\). At the center,
the absolutely convergent Euler product gives
\begin{equation}\label{eq:canonical-anchor}
 |F(c_v)|\ge c_0:=\frac35\frac{\zeta(5)}{\zeta(5/2)}>0.
\end{equation}
This safe-line estimate is unconditional and supplies no lower
modulus on the real segment.

Define
\begin{equation}\label{eq:local-canonical-constants}
 M_v=2|v|+26,\quad A_v=\log(M_v/c_0),\quad
 n_v=\max\{2,\lceil A_v/\kappa\rceil\},\quad
 C_v=2Bn_v\ee^{51A_v/n_v}.
\end{equation}
Apply Theorem~\ref{thm:local-log} to \(F/M_v\).

\begin{corollary}[Canonical local acquisition]\label{cor:canonical-log}
If \(\sup_{I_v}|F-\widehat F|\le\delta\), with \(\delta>0\), then
\begin{equation}\label{eq:canonical-envelope}
 0\le\int_{I_v}
 [\log(|\widehat F|+\delta)-\log|F|]\dd x
 \le\mathcal E_v(\delta)
 :=\frac{C_vn_v^2}{n_v-1}
       \left(\frac{2\delta}{M_v}\right)^{1/n_v}.
\end{equation}
Also
\begin{equation}\label{eq:canonical-absolute-log}
 \int_{I_v}|\log|F(x)||\dd x\le\log M_v+C_vn_v.
\end{equation}
\end{corollary}
\begin{proof}
Use \(f=F/M_v\), \(g=\widehat F/M_v\),
\(\varepsilon=\delta/M_v\). Logarithmic constants cancel
exactly in the difference. Equation~\eqref{eq:absolute-log-local}
gives the second assertion.
\end{proof}

Because \(A_v/n_v\le\kappa\), the prefactor in
\eqref{eq:canonical-envelope} is \(O(\log^2(2+|v|))\),
whereas \(n_v\) is of order \(\log(2+|v|)\). To make the local
error at most \(h>0\), it suffices to take
\begin{equation}\label{eq:precision}
 \delta\le\frac{M_v}{2}
 \left(\frac{h(n_v-1)}{C_vn_v^2}\right)^{n_v}.
\end{equation}
This is finite precision for each fixed interval, not a uniform
positive Hölder exponent or a useful complexity theorem.

The local double-zero model illustrates the relevant scale without
any numerical approximation:
\begin{equation}\label{eq:double-log}
 \int_\R[\log(x^2+\delta)-\log x^2]\dd x=2\pi\sqrt\delta.
\end{equation}
Substitute \(x=\sqrt\delta\,u\) and integrate
\(\log(1+u^{-2})\) by parts on the positive half-line; its integral
is \(2\int_0^\infty(1+u^2)^{-1}\dd u=\pi\).
Thus a real zero produces no fixed error floor, although a Lipschitz
error bound fails even in this simple model. This is an algebraic
scale calculation, not a canonical collision control, and the
conservative exponent in \eqref{eq:canonical-envelope} is not
claimed sharp.

The approximant in this theorem may have zeros, and its logarithmic
integral is retained. This is the precise distinction from a
zero-free harmonic-baseline replacement. Likewise no small
quadratic-residual hypothesis from Section~\ref{sec:quadratic}
is used. The theorem measures signed information already present
in the exact field; it does not establish its favorable value.


\section{Complete finite-source acquisition of the signed logarithm}
\label{sec:signed-acquisition}

The local theorem can be assembled into a complete one-sided Green
estimate. We give the full construction, rather than regard a uniformly
accurate approximation as an evaluated exclusion
\cite{OGAssemblyProof,OGAssemblyAssessment}.

\subsection{The literal source and its two independent tails}
For an integer \(N\ge1\), reflect the first \(N\) positive-half-line
terms evenly and set
\[
 \Phi_N^{\rm ref}(u)=\ee^{|u|}\sum_{n=1}^N
 [2(n^2\rho(u))^2-3n^2\rho(u)]\ee^{-n^2\rho(u)},\qquad
 \widehat H(z)=\frac12\int_{-V}^V\Phi_N^{\rm ref}(u)\ee^{\ii zu}\dd u.
\]
Let \(\widehat F=D\widehat H\), with \(D\) as in \eqref{eq:F}, and put
\[
 \mathcal Q=\int_\R\ee^{9|u|}\rho(u)^2\ee^{-\rho(u)}\dd u,
 \qquad \mathcal Q^{>V}=\int_{|u|>V}\ee^{9|u|}\rho(u)^2\ee^{-\rho(u)}\dd u.
\]
The arithmetic tail inside the radial cutoff and the complete source
outside it give, with \(B_0=5C_2(0)\),
\begin{equation}\label{eq:delta-H}
 \sup_{|\Im z|\le8}|H_0(z)-\widehat H(z)|
 \le\delta_H=\tfrac12[5C_2(N)\mathcal Q+B_0\mathcal Q^{>V}].
\end{equation}
This follows immediately from Lemma~\ref{lem:source} and
\(|\ee^{\ii zu}|\le\ee^{8|u|}\). The reflected finite source need not
be smooth at zero; this undifferentiated integral argument does not
require it to be. The finite integral is entire in \(z\).

These errors have terminating explicit majorants. Successive term
ratios decrease, so
\[
 C_2(N)\le
 \frac{(N+1)^4\ee^{-\pi((N+1)^2-1)}}
 {1-((N+2)/(N+1))^4\ee^{-\pi(2N+3)}}.
\]
With \(T_V=\pi\ee^{4V}\), substitution in the radial integral gives
\begin{equation}\label{eq:radial-majorant}
 \mathcal Q^{>V}=\frac1{2\pi^{9/4}}\int_{T_V}^\infty t^{13/4}\ee^{-t}\dd t
 \le\frac{\ee^{-T_V}}{2\pi^{9/4}}
       (T_V^4+4T_V^3+12T_V^2+24T_V+24).
\end{equation}
The same bound at \(T_V=\pi\) bounds \(\mathcal Q\). Both errors tend
to zero separately.

The normalization multiplier must also be paid. On
\(|\Re z-X|\le R\), \(0\le\Im z\le8\), one sufficient bound is
\begin{equation}\label{eq:multiplier}
 |D(z)|\le\mathcal B(X,R)
 =16\pi^{9/4}\exp\left\{\frac94+
                 \frac{\pi^2[81+(|X|+R)^2]}{192}\right\}.
\end{equation}
To derive it, use the reciprocal-gamma product
\cite[5.8.2]{DLMFReciprocalGammaOC}. For \(w=s/2\),
\[
 |\Gamma(w)^{-1}|
 \le|w|\exp\{\gamma\Re w+\pi^2|w|^2/12\},
\]
because \(\log|1+w/n|\le\Re(w/n)+|w|^2/(2n^2)\).
Now \(1/2\le\Re s\le9/2\), \(|s|\ge1/2\),
\(|s|^2\le[81+(|X|+R)^2]/4\), and \(\gamma<1\).
Substitution into \eqref{eq:F} proves \eqref{eq:multiplier}.
Thus the compact error in \(F\) is at most \(\mathcal B\delta_H\).
This deliberately coarse compact bound is never integrated along an
infinite remote boundary.

\subsection{The exact boundary tail}
Use \(B,\kappa,c_0\) from Section~\ref{sec:boundary-log}, and set
\[
 L_0=\log(26/c_0),\quad N_0=2+L_0/\kappa,\quad
 C_{\log}=\log26+2B\ee^{51\kappa}N_0^2,\quad q=\ee^{-a}.
\]
Equation~\eqref{eq:canonical-absolute-log} is bounded by
\(C_{\log}[1+\log(1+|v|)]^2\) on the unit interval \(I_v\).
For an integer \(R\ge1\), covering the two remote half-lines with
unit intervals consequently gives
\begin{equation}\label{eq:assembled-log-tail}
 \int_{|x-X|>R}K(x-X)|\log|F(x)||\dd x
 \le\mathcal T(X,R):=\frac{C_{\log}}4\ee^{-aR}
 [1+\log(|X|+R+2)]^2\frac{1+q}{(1-q)^3}.
\end{equation}
Indeed the two interval weights at offset \(k\ge R\) are each at
most \(\ee^{-ak}/8\). Bound the remaining logarithmic factor by its
value at \(R\) times \((1+k-R)^2\), and sum
\(\sum_{j\ge0}q^j(1+j)^2=(1+q)/(1-q)^3\).
This is a tail of the exact logarithm, not a tail of its approximant.

\subsection{A one-sided assembly with all constants retained}
Assume \(\sup_{|x-X|\le R}|F-\widehat F|\le\delta\),
where \(0<\delta\le1\). Write
\begin{align}
 \mathcal I&=\int_{X-R}^{X+R}K(x-X)
                    \log(|\widehat F(x)|+\delta)\dd x,\nonumber\\
 \mathcal E_{X,R}(\delta)&=\sum_{k=0}^{R-1}
 \frac{\mathcal E_{X+k+1/2}(\delta)+\mathcal E_{X-k-1/2}(\delta)}
 {16\cosh(ak)},\label{eq:compact-log-envelope}\\
 U_{R,Q,\delta}(X)&=\mathcal I+\mathcal T(X,R)-P_Q(X)+\beta_Q+G(X,1).
 \label{eq:signed-upper}
\end{align}
Here \(Q\ge2\) is an integer and \(\mathcal E_v\) is exactly
\eqref{eq:canonical-envelope}.

\begin{theorem}[Complete one-sided acquisition]\label{thm:signed-acquisition}
Under these assumptions,
\begin{equation}\label{eq:signed-acquisition-error}
 0\le U_{R,Q,\delta}(X)-J_0(X)
 \le\mathcal E_{X,R}(\delta)+2\mathcal T(X,R)+2\beta_Q.
\end{equation}
For each fixed \(X\) and every positive error allowance, finite
\(R,Q,N,V\) and a positive \(\delta\) can make this bound smaller
than the allowance. No conclusion about the size of \(J_0(X)\) follows
from this acquisition assertion alone.
\end{theorem}
\begin{proof}
The integrand difference from \(\log|F|\) is nonnegative almost
everywhere. On each unit interval apply Corollary~\ref{cor:canonical-log}
and multiply by the maximum kernel weight, obtaining
\eqref{eq:compact-log-envelope}. Let \(L_{\rm comp},L_{\rm tail}\)
be the exact compact and remote logarithmic integrals. Subtracting
\eqref{eq:F-prime} from \eqref{eq:signed-upper} gives
\[
 U-J_0=(\mathcal I-L_{\rm comp})+
        (\mathcal T-L_{\rm tail})+(P-P_Q+\beta_Q).
\]
The three terms lie respectively in
\([0,\mathcal E]\), \([0,2\mathcal T]\) and \([0,2\beta_Q]\).
This proves the inequality, including both factors of two.
Given an allowance \(h>0\), choose \(R\) so that \(\mathcal T<h/6\),
then \(Q\) so that \(\beta_Q<h/6\). The finite sum \(\mathcal E\)
tends to zero with \(\delta\), so choose it with \(\mathcal E<h/3\).
Finally choose \(N,V\) using \eqref{eq:delta-H}--\eqref{eq:multiplier}
to realize that compact error. This order avoids using a multiplier
bound on an unchosen or unbounded domain.
\end{proof}

For fixed \(\delta>0\), the compact approximate integrand is continuous
and has a finite modulus of continuity. Validated subdivision can
therefore evaluate its integral to any specified positive tolerance.
If a directed evaluation satisfies \(U\le\overline U\le U+\eta_{\rm eval}\),
the complete upper error adds \(\eta_{\rm eval}\). An interval's width,
rather than an unqualified midpoint radius, pays this one-sided step.
Finite termination at each prescribed tuple is not a common all-height
complexity estimate. The next section gives a different integration
method that exploits certified local root information.

\section{Deflated logarithms, contour floors and signed integration}
\label{sec:deflation}

The generic logarithmic envelope uses no real-boundary modulus floor.
When a certified local divisor is available, one can instead integrate
its singular logarithms explicitly and approximate the remaining regular
factor. The following construction is inherited from the assessed
deflation note and its later conditional contour design
\cite{OGDeflation,OGDeflationAssessment,OGContourDesign}.

\subsection{Stable factor acquisition by divided differences}
On a real patch \(I_v\), suppose the relevant actual roots of \(H_0\)
are the multiset \(z_1,\ldots,z_n\), and define the removable quotient
\[
 U_v(z)=H_0(z)\Big/\prod_{j=1}^n(z-z_j).
\]
Repeated roots are repeated in the list. Deflate \(H_0\), whose entire
domain is available, before multiplying by \(D\); the meromorphic
extension of \(F\) has a pole at \(-\ii\).

For a holomorphic \(f\), its divided difference has the simplex formula
\begin{equation}\label{eq:divided-simplex}
 f[z_1,\ldots,z_n,x]
 =\int_{\Delta_n}f^{(n)}\left(t_xx+\sum_{j=1}^nt_jz_j\right)\dd t,
 \qquad \operatorname{vol}(\Delta_n)=1/n!.
\end{equation}
It follows by iterating the fundamental theorem of calculus; continuity
in the knots supplies repeated-node values. When the knots are the
actual roots, with their multiplicities, the interpolation polynomial
vanishes and \eqref{eq:divided-simplex} equals \(U_v(x)\).

Suppose certified matching gives approximate knots \(w_j\) with
\(|z_j-w_j|\le\rho_j\), and let \(\widehat H\) approximate \(H_0\)
on the full convex domain of the knots and patch. Define
\(\widehat U_v(x)=\widehat H[w_1,\ldots,w_n,x]\).
If \(\epsilon_k\) bounds the error in derivative \(k\) of \(H_0\), and
\(\widehat M_k\) bounds \(|\widehat H^{(k)}|\), then
\begin{equation}\label{eq:deflated-errors}
 \|U_v^{(j)}-\widehat U_v^{(j)}\|_{I_v}
 \le e_j:=\frac{j!\epsilon_{n+j}}{(n+j)!}
       +\frac{j!\widehat M_{n+j+1}}{(n+j+1)!}\sum_{l=1}^n\rho_l,
 \quad j=0,1,2.
\end{equation}
Differentiate the simplex integral: the factor is \(t_x^j\).
Moving a knot adds \(t_l\rho_l\) and one derivative. The simplex
moments are \(j!/(n+j)!\) and \(j!/(n+j+1)!\), respectively.
This proves the estimate without division by a root separation.
Approximate knots require derivatives through order \(n+3\).
The function \(\widehat U_v\) is a divided difference, not
\(\widehat H/\prod(z-w_j)\); the latter would leave spurious poles.

If \(|\widehat U_v|\ge\mu_v>e_0\) on the real patch, integration of
the derivative of the real logarithm along a segment of moduli gives
\begin{equation}\label{eq:factor-log-error}
 |\log|U_v|-\log|\widehat U_v||\le b_v:=\frac{e_0}{\mu_v-e_0}.
\end{equation}
With \(W_v=\int_{I_v}K(x-X)\dd x\), the factor error is
\(E_{\rm fac}=\sum_vW_vb_v\), and \(\sum_vW_v\le1/2\).
The floor is an additional premise, not an automatic consequence of
divided differences.

\subsection{Paying root uncertainty through real zeros}
For an interval of length \(L\), two complex points within distance
\(\rho\) satisfy the root-log estimate from the deflation source
\cite[Section 3]{OGDeflation}
\begin{equation}\label{eq:root-log-error}
 \int_I|\log|x-z|-\log|x-w||\dd x
 \le A(L,\rho):=4\rho\left[1+\log\left(1+\frac{L}{2\rho}\right)\right],
 \quad A(L,0)=0.
\end{equation}
For a direct proof, the triangle inequality gives
\[
 |\log|x-z|-\log|x-w||
 \le\log\left(1+\frac{\rho}{\min(|x-z|,|x-w|)}\right)
 \le\sum_{a\in\{\Re z,\Re w\}}\log(1+\rho/|x-a|).
\]
For each center the integral is largest when the interval is centered
there, because the integrand is even and decreases with distance. Its
value is
\[
 2\int_0^{L/2}\log(1+\rho/t)\dd t
 =L\log(1+2\rho/L)+2\rho\log(1+L/(2\rho))
 \le2\rho[1+\log(1+L/(2\rho))].
\]
Adding the two centers proves \eqref{eq:root-log-error}, including real
roots by integrable limiting values. No simple-zero slope is used.
Since \(K\le1/16\), one may charge
\[
 E_{\rm roots}=\frac1{16}\sum_{v,j}A(|I_v|,\rho_{vj}).
\]
Repeated uses of a root in different patches are charged again.
This is stability in root location. Source perturbations can move
roots by a square root, so it is not a Lipschitz theorem in source error.

\subsection{Product integration with exact singular primitives}
Set \(q_v(x)=\log|D(x)|+\log|\widehat U_v(x)|\) and assume
\(|q_v|\le H_v\), \(|q_v''|\le B_v\). On panels of maximum width
\(\Delta\), linearly interpolate both \(K\) and \(q_v\), denoting
the interpolants by \(K_h,q_h\). Use
\begin{equation}\label{eq:product-integration}
 Q_h=\sum_v\int_{I_v}K_h(x)
       \left(q_h(x)+\sum_j\log|x-w_{vj}|\right)\dd x.
\end{equation}
The regular part is polynomial integration. For a root with real part
\(a_j\) and \(b_j=|\Im w_j|\), put \(t=x-a_j\). The two singular
primitives needed for the linear kernel are
\begin{align}
 A_0(t,b)&=\tfrac t2\log(t^2+b^2)-t+b\arctan(t/b),\nonumber\\
 A_1(t,b)&=\tfrac14[(t^2+b^2)\log(t^2+b^2)-(t^2+b^2)].
 \label{eq:root-primitives}
\end{align}
Their derivatives are \(\log\sqrt{t^2+b^2}\) and
\(t\log\sqrt{t^2+b^2}\). Continuous limits at \(b=0\) include roots
at panel endpoints; there is no exclusion of multiple real roots.

An explicit common bound for the unweighted absolute integrals of the
regular part and singular logarithms is
\[
 \mathcal S:=\sum_v\left[|I_v|H_v+
       \sum_j\{|I_v|\log^+D_{vj}+2\}\right],\qquad
 D_{vj}=\max_{x\in I_v}|x-w_{vj}|.
\]
The positive logarithm uses its maximum; the integral of the negative
logarithm is at most \(\int_{-1}^1-\log|t|\dd t=2\).
Linear interpolation has uniform error at most \(\Delta^2/8\) times
the second-derivative bound, and \(\|K''\|_\infty\le a^2/16\).
Split the regular integrand as
\(Kq_v-K_hq_h=K(q_v-q_h)+(K-K_h)q_h\).
Linear interpolation preserves the bound \(|q_h|\le H_v\), so the same
\(\mathcal S\) also bounds the integral of the interpolated regular
part plus the singular logarithms. Consequently
\begin{equation}\label{eq:product-error-safe}
 E_{\rm quad}\le\frac{\Delta^2}{8}
       \left(\sum_v W_v B_v+\frac{a^2}{16}\mathcal S\right)
 \le\Delta^2\left(\frac{B_*}{16}+\frac{a^2\mathcal S}{128}\right),
 \qquad B_*=\max_v B_v.
\end{equation}
Only the first term uses exact Poisson weights \(W_v\). Numerical
node errors, when later multiplied by the interpolated kernel, instead
use its explicitly computable weights \(\int K_h\). This distinction
avoids an unproved assumption that interpolation preserves Poisson mass.
For fixed paid constants the panel-width cost is of order
\(E_{\rm quad}^{-1/2}\); the source, count and factor costs remain.

\subsection{Contour acquisition of the factor floor}
Suppose a closed disk of radius \(A\), centered at real \(c\), contains
exactly \(n\) roots with multiplicity, all at distance at most \(s<A\)
from \(c\), and
\(0<m\le|H_0|\le M\) on its boundary. Deflating all these roots gives
a holomorphic, zero-free function \(U\) on the closed disk. The maximum
principle applied to \(U\) and \(1/U\) proves
\begin{equation}\label{eq:contour-factor-floor}
 \mu_0:=\frac{m}{(A+s)^n}\le|U(z)|\le
 V_0:=\frac{M}{(A-s)^n}.
\end{equation}
Thus a paid boundary margin and complete root count can replace a
direct search for a small real-axis denominator. They cannot be omitted.
For a polynomial source model \(p\), Rouch\'e's theorem on the outer
boundary, together with disjoint inner cluster counts summing to the
outer count, supplies the required completeness. A list of located
roots without this count does not rule out missed roots in the annulus.

If \eqref{eq:deflated-errors} holds on the entire disk and
\(e_0<\mu_0\), the divided-difference model \(\widetilde U\) is zero-free
there and lies between \(\mu=\mu_0-e_0\) and \(V=V_0+e_0\).
For a real subinterval \(|x-c|\le\ell<A\), harmonicity of
\(\log|\widetilde U|\) gives
\begin{equation}\label{eq:factor-curvature}
 |(\log|\widetilde U|)''(x)|
 \le\frac{2\log(V/\mu)}{(A-\ell)^2}.
\end{equation}
For example, on a circle of radius \(r<A-\ell\), subtract the midpoint
of its allowed logarithmic range. The second Fourier coefficient of
this bounded harmonic function gives a bound
\(4\log(V/\mu)/(\pi r^2)\le2\log(V/\mu)/r^2\); let \(r\uparrow A-\ell\).
On the real axis, direct differentiation of \eqref{eq:F} gives
\[
 (\log|D|)''=\Re\{-1/(2s^2)+\psi'(s/2)/16\},\qquad
 |(\log|D|)''|\le13/4.
\]
Indeed \(|s|\ge1/2\) and the trigamma series at \(\Re(s/2)=1/4\)
is bounded by \(16+\int_0^\infty(t+1/4)^{-2}\dd t=20\).
Together these bounds supply the \(B_v\) needed in
\eqref{eq:product-error-safe}.

All derivatives used above can be acquired from the original source.
For a degree-\(L\) Taylor polynomial about \(c\), if coefficient
errors are \(\delta_j\), finite-source derivative bound is
\(M_{L+1,A}^{\rm fin}\), and the complete omitted-source derivative
allowance is \(E_{{\rm src},k}\), then
\begin{equation}\label{eq:taylor-source-error}
 \epsilon_k\le E_{{\rm src},k}
 +\sum_{j=k}^L(j)_k\delta_jA^{j-k}
 +M_{L+1,A}^{\rm fin}\frac{A^{L+1-k}}{(L+1-k)!}.
\end{equation}
The moment integral of \(|u|^k\ee^{A|u|}\) against the arithmetic and
radial tails proves the first term. Taylor's integral remainder proves
the last. Increasing arithmetic precision does not remove either of
these analytic remainders.

\subsection{The signed error and an identical-information comparison}
Let a numerical enclosure of \(Q_h-P_Q+G(X,1)\) have center \(S_c\)
and radius \(e_{\rm num}\). Combining the preceding estimates yields
\begin{equation}\label{eq:deflation-ledger}
 |J_0(X)-S_c|\le\mathcal R
 :=\mathcal T+\beta_Q+E_{\rm fac}+E_{\rm roots}+E_{\rm quad}+e_{\rm num}.
\end{equation}
Consequently \(J_{\rm up}=S_c+\mathcal R\) satisfies
\(0\le J_{\rm up}-J_0\le2\mathcal R\). Small acquisition error is
separate from the strict value test \(J_{\rm up}<m(d,r)\).

There is also a useful scoped comparison with direct zero information.
Suppose a complete collection of rectangular cells
\(C_i=[l_i,u_i]\times[b_i,c_i]\) carries integer counts \(n_i\).
For each cell in the upper strip, the independent Green maximum is
\[
 U_i=G(\operatorname{dist}(X,[l_i,u_i]),\max(c_i,0)).
\]
If an occupied cell intersects the target
\([X-r,X+r]+\ii[d,1]\), monotonicity of \(G\) gives
\(U_i\ge G(r,d)\ge m(d,r)\). Therefore a strict pass of
\(\sum_i n_iU_i+T_G<m(d,r)\), with \(T_G\ge0\) a complete tail,
already implies that every occupied cell is disjoint from the target.
The same data thus directly exclude zeros there. This proves dominance
only for independent rectangular maxima with identical complete cell
information. It is neither a cost theorem nor an obstruction to coupled
signed estimates, different information, or reuse of an acquired observable.

\section{Riemann--Siegel structure and Newman-time evolution}
\label{sec:RS}

The preceding acquisition theorems accept complete approximants with
paid errors. Riemann--Siegel theory supplies such approximants by
preserving oscillatory arithmetic terms and their remainder. It is
useful to separate three questions: where the classical saddles lie,
how the complete expression evolves in Newman time, and how its error
passes through an integrated logarithm. The first two identities alone
do not answer the signed upper-value question.

\subsection{The critical strip comes before the critical line}
For the standard oriented line \(\mathcal C\) from northeast to
southwest crossing the real interval \((0,1)\), with the usual
continuous branch of \(\log w\), Riemann's auxiliary function is
\begin{equation}\label{eq:RS-auxiliary}
 \mathscr R(s)=\int_{\mathcal C}
 \frac{w^{-s}\ee^{\pi\ii w^2}}{\ee^{\pi\ii w}-\ee^{-\pi\ii w}}\dd w.
\end{equation}
Its reflected identity is
\begin{equation}\label{eq:RS-reflection}
 \zeta(s)=\mathscr R(s)+\chi(s)\overline{\mathscr R(1-\bar s)},
 \qquad
 \chi(s)=\pi^{s-1/2}\frac{\Gamma((1-s)/2)}{\Gamma(s/2)}.
\end{equation}
These are the classical entire auxiliary integral and reflected zeta
formula; at exceptional factors they are understood by continuation
\cite[Section 3, Definition 4 and Theorem 5]{Arias2024AuxiliaryBasic}.
Contour displacement across integer poles produces finite Dirichlet
sums plus a remaining saddle contour. The corresponding two-sum
approximate functional equation is valid off the line as well. A
balanced cutoff near \(\sqrt{t/(2\pi)}\) does not require
\(\Re s=1/2\) \cite[Sections 2 and 8]{Arias2024AuxiliaryExpansions}.

On \(s=1/2+\ii t\), the two terms are related by conjugacy and
\(|\chi(s)|=1\). Choose the continuous gamma phase \(\theta(t)\) so
that \(\chi(1/2+\ii t)=\ee^{-2\ii\theta(t)}\). Multiplying by
\(\ee^{\ii\theta(t)}\) makes the full expression real, and its balanced
finite sums combine into
\begin{equation}\label{eq:RS-cosine}
 Z(t)=2\sum_{n\le\lfloor\sqrt{t/(2\pi)}\rfloor}
 n^{-1/2}\cos(\theta(t)-t\log n)+\mathcal E_{\rm RS}(t).
\end{equation}
The correction \(\mathcal E_{\rm RS}\) is retained. This is a
large-height expansion specialized by symmetry, not an expansion
assuming all zeros lie on the line. Its reality follows from the
functional equation whether or not RH holds. Away from the line the
two complete contributions must be kept separately.

\subsection{Off-line zeros do not force degenerate classical saddles}
To specify what is meant by a saddle, place the denominator of
\eqref{eq:RS-auxiliary} in the amplitude and choose the phases
\[
 \varphi_+(w;s)=\pi\ii w^2-s\log w,\qquad
 \varphi_-(w;s)=-\pi\ii w^2+(s-1)\log w.
\]
The second is the reflected phase in \eqref{eq:RS-reflection}.
For \(s=\sigma+\ii t\), \(t>0\), their near-positive saddle branches obey
\begin{equation}\label{eq:RS-saddles}
 w_+^2=\frac{t}{2\pi}-\frac{\ii\sigma}{2\pi},\qquad
 w_-^2=\frac{t}{2\pi}+\frac{\ii(1-\sigma)}{2\pi},\qquad
 \varphi_+''(w_+)=4\pi\ii,\quad \varphi_-''(w_-)=-4\pi\ii.
\end{equation}
These formulas follow by differentiating the displayed phases, then
substituting their stationary equations into the second derivatives.
For bounded \(\sigma\), with \(a_t=\sqrt{t/(2\pi)}\), Taylor expansion
of the square root gives
\[
 w_+=a_t-\frac{\ii\sigma}{4\pi a_t}+O(t^{-3/2}),\qquad
 w_-=a_t+\frac{\ii(1-\sigma)}{4\pi a_t}+O(t^{-3/2}).
\]
On the line they are conjugates. Changing \(\sigma\) moves them
smoothly, and the chosen phase Hessians remain nonzero throughout the
high critical strip. The saddle definition is tied to this phase/amplitude
split; it is not an invariant zero detector for an arbitrary integrand
factorization.

If a hypothetical off-line zero \(s_0\) exists, the exact conclusion is
\(\mathscr R(s_0)=-\chi(s_0)\overline{\mathscr R(1-\bar s_0)}\).
After contour displacement this is cancellation of the finite sums,
gamma factors and contour remainders together. It imposes no extra
stationary equation on \(w_\pm\). Saddles can be individually
nondegenerate while their complete contributions cancel. No zero-free
claim can be deduced by inspecting only the leading saddle or by
discarding a remainder near a small value of the total expression.

\subsection{Exact evolution of a complete Riemann--Siegel formula}
Riemann--Siegel is not restricted to \(\lambda=0\). That is the slice
where the classical zeta formula and the prime baseline apply directly.
For \(\lambda>0\), the complete source proves the exact identity
\begin{equation}\label{eq:imaginary-Gaussian}
 H_\lambda(z)=\frac1{\sqrt\pi}\int_\R
       \ee^{-v^2}H_0(z-2\ii\sqrt\lambda\,v)\dd v.
\end{equation}
For \(T>0\), forward heat similarly gives
\begin{equation}\label{eq:real-Gaussian}
 H_{-T}(z)=\frac1{\sqrt\pi}\int_\R
       \ee^{-v^2}H_0(z+2\sqrt T\,v)\dd v.
\end{equation}
To prove both, substitute \eqref{eq:H}, interchange integrals, and use
\(\pi^{-1/2}\int\ee^{-v^2+2av}\dd v=\ee^{a^2}\).
For the imaginary shifts Tonelli's absolute majorant is
\(\int\Phi(u)\ee^{\lambda u^2+|\Im z||u|}\dd u<\infty\).
For real shifts, use the same bound with \(\lambda=0\); the oscillatory
Gaussian integral supplies \(\ee^{-Tu^2}\). Thus these are actual
identities, not formal application of an unbounded reverse heat
operator to arbitrary data.

Write \(M(s)=s(s-1)\pi^{-s/2}\Gamma(s/2)/16\), so that
\(H_0(z)=M(s(z))\zeta(s(z))\). Under the shift in
\eqref{eq:imaginary-Gaussian}, \(s(z)\) becomes \(s(z)-\sqrt\lambda v\).
Consequently a complete finite Riemann--Siegel decomposition can be
evolved by Gaussian integration of its arithmetic terms and its
remainder, provided their individual integrability or a justified
combined continuation is retained. Polymath explicitly carries this
out in the heat-evolved decomposition
\cite[Section 4, equations (35)--(41)]{polymath2019heat}.

The arithmetic weight is visible even before treating the gamma factor:
\begin{equation}\label{eq:heat-Dirichlet-weight}
 \frac1{\sqrt\pi}\int_\R\ee^{-v^2}
       n^{-s+\sqrt\lambda v}\dd v
 =n^{-s}\ee^{\lambda(\log n)^2/4}.
\end{equation}
This is why heat-evolved Riemann--Siegel formulas contain
\(b_n^\lambda=\ee^{\lambda(\log n)^2/4}\). For the actual completed
field the moving \(M(s-\sqrt\lambda v)\), reflected contribution,
cutoff and errors remain. The displayed scalar identity does not
justify multiplying the old cosine formula termwise by \(b_n^\lambda\)
and leaving everything else unchanged. In particular the corresponding
infinite weighted Dirichlet series is not asserted to converge for
positive \(\lambda\).

The published explicit approximation is stated in a definite region:
\(0<\lambda\le1/2\), \(0\le\Im z\le1\), \(\Re z\ge200\), with
gamma corrections, shifted arguments, finite cutoff and error bounds
\cite[Theorem 1.3]{polymath2019heat}. Its exact decomposition has a wider
scope than that asymptotic theorem. A Gaussian average samples points
outside a bounded strip, so a strip-only error estimate cannot be
integrated over the whole Gaussian line without a separate tail payment.
Neither these identities nor the heat equation assign the time-zero
Euler product to \(H_\lambda\) for \(\lambda\ne0\).

\subsection{From heat-density asymptotics to a logarithmic error}
The heat-density investigation uses the same complete source with
forward time \(T=-\lambda\). Its positive-time estimates and its
time-zero Riemann--Siegel estimates have different regimes. For example,
a remainder proved on \(1/\log X\le T\le2\) cannot be specialized to
\(T=0\); the separate time-zero formula must be used. The positive
Newman-time formulas and crossover bounds likewise retain their stated
parameter bands. The source-use and transfer note records these domain
distinctions \cite{GNPHDATransfer}.

At zero Newman time a useful quantitative bridge is available. For
real \(x>0\), put \(s_\pm=1/2\pm\ii x/2\),
\(a_x=\sqrt{x/(4\pi)}\), \(N_x=\lfloor a_x\rfloor\), and
\(\vartheta_x=1-2(a_x-N_x)\). Let
\[
 U_x=\exp\{-\ii[(x/4)\log(x/(4\pi))-x/4-\pi/8]\}.
\]
With the coefficient convention \(C_k\) of the cited expansion, take
\begin{equation}\label{eq:RS-approximant}
 f_K(x)=2\Re\left\{M(s_+)\left[
 \sum_{n\le N_x}n^{-s_+}+(-1)^{N_x-1}U_xa_x^{-1/2}
             \sum_{k=0}^KC_k(\vartheta_x,1/2)a_x^{-k}\right]\right\}.
\end{equation}
The imported explicit remainder, for every \(x>0\) and integer \(K\ge1\), is
\begin{equation}\label{eq:RS-remainder}
 |H_0(x)-f_K(x)|\le\frac{2^{7/4}}7|M(s_+)|
       \Gamma((K+1)/2)1.1^{K+1}a_x^{-K-3/2}.
\end{equation}
This is the normalization recorded and checked in the inherited
transfer, using Polymath's Proposition 6.2
\cite{GNPHDATransfer,polymath2019heat}; the coefficients are those of
that exact version, not a new recurrence asserted here.

The gamma scale cancels in the logarithmic consumer. Indeed
\(D(x)M(s_-)=q(s_-)\), \(|q(s_-)|=1\) and
\(|M(s_+)|=|M(s_-)|\). Thus on the unit interval \(I_v\), with
\(a_- =\sqrt{(v-1/2)/(4\pi)}\), \(D f_K\) approximates \(F\) with
the explicit uniform error
\begin{equation}\label{eq:RS-normalized-error}
 \delta_{v,K}=\frac{2^{7/4}}7
       \Gamma((K+1)/2)1.1^{K+1}a_-^{-K-3/2}.
\end{equation}
The approximant need not be zero-free or holomorphic in a whole disk
for Corollary~\ref{cor:canonical-log}; holomorphy is required of the
exact field. Changes of the integer cutoff do not authorize omitting
the compensating remainder.

The distinction between fixed and increasing expansion order is
quantitative. In \eqref{eq:local-canonical-constants},
\(n_v\sim\log v/\kappa\) and the prefactor of the envelope is
\(O(\log^2v)\). For \(K=2p\), use
\(\Gamma(p+1/2)\le\sqrt\pi\,p^p\) to obtain
\[
 \log(2\delta_{v,2p}/M_v)
 \le-(p+7/4)\log v+O(p\log p).
\]
Choosing \(p=\lceil(4/\kappa)\log\log v\rceil\) therefore gives
\begin{equation}\label{eq:RS-log-transfer}
 \mathcal E_v(\delta_{v,2p})
 =O(\log^2v\,\ee^{-\kappa p+o(1)})=O((\log v)^{-2})
 \qquad(v\to\infty).
\end{equation}
This proves that an increasing-order, complete Riemann--Siegel error
can pass through the deteriorating local logarithmic exponent. A
fixed algebraic source-error rate does not make this particular
sufficient envelope vanish. At fixed height the remainder hierarchy
is asymptotic: the ratio of the bounds at orders \(2p+2\) and \(2p\)
is \((p+1/2)1.1^2/a_-^2\). Sending \(p\to\infty\) at fixed \(v\)
is therefore not a convergence claim. The finite initial-height range
requires its own acquisition, and the arithmetic sum still contains
order \(\sqrt v\) modes.

Equation~\eqref{eq:RS-log-transfer} controls an error allowance. The
remaining time-zero arithmetic task is to bound the signed logarithm of the
retained expression against \(P(X)-G(X,1)\) plus the chosen defect margin, including exceptional
centers. A saddle description, or a remainder tending to zero, does
not supply that comparison. This is the precise boundary between the
background established here and the transition-scale development below.
The time-zero gamma normalization and its pole correction must be
distinguished from the harmonic carrier used there.

\section{The transition-scale Green reduction}\label{sec:transition}

We now keep the complete physical field and choose a quantitative
consumer. This section integrates the child's B01, E01 and E04 results
\cite{GNPRSOuter,GNPRSLogTransfer}; their inherited field input is stated
explicitly below. All constants are finite, but their numerical values
and the first admissible height have not been evaluated.

\subsection{One clock, one carrier, one defect margin}
For a large positive center \(X\), set
\begin{equation}\label{eq:transition-band}
 L=L_X=\log(X/(4\pi)),\quad \lambda=c/L,\quad
 \tfrac12\le c\le1,\quad R_X=\frac{32}{\pi}\log X,
 \quad W_X=[X-R_X,X+R_X].
\end{equation}
The time \(\lambda\) and parameter \(c\) are fixed throughout every
spatial integral at this center. Write \(d\omega_X=K(x-X)\dd x\) on
\(W_X\) and \(q_X=\omega_X(W_X)\). This one-boundary measure is distinct
from the two-boundary probability measure in Section~\ref{sec:green}:
\begin{equation}\label{eq:window-mass}
 0<q_X\le\tfrac12,\qquad
 0\le\tfrac12-q_X\le\frac2\pi X^{-4}.
\end{equation}
Indeed \(K(t)\le\ee^{-\pi|t|/8}/8\), and integration of its two
exterior tails proves the estimate.

Use the upper-half-plane branches in
\begin{align}\label{eq:transition-carrier}
 L(z)&=\Log(z/(4\pi)),\nonumber\\
 \mathcal A(\lambda,z)&=\frac{\pi^{1/4}}{16}z^{7/4}
 \exp\left\{-\frac{\pi z}{8}
       +\frac{\lambda}{16}\left(L(z)^2-\frac{\pi^2}{4}\right)\right\},\nonumber\\
 \Theta(\lambda,z)&=\frac z4(L(z)-1)
                  +\frac{\pi\lambda L(z)}{16}+\frac{7\pi}{8},\qquad
 C_\pm=\tfrac12\mathcal A\ee^{\pm\ii\Theta},\qquad
 V_\lambda=H_\lambda/C_-.
\end{align}
These are explicit nonvanishing carriers on the open strip; they are
not definitions through the unknown zeros. The fixed consumer is
\begin{equation}\label{eq:transition-margin}
 m_*:=m(1/2,1)=\frac{\sin(\pi/8)}{2\cosh(\pi/8)}.
\end{equation}
By Corollary~\ref{cor:margin}, \(J_\lambda(X)<m_*\) excludes a zero
with \(1/2\le\Im z\le1\), \(|\Re z-X|\le1\), at this time.
This fixed-depth consumer is not asserted equivalent to RH.

\begin{lemma}[Harmonic normalization]\label{lem:carrier-harmonic}
The carrier has \(\calJ[C_-]=0\). Its real-boundary logarithm and that
of \(V_\lambda\) are absolutely integrable against \(K\), including
the origin and the negative half-line. Consequently, with
\(I_\lambda(X)=\int_\R K(x-X)\log|V_\lambda(x+\ii0)|\dd x\),
\begin{equation}\label{eq:normalized-green-exact}
 J_\lambda(X)-I_\lambda(X)
 =\int_\R K(x-X)\log|V_\lambda(x+8\ii)|\dd x
       -\log|V_\lambda(X+4\ii)|.
\end{equation}
\end{lemma}
\begin{proof}
The displayed analytic logarithm of \(C_-\) is harmonic in the strip.
At the boundary origin it is \(O(1+|\log|z||^2)\), with gradient
\(O((1+|\log|z||)/|z|)\). Remove an upper semicircle of radius
\(\epsilon\). The Green kernel vanishes linearly there, its normal
derivative is bounded, and the extra Green-formula terms are
\(O(\epsilon(1+\log^2\epsilon))\). On the far vertical sides the
logarithm grows at most linearly plus logarithmically squared, while
harmonic measure decays exponentially. Both contributions vanish.
The boundary modulus, using argument \(\pi\) at negative \(x\), is
\[
 |C_-(\lambda,x+\ii0)|=\frac{\pi^{1/4}}{32}|x|^{7/4}
 \exp\left\{-\frac{\pi|x|}{8}
   +\frac{\lambda}{16}\left(\log^2\frac{|x|}{4\pi}-\frac{\pi^2}{4}\right)\right\}.
\]
Its logarithm has a uniform unit-interval bound of the form in
Lemma~\ref{lem:unit-log}, on both boundaries. Subtracting it from the
complete field's logarithm proves the integrability assertion. Subtract
its zero Green functional from Theorem~\ref{thm:green} to obtain
\eqref{eq:normalized-green-exact}.
\end{proof}

\subsection{The complete field input and the outer baseline}
We use the fixed-buffer heat-density approximation recorded in
\cite{GNPRSHDA,GNPRSOuter}. It is an inherited internally assessed
input, not re-proved from saddle asymptotics here. On the fixed charts
recentered at each large real \(x\), \(0\le y\le8\), put
\[
 N_x=\left\lfloor\sqrt{x/(4\pi)}\right\rfloor,\qquad
 \nu_\lambda(z)=z/2+\pi\lambda/8,\qquad
 w_n(\lambda,z)=n^{-1/2}
   \exp\{\lambda[(\log n)^2-L(z)\log n]/4\}.
\]
The complete approximant is
\begin{equation}\label{eq:transition-provider}
 F_{N_x}(\lambda,z)=\mathcal A(\lambda,z)
       \sum_{n\le N_x}w_n(\lambda,z)
       \cos(\Theta(\lambda,z)-\nu_\lambda(z)\log n)
       +\ee^{-\pi^2\lambda/64}\mathcal R_{N_x}(z),
\end{equation}
where the retained correction is
\begin{align*}
 \mathcal R_N&=(-1)^N\mathcal G(z)a(z)^{-1/2}\mathcal C(a(z)-N),
 &a(z)&=\sqrt{z/(4\pi)},\\
 \mathcal G(z)&=[M(s_+)M(s_-)]^{1/2},
 &\mathcal C(p)&=\frac{\cos(2\pi(p^2-p)-\pi/8)}{\cos(2\pi p)}.
\end{align*}
Here \(M,s_\pm\) are as in Section~\ref{sec:RS}, with compatible
branches. The apparent singularities of \(\mathcal C\) are removable.
The input supplies fixed constants \(h_9,g_9,f_9\) such that, for
\(0\le\lambda L_x\le8\) after its fixed chart onsets,
\begin{align}\label{eq:transition-provider-bounds}
 |H_\lambda-F_{N_x}|&\le h_9x^{-3/4}(|C_+|+|C_-|),\nonumber\\
 |\mathcal G(z)|&\le g_9|\mathcal A(0,z)|/2,
 &|a(z)^{-1/2}\mathcal C(a(z)-N_x)|&\le f_9x^{-1/4}.
\end{align}
The correction bound is uniform in the fractional cutoff. The
fixed-buffer input includes its full Gaussian exterior. The charts
are recentered, never enlarged to radius \(R_X\).

\begin{theorem}[Uniform outer prime baseline]\label{thm:outer-baseline}
There are constants \(C_{\rm out},X_*<\infty\), independent of
\(X,c\) in \eqref{eq:transition-band}, such that
\begin{equation}\label{eq:outer-baseline}
 |J_\lambda(X)-I_\lambda(X)+P_{c,\lambda}(X)|
       \le C_{\rm out}/L_X\qquad(X\ge X_*),
\end{equation}
where, writing \(\Lambda_{\rm vM}\) for the von Mangoldt function,
\begin{equation}\label{eq:transition-prime}
 P_{c,\lambda}(X)=\sum_{n\ge2}\frac{\Lambda_{\rm vM}(n)}{\log n}
 \frac{n^{-5/2-c/4}}{1+n^{-4}}
 \cos\bigl((X/2+\pi\lambda/8)\log n\bigr).
\end{equation}
The omitted integers beyond \(Q\ge2\) cost at most
\(Q^{-3/2-c/4}/(3/2+c/4)\le(2/3)Q^{-3/2}\).
\end{theorem}
\begin{proof}
On \(|x-X|\le R_X\), \(y=4,8\), the eventual gates
\(R_X+8\le X/2\), \(L_X\ge16\) give
\(|L(x+\ii y)-L_X|\le2(R_X+8)/X\) and \(\lambda L_x<2\).
Dividing \eqref{eq:transition-provider} by \(C_-\), the dominant arm is
\[
 D_N(z)=\sum_{n\le N_x}n^{-s_\lambda(z)}
                         \ee^{\lambda(\log n)^2/4},\qquad
 s_\lambda(z)=\frac{1-\ii z}{2}+\frac{\lambda L(z)}4
                               -\frac{\ii\pi\lambda}{8}.
\]
The reflected carrier satisfies \(|\ee^{2\ii\Theta}|\ll x^{-y/2}\).
Since \(|w_n|\le n^{-1/2}\), summing the entire reflected population
costs \(O(x^{(1-y)/4})\). The normalized correction costs
\(O(x^{-(y+1)/4})\), and the remaining normalized error costs
\(O(x^{-3/4})\), by \eqref{eq:transition-provider-bounds}.

Let \(D_j=\sum_{n\ge2}n^{-5/2}(\log n)^j\), \(j=0,1,2\).
On these outer lines, \(\log n\le\Re L(z)/2\), so
\(|n^{-s_\lambda(z)}\ee^{\lambda(\log n)^2/4}|\le n^{-5/2}\).
Using \(\ee^u-1\le u\ee^u\) pays removal of the heat weight by
\(\lambda D_2/4\). The ordinary Dirichlet-series tail costs
\((2/3)N_x^{-3/2}=O(x^{-3/4})\). Replacing \(s_\lambda(z)\) by
\[
 s_y^0(x)=\frac{1+y}{2}+\frac c4
                   -\ii\left(\frac x2+\frac{\pi\lambda}{8}\right)
\]
costs at most \(D_1\lambda(R_X+8)/(2X)\). These uses of the
ordinary zeta series are all in \(\Re s\ge5/2\)
\cite{DLMFZetaOuter2026}. Thus a fixed \(C_F\) gives
\begin{equation}\label{eq:outer-field-comparison}
 |V_\lambda(x+\ii y)-\zeta(s_y^0(x))|\le C_F/L_X,
 \qquad y=4,8,\quad |x-X|\le R_X.
\end{equation}
No such absolutely convergent comparison has been used at \(y=0\).

The independent floor
\(|\zeta(s)|\ge1-D_0>1/2\) follows from
\(D_0\le7/(12\sqrt2)<1/2\). Once \(L_X>4C_F\), the logarithmic
error in \eqref{eq:outer-field-comparison} is at most \(4C_F/L_X\).
The unit-log bound in Lemma~\ref{lem:carrier-harmonic}, summed in
unit shells outside \(W_X\), gives on either boundary
\begin{equation}\label{eq:transition-log-tail}
 \int_{\R\setminus W_X}K(x-X)|\log|V_\lambda(x+\ii y)||\dd x
 \le T_X:=D_VX^{-3}\log^2X,\qquad y=0,8.
\end{equation}
This is an estimate for the exact logarithm. The comparison zeta log
has a separate exterior bound \(D_ZX^{-4}\), since its absolute
prime-power series is bounded by \(D_0\).

Finally \(\widehat K(\omega)=1/[2\cosh(4\omega)]\), so at
\(\omega=(\log n)/2\) it equals \(1/(n^2+n^{-2})\).
The upper-boundary term less the center term has coefficient
\[
 \frac{n^{-9/2-c/4}}{n^2+n^{-2}}-n^{-5/2-c/4}
       =-\frac{n^{-5/2-c/4}}{1+n^{-4}}.
\]
Absolute convergence permits termwise integration. Equation
\eqref{eq:normalized-green-exact} now has error at most
\(6C_F/L_X+T_X+D_ZX^{-4}\), proving the assertion after enlarging
the onset and taking \(C_{\rm out}=6C_F+D_V+D_Z\).
The prime tail follows from the integral test.
\end{proof}

The constants \(C_{\rm out}\) and \(D_V\) denote the child's
\(C_*\) and \(D_V\); this avoids collision with the earlier
unit-log constant. The onset includes the inherited fixed-buffer
onsets, \(X\ge26\), \(R_X+8\le X/2\),
\(L_XX^{-3/4}\le1\), \(L_X>4C_F\), and
\(L_XX^{-3}\log^2X\le1\). Each gate is eventual after the constants
are fixed. The result holds at every sufficiently large center, but
estimates the outer reduction error, not the sign of \(I-P\).
There is no extra time-zero pole correction in
\eqref{eq:outer-baseline}: the carrier here is harmonic and nonvanishing.

\subsection{Logarithmic transfer without a real-line floor}
On the real window define the two retained arms by
\begin{equation}\label{eq:real-two-arms}
 A_\lambda(x)=\sum_{n\le N_x}w_n(\lambda,x)
                  \ee^{\ii\nu_\lambda(x)\log n},\qquad
 S_\lambda=A_\lambda+\ee^{\ii\phi_\lambda}\overline{A_\lambda},
 \quad\phi_\lambda=2\Theta(\lambda,x).
\end{equation}
The exact field is \(V_\lambda=S_\lambda+e_\lambda\), with
\begin{equation}\label{eq:real-provider-error}
 \sup_{W_X}|e_\lambda|\le C X^{-1/4},\qquad
 C=2^{3/4}(2h_9)+2^{1/4}g_9f_9.
\end{equation}
Indeed the two carriers have equal modulus, the full correction costs
\(g_9f_9x^{-1/4}\), and the provider remainder costs
\(2h_9x^{-3/4}\). Thus the correction is paid in \(e_\lambda\),
not discarded when it is absent from \(S_\lambda\).

\begin{lemma}[One-sided logarithmic transfer]\label{lem:upper-log-transfer}
For \(v,s\in\C\) and \(\delta>0\),
\begin{equation}\label{eq:upper-log-transfer}
 \log|v|\le\tfrac12\log(|s|^2+\delta^2)
       +\tfrac12\log(1+|v-s|^2/\delta^2).
\end{equation}
\end{lemma}
\begin{proof}
Use \(|v|\le|s|+|v-s|\) and Cauchy--Schwarz on
\((|s|,\delta)\) and \((1,|v-s|/\delta)\), then take logarithms.
At \(v=0\) the inequality is understood with \(\log0=-\infty\).
\end{proof}
Combining this lemma with \eqref{eq:real-provider-error}, the exact
tail and Theorem~\ref{thm:outer-baseline}, and choosing
\(\delta=X^{-1/8}\), gives
\begin{align}\label{eq:full-regularized-transfer}
 J_\lambda(X)\le{}&\int_{W_X}\tfrac12
           \log(|S_\lambda|^2+X^{-1/4})\,d\omega_X-P_{c,\lambda}(X)
            +T_X+\frac{C_{\rm out}}{L_X}\nonumber\\
 &+\frac{q_X}{2}\log(1+C^2X^{-1/4}).
\end{align}
The field-transfer cost is \(O(X^{-1/4})\). Neither the exact field
nor the reflected approximant is assumed bounded away from zero.

\subsection{A shifted-zeta comparison retains its unknown divisor}
For comparison with a single ordinary zeta function, let
\(\eta=\pi\lambda/4\) and
\(Z_{c,\lambda}(z)=\zeta((1-\ii z)/2+c/4-\ii\pi\lambda/8)\).
The exact time-zero normalization gives
\[
 M((1-\ii z)/2+c/4-\ii\pi\lambda/8)Z_{c,\lambda}(z)
       =H_0(z+\eta+\ii c/2).
\]
In the strip the zeta factor has a simple pole
\(p=-\eta+\ii(1-c/2)\); its gamma completion has precisely the
cancelling simple zero there. If \(x_j(0)+\ii y_j(0)\) are the
time-zero zeros, with multiplicity, set
\[
 \mathcal Z_{c,\lambda}(X)=\sum_{y_j(0)>c/2}
       G(X+\eta-x_j(0),y_j(0)-c/2),\qquad
 \Delta=I_\lambda-\int_\R K(x-X)\log|Z_{c,\lambda}(x)|\dd x.
\]
The complete-product proof of Theorem~\ref{thm:green} applies to this
fixed affine translate: local distance logarithms remain integrable,
the count bound is unchanged, and the boundary tails remain summable.
The pole enters with negative sign. On the two outer lines the zeta
combination is exactly \(-P_{c,\lambda}\); therefore
\begin{equation}\label{eq:shifted-zeta-divisor}
 I_\lambda-P_{c,\lambda}
    =\Delta+\mathcal Z_{c,\lambda}-G(X+\eta,1-c/2).
\end{equation}
This is B02's meromorphic identity \cite{GNPRSComparisons}. The
difference \(\Delta\) is a difference of integrable logarithms, not
an assumed holomorphic logarithm of a zero-free quotient. A useful
upper bound still requires a coupled estimate for
\(\Delta+\mathcal Z_{c,\lambda}\). The identity alone supplies none;
the isolated pole term decays exponentially with \(X\).

\section{Coherent moments and the limits of quadratic acquisition}
\label{sec:coherent-moments}

The logarithmic transfer preserves cancellations. Replacing its
integral by a quadratic moment loses information in two stages: the
oscillating reflected phase and the variation of the arithmetic arm.
We first compute the exact coherent moment, then quantify both losses.
The statements in this section are the E03 and B03 results
\cite{GNPRSMoment,GNPRSPhaseObstruction}.

\subsection{The exact local moment law}
Set \(M(X,c)=\int_{W_X}|V_{c/L_X}(x)|^2\,d\omega_X\).
\begin{theorem}[Coherent moment at \(c=1\)]\label{thm:coherent-moment}
For the complete field, with time fixed in each defining integral,
\begin{equation}\label{eq:coherent-moment}
 \lim_{T\to\infty}\frac1T\int_T^{2T}M(X,1)\dd X=\zeta(3/2).
\end{equation}
\end{theorem}
\begin{proof}
Introduce only for the height average
\(v(u)=V_{1/L_u}(u)\). Its auxiliary arm weights and phases are
\[
 w_n(u)=n^{-3/4}\ee^{(\log n)^2/(4L_u)}\le n^{-5/8},\quad
 \nu(u)=u/2+\pi/(8L_u),\quad
 \phi(u)=u(L_u-1)/2+\pi/8+7\pi/4.
\]
Extend weights by zero above \(N_u\). Each arm is \(O(u^{3/16})\).
By \eqref{eq:real-provider-error}, replacing \(|v|^2\) by the
square of the two arms costs \(O(u^{-1/16}+u^{-1/2})\).
The diagonal of the two arm squares is \(2\sum_{n\le N_u}w_n(u)^2\).
For every fixed \(n\), its summand tends uniformly on \([T,2T]\)
to \(2n^{-3/2}\); the dominating series \(2\sum n^{-5/4}\)
is summable. Thus this diagonal tends uniformly to \(2\zeta(3/2)\).

For every pair, integrate on its actual activation interval
\[
 I_{mn}=[T,2T]\cap[4\pi\max(m,n)^2,\infty).
\]
There \(w_mw_n\) decreases, and its supremum plus total variation
is at most \(2(mn)^{-5/8}\). A phase with monotone derivative of
fixed sign and modulus at least \(\alpha\) has every partial
integral bounded by \(3/\alpha\): integration by parts pays two
endpoints and the variation of the reciprocal derivative. Applying
this bound through the amplitude primitive pays the stated variation.

For same-arm pairs \(m\ne n\), eventually \(\nu'\ge1/4\) and
\(\nu''>0\). The cost is
\(O((mn)^{-5/8}/|\log(n/m)|)\). For \(m\le n/2\) use
\(|\log(n/m)|\ge\log2\); for \(n/2<m<n\) use
\(\log(n/m)\ge(n-m)/n\). Summing through \(N_{2T}\ll\sqrt T\)
gives an averaged error \(O(T^{-5/8}(1+\log T))\).

The cross-arm phases are \(\phi(u)-\nu(u)\log(mn)\), including
\(m=n\). On their activation intervals,
\[
 \frac1{8T}\le(\phi-\nu\log(mn))''\le\frac1{2T}
\]
eventually. The part where the first derivative has modulus at most
\(T^{-1/2}\) has length at most \(16\sqrt T\). Its complement
has at most two pieces, each costing \(3\sqrt T\). This gives
\(22\sqrt T\) for every partial integral, including a stationary
point at a cutoff endpoint. The complete averaged cross-arm cost is
\[
 O\left(T^{-1/2}\left(\sum_{n\le N_{2T}}n^{-5/8}\right)^2\right)
       =O(T^{-1/8}).
\]
We have proved \(T^{-1}\int_T^{2T}|v(u)|^2\dd u\to2\zeta(3/2)\).

It remains to restore the original clock and convolution. For
\(x\in W_X\), \(|1/L_X-1/L_x|\le4R_X/(XL_X^2)\).
At intermediate times \(1/2\le\lambda L_x\le2\), the common
weight bound is \(n^{-9/16}\); hence
\(S_\lambda=O(x^{7/32})\) and
\(\partial_\lambda S_\lambda=O(L_x^2x^{7/32})\).
The squared-arm comparison costs \(O(R_XX^{-9/16})\); paying the
full correction at both endpoint times adds \(O(X^{-1/32})\).
Thus it changes the local moment by a uniformly vanishing quantity.
No time-dependent field has been substituted without this payment.

Fubini now gives an auxiliary weight
\[
 W_T(x)=\int_T^{2T}K(x-X)\mathbf1_{|x-X|\le R_X}\dd X.
\]
It lies in \([0,1/2]\), and on
\([T+R_{2T},2T-R_{2T}]\) it is at least \(q_T\).
The remaining edge intervals have total length at most \(4R_{2T}\).
Since \(|v|^2=O(T^{3/8})\) there, replacing \(W_T\) by
\(\tfrac12\mathbf1_{[T,2T]}\) costs
\(O(T^{-29/8}+T^{-5/8}\log T)\) in the average.
The kernel mass \(1/2\) proves \eqref{eq:coherent-moment}.
\end{proof}

\subsection{An upper allowance that cannot succeed at every center}
Jensen's inequality applied with probability measure
\(d\omega_X/q_X\), together with the exact exterior, gives
\begin{equation}\label{eq:quadratic-allowance}
 J_\lambda(X)\le\mathcal B_J(X,c):=
 \frac{q_X}{2}\log\frac{M(X,c)}{q_X}
      -P_{c,\lambda}(X)+T_X+\frac{C_{\rm out}}{L_X}.
\end{equation}
Here \(M>0\), since an entire nonzero field cannot vanish on an
interval. The following failure concerns this upper allowance.

\begin{corollary}[Exact quadratic-Jensen obstruction]\label{cor:Jensen-stop}
There are arbitrarily large \(X\) such that
\(\mathcal B_J(X,1)\ge m_*\). The same centers fail if the exact
moment is replaced by any valid upper bound for it.
\end{corollary}
\begin{proof}
The elementary bounds needed are
\[
 m_*<\frac9{50},\qquad |P_{1,\lambda}|<\frac{23}{100},\qquad
 \zeta(3/2)>\frac{13}{5},\qquad \ee^{41/25}<\frac{26}{5}.
\]
For the first, use \(\sin(\pi/8)<77/200\) and
\(\cosh(\pi/8)>137/128\). For the second, the prime powers through
16 are \(2,3,4,5,7,8,9,11,13,16\); respective rational upper
budgets for their coefficients are
\(.14,.049,.012,.012,.005,.0014,.0013,.0014,.001,.00013\).
They sum to \(.22323\), and the remaining integer tail is at most
\(1/224<.0045\). These rational budgets follow by raising the
elementary fourth-root comparisons to the fourth power, as detailed
in \cite{GNPRSMoment}. Convexity of \(u^{-3/2}\) gives
\(\zeta(3/2)\ge1+1/(2\sqrt2)+13/(6\sqrt3)>177/68>13/5\).
Finally \(\ee^a\le1+a+a^2/[2(1-a/3)]\), \(0\le a<3\), follows
from \(k!\ge2\cdot3^{k-2}\); apply it at \(a=41/100\) and
raise to the fourth power.

If the allowance were below \(m_*\) eventually, dropping its
nonnegative errors would give
\[
 M(X,1)<q_X\exp\{(41/50)/q_X\}\longrightarrow
                  \tfrac12\ee^{41/25}<13/5.
\]
This eventual bound contradicts Theorem~\ref{thm:coherent-moment}.
Monotonicity of the logarithm proves the assertion for larger moments.
\end{proof}

\subsection{Why the ideal cosine correction is still insufficient}
For an independent uniform phase \(\vartheta\) and \(A\ne0\),
\[
 \mathbb E_\vartheta\log|2A\cos\vartheta|=\log|A|,
 \qquad \mathbb E_\vartheta|2A\cos\vartheta|^2=2|A|^2.
\]
The logarithmic identity follows from the sine/cosine doubling
formula with its integrable zero logarithms. If \(A\) also varies
under \(d\omega_X/q_X\), the product-average Jensen gap is
\begin{equation}\label{eq:phase-amplitude-gap}
 \frac{q_X}{2}\log2+
 q_X\left(\frac12\log\mathbb E_X|A|^2-\mathbb E_X\log|A|\right).
\end{equation}
The second term is nonnegative. This identity assumes an independent
phase average, not the actual spatial phase. Even granting its first
term for free does not repair the quadratic method.

\begin{theorem}[Phase-only correction obstruction]\label{thm:phase-stop}
Define the diagnostic
\(\mathcal B_{\rm ph}=\mathcal B_J-(q_X/2)\log2\).
There are arbitrarily large \(X\) with
\begin{equation}\label{eq:phase-stop}
 \mathcal B_{\rm ph}(X,1)\ge m_*+1/1000.
\end{equation}
No assertion that \(\mathcal B_{\rm ph}\) bounds \(J\) is needed
or made. The failure persists for any upper replacement of \(M\).
\end{theorem}
\begin{proof}
At \(c=1\), write \(P_0(X)=P_{1,0}(X)=\sum_p F_p((X/2)\log p)\),
where
\[
 F_p(\theta)=\sum_{k\ge1}
    \frac{p^{-11k/4}\cos(k\theta)}{k(1+p^{-4k})},\qquad
 b_p=\sum_{k\ge1}\frac{p^{-11k/4}}{k(1+p^{-4k})}.
\]
Each block has mean zero and lies in \([-b_p,b_p]\).
For a finite set of primes the height average of \(\exp(4\sum F_p)\)
equals the product circle average. Indeed the absolutely convergent
Fourier expansion has nonzero frequency
\(\tfrac12\sum j_p\log p\) for every nonconstant mode, by unique
factorization, and such a mode has vanishing dyadic mean.
The bound \(\sum b_p<\infty\) permits removal of the prime tail
uniformly. Convexity's chord bound and \(\cosh u\le\ee^{u^2/2}\)
therefore give
\begin{equation}\label{eq:prime-exponential-mean}
 \mathfrak E:=\lim_{T\to\infty}\frac1T\int_T^{2T}\ee^{4P_0(X)}\dd X
 \le\exp\left(8\sum_p b_p^2\right)<\ee^{27/125}.
\end{equation}
For completeness the last constant uses
\[
 b_2<77/500,\quad b_3<51/1000,\quad
 \sum_{p\ge5}b_p^2<\frac{2432}{7021125},\quad
 \sum_p b_p^2<\frac{27}{1000}.
\]
The first coefficients at 2 and 3 are below \(.14,.049\).
Writing \(r=p^{-11/4}\), the higher powers cost at most
\(r^2/[2(1-r)]\), with \(r_2<3/20,r_3<1/20\).
For \(p\ge5\), \(b_p<(80/79)p^{-11/4}\); sum its square over
integers using
\(\sum_{n\ge5}n^{-11/2}<1/6250+1/5625\).
This proves the displayed rational bounds.

Suppose \eqref{eq:phase-stop} fails eventually. Rearranging gives
\[
 M(X,1)<2q_X\exp\{2(m_*+1/1000+P_{1,1/L_X}(X))/q_X\}.
\]
The differentiated absolutely convergent prime series shows
\(P_{1,1/L_X}-P_0=O(1/L_X)\). Thus the right side is
\(\ee^{4(m_*+1/1000)}\ee^{4P_0(X)}+o(1)\), uniformly in \(X\).
Taking means and using \eqref{eq:prime-exponential-mean} forces
\[
 \zeta(3/2)<\ee^{47/50}
       \le\frac{26609}{10300}<\frac{13}{5},
\]
contrary to the preceding lower bound. The exponential bound uses
the same factorial inequality at \(47/50\).
\end{proof}

These obstructions locate information lost by the specified quadratic
expressions. They do not lower-bound the true logarithm or the Green
potential at the failing centers. In particular, no collision exclusion
or obstruction to all signed nonlinear estimates follows.

\subsection{Exact regularized phase and the fixed-prefix limit}
For any complex \(A\), put
\[
 B_\delta(A)=\frac{\delta+\sqrt{\delta^2+4|A|^2}}2,
 \qquad r_\delta(A)=\frac{\overline A^{,2}}{B_\delta(A)^2}.
\]
Since \(B_\delta-|A|^2/B_\delta=\delta\), direct expansion gives
\begin{equation}\label{eq:regularized-phase-factor}
 \tfrac12\log(|A+\ee^{\ii\phi}\overline A|^2+\delta^2)
 =\log B_\delta(A)+\Re\log(1+\ee^{\ii\phi}r_\delta(A)).
\end{equation}
Here \(|r_\delta|<1\), including \(A=0\), so the second logarithm
has its convergent power series. Its independent phase mean is zero.

For a fixed prefix \(A_{\lambda,M}\) and fixed \(\delta>0\),
\(|r_\delta|\le\rho<1\) and \(|\partial_x r_\delta|=O_{M,\delta}(1)\),
while \(\phi_\lambda'\ge L_X/4\) and
\(\phi_\lambda''=O(X^{-1})\). Integration by parts in the \(k\)-th
harmonic, including both endpoints, bounds its integral by
\[
 \frac{\rho^k(2\|K\|_\infty+\|K'\|_1)
       +k\rho^{k-1}q_X\|r_\delta'\|_\infty}{k\min\phi_\lambda'}
 +\frac{q_X\rho^k\|\phi_\lambda''\|_\infty}
       {k(\min\phi_\lambda')^2}.
\]
Summing with the logarithmic coefficient \(1/k\) gives
\(O_{M,\delta}(1/L_X)\). If
\(E_M=\int_{W_X}|V_\lambda-S_{\lambda,M}|^2\,d\omega_X\),
Lemma~\ref{lem:upper-log-transfer} consequently yields
\begin{align}\label{eq:prefix-transfer}
 J_\lambda(X)\le{}&\int_{W_X}\log B_\delta(A_{\lambda,M})\,d\omega_X
       -P_{c,\lambda}(X)
       +\frac{q_X}{2}\log\left(1+\frac{E_M}{q_X\delta^2}\right)
       \nonumber\\
 &+T_X+\frac{C_{\rm out}+C_{M,\delta}}{L_X}.
\end{align}
Restricting the moment proof to modes \(n>M\) gives exactly
\(T^{-1}\int_T^{2T}E_M(X,1)\dd X\to\sum_{n>M}n^{-3/2}\).
The tail is positive for each fixed prefix. This limits a vanishing
residual claim but does not decide the fixed defect margin for a
chosen prefix and regulator.

The fixed-prefix phase constants cannot be used for the full growing
cutoff. Already the relative phase at \(m=n=N_x\) has derivative
\((L_x-2\log N_x)/2+\pi\lambda/(8x)=O(x^{-1/2})\).
A large carrier derivative alone does not bound this derivative away
from zero. The next section chooses a different reference for which
the complete phase can be controlled.

\section{Euler references, density and exceptional centers}
\label{sec:euler-reference}

A finite Euler reference has two advantages over a fixed arithmetic
prefix: its logarithm calibrates exactly to a prime subseries, and
its modulus never vanishes. The reflected reference does have real
zeros; these will be retained inside integrable logarithms. The
development combines E05's fixed-reference moment analysis and E06's
joint-phase estimate \cite{GNPRSEuler,GNPRSAdaptive}.

\subsection{Exact calibration and a regulator-uniform phase lemma}
For \(P\ge2\), fixed across the spatial integral, define
\begin{equation}\label{eq:euler-reference}
 E_P(x)=\prod_{p\le P}(1-p^{-\sigma_c}
                   \ee^{\ii\nu_\lambda(x)\log p})^{-1},\qquad
 \sigma_c=\tfrac12+c/4,\qquad
 S_P^{\rm ref}=E_P+\ee^{\ii\phi_\lambda}\overline E_P.
\end{equation}
Use the canonical logarithm
\(\ell_P=\sum_{p\le P}\sum_{k\ge1}
 p^{-k\sigma_c}\ee^{\ii k\nu_\lambda\log p}/k\).
The finite product and its inverse are bounded on the entire real
line for fixed \(P\). Absolute convergence and the Fourier transform
of \(K\) give
\begin{equation}\label{eq:euler-calibration}
 \int_\R K(x-X)\log|E_P(x)|\dd x
 =P_{c,\lambda}^{(\le P)}(X)
 :=\sum_{p\le P}\sum_{k\ge1}
 \frac{p^{-k(5/2+c/4)}\cos(k\nu_\lambda(X)\log p)}{k(1+p^{-4k})}.
\end{equation}
All powers of the retained primes are included. The omitted prime
subseries costs at most \((2/3)P^{-3/2}\), since every omitted
prime power exceeds \(P\).

Put \(\kappa_E=1-2^{-5/8}>0\). Differentiating the absolutely
convergent logarithm, using \(\nu_\lambda'=1/2\), yields
\begin{align}\label{eq:euler-phase-derivatives}
 2|\Im\ell_P'|&\le D_1(P):=\frac8{3\kappa_E}P^{3/8}\log P,\nonumber\\
 2|\Im\ell_P''|&\le D_2(P):=\frac4{3\kappa_E^2}P^{3/8}(\log P)^2.
\end{align}
For example the first sum is bounded by
\(\sum_{p\le P}p^{-5/8}\log p/(1-p^{-5/8})\); the second has
coefficient \(1/2\) and denominator squared. In both, use
\(\sum_{n\le P}n^{-5/8}\le(8/3)P^{3/8}\).

\begin{lemma}[The complete reference phase]\label{lem:joint-phase}
Let \(\psi=\phi_\lambda-2\Im\ell_P\), so that
\(a:=S_P^{\rm ref}/E_P=1+\ee^{\ii\psi}\). Assume the eventual
window gates and \(D_1(P)\le L_X/8\), \(R_X/X\le L_X/8\).
Then \(\psi'\ge L_X/4\) on \(W_X\), and, writing
\(D=\|\psi''\|_{\infty,W_X}\),
\begin{align}\label{eq:joint-phase}
 D&\le D_2(P)+X^{-1}+\pi\lambda/(2X^2),\nonumber\\
 \left|\int_{W_X}\tfrac12\log(|a|^2+\delta^2)\,d\omega_X
                   -q_X\log b_\delta\right|
 &\le\frac2{L_X}+\frac{16D}{L_X^2},
 \qquad b_\delta=\frac{\delta+\sqrt{\delta^2+4}}2.
\end{align}
The bound holds uniformly for every \(\delta>0\) fixed within
the integral, and also at \(\delta=0\), with \(\log b_0=0\).
\end{lemma}
\begin{proof}
At fixed time \(\phi_\lambda'=L_x/2+\pi\lambda/(8x)\) and
\(\phi_\lambda''=1/(2x)-\pi\lambda/(8x^2)\). The derivative
bounds and \(L_x\ge L_X-2R_X/X\) prove the assertions about \(\psi\).
The unit-arm instance of \eqref{eq:regularized-phase-factor} is
\[
 \tfrac12\log(|a|^2+\delta^2)=\log b_\delta+
  \Re\sum_{k\ge1}\frac{(-1)^{k+1}\rho^k}{k}\ee^{\ii k\psi},
 \qquad \rho=b_\delta^{-2}<1.
\]
Its coefficient modulus is constant in \(x\). Integration by parts
in the entire phase, with \(m=\min\psi'\), gives
\[
 \left|\int_{W_X}\ee^{\ii k\psi}\,d\omega_X\right|
 \le\frac{2\|K\|_\infty+\|K'\|_1}{km}
          +\frac{q_XD}{km^2}.
\]
Here \(\|K\|_\infty=1/16\), \(\|K'\|_1=1/8\) and
\(q_X\le1/2\). Summing with \(\rho^k/k\) and
\(\sum\rho^k/k^2\le2\) proves \eqref{eq:joint-phase}.
For \(0<\delta\le1\), the regularized log decreases to
\(\log|a|\) and is bounded above by \(\tfrac12\log5\).
Monotone convergence applied to this constant minus the log, and
the uniform two-sided integral bound, give a finite limit with the
same estimate. Thus the real reference zeros are included.
\end{proof}

Define the exact relative residual
\begin{equation}\label{eq:exact-relative-residual}
 R_P(X,c)=\int_{W_X}\frac{|V_\lambda-S_P^{\rm ref}|^2}{|E_P|^2}
                   \,d\omega_X.
\end{equation}
Only the unreflected nonvanishing product is divided out. Also
\(|\ell_P|\le(8/(3\kappa_E))P^{3/8}\), so its omitted exterior
costs at most
\(B_P^{\rm ext}X^{-4}\), where
\(B_P^{\rm ext}=16P^{3/8}/(3\pi\kappa_E)\).
Applying Lemma~\ref{lem:upper-log-transfer} to \(V/E_P,a\), then
Jensen only to the residual payment, proves the complete bound
\begin{align}\label{eq:euler-master}
 J_\lambda(X)\le{}&P_{c,\lambda}^{(\le P)}-P_{c,\lambda}
       +q_X\log b_\delta
       +\frac{q_X}{2}\log\left(1+\frac{R_P(X,c)}{q_X\delta^2}\right)
       \nonumber\\
 &+T_X+\frac{C_{\rm out}+2}{L_X}
       +\frac{16D}{L_X^2}+B_P^{\rm ext}X^{-4}.
\end{align}
No field floor and no omitted true logarithmic exterior enter this
inequality. For fixed \(P\), its phase error is \(O_P(1/L_X)\)
uniformly over regulators, strengthening the older fixed-regulator
form of E05 by E06's phase lemma.

\subsection{The fixed-reference residual mean and density gain}
\begin{theorem}[Fixed-prime residual law]\label{thm:residual-mean}
For each fixed finite \(P\),
\begin{equation}\label{eq:residual-mean}
 \lim_{T\to\infty}\frac1T\int_T^{2T}R_P(X,1)\dd X
 =\mu_P:=\zeta(3/2)\prod_{p\le P}(1-p^{-3/2})-1,
 \qquad \mu_P\longrightarrow0\quad(P\to\infty).
\end{equation}
The height limit is taken before the prime-cutoff limit.
\end{theorem}
\begin{proof}
At the auxiliary time \(1/L_u\), let \(D_P=\prod_{p\le P}p\) and
\[
 Q(u)=E_P(u)^{-1}=\sum_{d\mid D_P}\mu(d)d^{-3/4}
                  \ee^{\ii\nu(u)\log d},\qquad B(u)=A(u)Q(u)-1,
\]
where \(\mu(d)\) is the M\"obius function. The coefficient at 1
vanishes. For \(k>1\), the coefficient of
\(\ee^{\ii\nu(u)\log k}\) is
\[
 b_k(u)=\sum_{\substack{d\mid D_P,\ d\mid k\\ k/d\le N_u}}
                     \mu(d)d^{-3/4}w_{k/d}(u).
\]
Its supremum plus total variation on \([T,2T]\) is
\(O_P(k^{-5/8})\), including each cutoff activation. For fixed
\(k>1\) it tends uniformly to \(k^{-3/4}\mathbf1_{(k,D_P)=1}\).
The diagonal thus tends to
\(\sum_{k>1,(k,D_P)=1}k^{-3/2}=\mu_P\), by the summable square
majorant \(O_P(k^{-5/4})\). The same-arm argument in
Theorem~\ref{thm:coherent-moment}, now through \(D_PN_{2T}\),
pays all off-diagonal terms by \(O_P(T^{-5/8}(1+\log T))\).

The reflected relative square equals
\[
 2|B|^2+2\Re\{\ee^{-\ii\phi}B^2\overline Q/Q\}.
\]
The finite-prime function \(\overline Q/Q=\overline Q E_P\) has
an absolutely uniformly convergent expansion in rational frequencies
\(\ee^{\ii\nu\log r}\). Truncate it with uniform error
\(\epsilon\). For each of the finitely retained \(r\), the phase
\(-\phi+\nu\log(klr)\) has second derivative of fixed sign and
magnitude \(\asymp1/T\), since \(k,l\ll_P\sqrt T\).
The previous second-derivative argument costs \(O_{P,r}(T^{-1/8})\)
after summation. The removed frequencies cost at most
\(2\epsilon\) times the bounded mean of \(|B|^2\).
Take \(T\to\infty\), then \(\epsilon\downarrow0\); the cross
mean is zero.

Dividing the complete field error by \(E_P\) costs
\(O_P(u^{-1/4})\), hence \(O_P(u^{-1/16})\) in the square.
The auxiliary mean is therefore \(2\mu_P\).
The fixed-clock comparison has the same
\(O_P(R_XX^{-9/16}+X^{-1/32})\) bound as before, because the fixed
product and its inverse have bounded modulus and time derivative;
the reflected phase derivative is already included in the arm bound.
The same Fubini edges and kernel mass \(1/2\) prove
\eqref{eq:residual-mean}. Finally \(\mu_P\to0\) by domination in
the displayed coprime sum.
\end{proof}

\begin{corollary}[Density of centers with a large Green value]
\label{cor:green-density}
For every fixed \(h>0\),
\begin{equation}\label{eq:green-density}
 \lim_{T\to\infty}\frac1T
   |\{X\in[T,2T]:J_{1/L_X}(X)\ge h\}|=0.
\end{equation}
In particular this holds at the original defect margin \(h=m_*\).
There is no rate or assertion about a fixed time, all \(c\), or
the density of zeros.
\end{corollary}
\begin{proof}
Choose a fixed \(\delta>0\) so \(q_X\log b_\delta\le h/5\),
and then arbitrarily large fixed \(P\) so the prime tail is at most
\(h/5\). Set \(r=\delta^2h/5\). If \(R_P\le r\), its
logarithmic payment is at most \(r/(2\delta^2)=h/10\).
The remaining errors in \eqref{eq:euler-master} are eventually at
most \(h/5\) for each fixed \(P\). Hence the bad set is contained
in \(\{R_P>r\}\). Markov's inequality and
Theorem~\ref{thm:residual-mean} bound its upper dyadic density by
\(\mu_P/r\). Let fixed \(P\) tend to infinity after the height
limit. This proves the statement for arbitrary \(h\), explicitly
extracting the arbitrary-threshold consequence of E05's proof at
\(m_*\).
\end{proof}

\subsection{An averaged consequence for the actual Green value}
The density conclusion can be strengthened in mean by the coherent
moment bound. The following consequence is derived in this synthesis
\cite{GNPRSMeanConsequences}; it adds no exceptional-center arithmetic
estimate to E05.

\begin{theorem}[Power means and subcritical exponential means]
\label{thm:green-power-means}
Write \(J(X)=J_{1/L_X}(X)\). For all sufficiently large \(T\),
\begin{equation}\label{eq:green-exponential-bound}
 \frac1T\int_T^{2T}\ee^{4J(X)}\dd X\le39.
\end{equation}
For every fixed \(p>0\) and every fixed \(0<\alpha<4\),
\begin{equation}\label{eq:green-power-means}
 \frac1T\int_T^{2T}J(X)^p\dd X\longrightarrow0,
 \qquad
 \frac1T\int_T^{2T}\ee^{\alpha J(X)}\dd X\longrightarrow1.
\end{equation}
The constant 39 is an explicit convenient bound; its finite onset is
unevaluated. Neither convergence rate nor convergence at \(\alpha=4\)
is asserted.
\end{theorem}
\begin{proof}
Eventually \(q_X\ge1/4\) and
\(T_X+C_{\rm out}/L_X\le1/100\). Using
\(|P_{1,1/L_X}|<23/100\) in \eqref{eq:quadratic-allowance} gives
\[
 \ee^{4J(X)}\le\ee^{24/25}
       \left(\frac{M(X,1)}{q_X}\right)^{2q_X}
       \le\ee\,(1+4M(X,1)).
\]
The last inequality uses \(t^r\le1+t\) for \(t\ge0\),
\(0<r\le1\). Also
\(\zeta(3/2)\le1+2^{-3/2}+\int_2^\infty u^{-3/2}\dd u<3\).
Theorem~\ref{thm:coherent-moment} thus makes the averaged right side
at most \(13\ee<39\) eventually.

Let \(d_T(t)=T^{-1}|\{X\in[T,2T]:J(X)>t\}|\).
Positivity of the Green representation and Markov's inequality give
\(d_T(t)\le\min(1,39\ee^{-4t})\). For each fixed \(t>0\),
Corollary~\ref{cor:green-density} gives \(d_T(t)\to0\).
Tonelli's theorem yields
\[
 \frac1T\int_T^{2T}J^p\dd X
     =p\int_0^\infty t^{p-1}d_T(t)\dd t,
 \qquad
 \frac1T\int_T^{2T}(\ee^{\alpha J}-1)\dd X
     =\alpha\int_0^\infty\ee^{\alpha t}d_T(t)\dd t.
\]
The displayed tail bound is an integrable majorant in the first
identity for every \(p>0\), and in the second for \(\alpha<4\).
Dominated convergence proves \eqref{eq:green-power-means}.
\end{proof}

In particular the actual Green potential tends to zero in every
finite \(L^p\) mean on dyadic intervals along this moving-time slice.
Narrow exceptional sets can survive every such mean statement.
There is no assertion of pointwise convergence or a supremum bound.
For an abstract limit check on a unit probability interval, take
\(Y_n=(\log n)\mathbf1_{[0,1/n]}/4\). Then every power mean
tends to zero and every exponential mean of exponent less than four
tends to one, but \(\mathbb E\ee^{4Y_n}=2-1/n\) and
\(\|Y_n\|_\infty\to\infty\). This is a measure-theoretic
example, not a theta field or a zero construction; it identifies the
logical limit of the mean statements alone.

\subsection{Exceptional centers for every fixed reference}
The mean law is compatible with unbounded exceptional residuals.
The reason is arithmetic coherence, which can be detected by fixed
rational-frequency tests without any quantitative equidistribution
rate.

\begin{lemma}[Twisted complete moments]\label{lem:twisted-moments}
For fixed coprime positive integers \(a,b\), let \(s=\log(a/b)\).
Then
\begin{equation}\label{eq:twisted-moment}
 \lim_{T\to\infty}\frac1T\int_T^{2T}\ee^{\ii Xs/2}M(X,1)\dd X
     =2\zeta(3/2)(ab)^{-3/4}\widehat K(s/2).
\end{equation}
\end{lemma}
\begin{proof}
In the auxiliary same-arm square the phase is
\(\nu(u)\log(n/m)+us/2\). Resonance requires \(an=bm\),
so \(n=bj,m=aj\), and its coefficient sum tends to
\((ab)^{-3/4}\zeta(3/2)\) for each of the two arm squares.
For nonresonant pairs
\(|\log(an/(bm))|\gg_{a,b}T^{-1/2}\).
The \(O(1/(T\log T))\) derivative perturbation preserves that
lower bound; the monotone first-derivative estimate costs
\(O_{a,b}(\sqrt T)\) per partial integral. The summed normalized
error is \(O_{a,b}(T^{-1/8})\). The linear twist leaves the
cross-arm second derivative unchanged. The correction and fixed-clock
payments are unchanged in absolute value. In the Fubini step the
twist contributes \(\ee^{-\ii ts/2}\) to the kernel at \(t=x-X\),
giving \(\widehat K(s/2)\) in place of its mass. The same edges
and tails finish the proof.
\end{proof}

\begin{theorem}[Fixed-reference exceptional residuals]
\label{thm:fixed-residual-unbounded}
The exact moment \(M(X,1)\) is unbounded at arbitrarily high centers.
For every fixed \(P\), so is \(R_P(X,1)\).
\end{theorem}
\begin{proof}
Choose a second fixed prime cutoff \(U\), independent of \(P\).
On its finite torus, with probability Haar measure \(dz\), put
\[
 E_U(z)=\prod_{p\le U}(1-p^{-3/4}z_p)^{-1},\quad
 C_U=\zeta(3/2)\prod_{p\le U}(1-p^{-3/2}),\quad
 \tau_p(t)=\ee^{\ii t\log p/2},
\]
\[
 F_U(z)=2C_U\int_\R K(t)|E_U(z\tau(t))|^2\dd t.
\]
This is continuous. For a character corresponding to
\(a/b=\prod p^{\ell_p}\), expansion of the finite-prime product
shows that \(\int z^\ell F_U(z)\dd z\) is the right side of
\eqref{eq:twisted-moment}. Along \(z_p(X)=\ee^{\ii X\log p/2}\),
every nonconstant character has zero limiting height mean, by unique
factorization. Thus each fixed trigonometric polynomial has the
corresponding weighted height mean against \(M\).

Use the nonnegative product Fej\'er polynomials
\[
 W_N(z)=\prod_{p\le U}\frac1N\left|\sum_{j=0}^{N-1}z_p^j\right|^2.
\]
Their Haar mean is one and they concentrate at \(z=1\): outside
\(|\arg z_p|<\eta\), the relevant factor is at most
\(1/[N\sin^2(\eta/2)]\). A finite union bound and continuity give
\(\int W_NF_U\dd z\to F_U(1)\).
If \(M\le H\) eventually, take the height limit first for fixed
\(U,N\) in the nonnegative inequality \(W_NM\le HW_N\), then
let \(N\to\infty\). This would imply \(F_U(1)\le H\) for every
fixed \(U\).

But for \(|t|\le1/\log U\) the prime cosines are at least \(7/8\),
and
\[
 |E_U(\tau(t))|^2\ge
        \exp\left(\frac34\sum_{p\le U}p^{-3/4}\right).
\]
The harmonic sum through \(\lfloor U\rfloor\) is at most
\(\prod_{p\le U}(1-1/p)^{-1}\); taking logarithms and paying the
summable higher powers gives
\(\sum_{p\le U}1/p\ge\log\log U-O(1)\).
Since \(p^{1/4}\ge4\) above a fixed cutoff,
\(\sum_{p\le U}p^{-3/4}\ge4\log\log U-O(1)\).
As \(C_U\ge1\) and \(K(t)\ge K(1)\) on this short interval,
\(F_U(1)\gg(\log U)^2\), a contradiction.

Finally \(B_P:=\sup_x|E_P(x)|<\infty\) for fixed \(P\). The
weighted triangle inequality gives
\begin{equation}\label{eq:residual-unbounded-transfer}
 R_P(X,1)\ge B_P^{-2}
        \bigl(\sqrt{M(X,1)}-2B_P\sqrt{q_X}\bigr)_+^2,
\end{equation}
which proves the residual assertion.
\end{proof}

This also obstructs the specified separate residual/Jensen payment
even with a regulator chosen at each center. Before any phase estimate,
the regularized reference log is at least \(q_X\log\delta\); hence
its sum with the residual payment is at least
\[
 \frac{q_X}{2}\log(\delta^2+R_P/q_X)
      \ge\frac{q_X}{2}\log(R_P/q_X)
\]
when \(R_P>0\). Restoring the bounded fixed-reference prime
difference cannot prevent divergence at those centers. This is a
statement about that allowance, not about the actual Green value.

\subsection{A growing reference with paid comparison costs}
Now choose \(P=P(X)=\lfloor L_X^2\rfloor\), retaining this one
value throughout each window. Equations~\eqref{eq:euler-phase-derivatives}
give \(D_1(P),D_2(P)=O(L_X^{3/4}(\log L_X)^2)=o(L_X)\).
Thus the phase gates hold and \(D\le L_X\) eventually, uniformly
in \(c\). The joint phase costs at most \(18/L_X\), for every
positive regulator. More generally the same conclusion holds for
\(2\le P\le L_X^\gamma\), with any fixed \(0<\gamma<8/3\).
At the selected cutoff the prime tail is \(O(L_X^{-3})\) and the
reference-log exterior is \(O(L_X^{3/4}X^{-4})\).

The full correction and provider error remain small after normalization:
\begin{equation}\label{eq:adaptive-provider-error}
 \sup_{W_X}\left|\frac{V_\lambda-S_\lambda}{E_P}\right|
 \le CX^{-1/4}\exp\{(8/3)L_X^{3/4}\}=X^{-1/4+o(1)}.
\end{equation}
Indeed \(|E_P|^{-1}\le\prod_{p\le P}(1+p^{-\sigma_c})
\le\exp((8/3)P^{3/8})\). The complete adaptive inequality is
\begin{align}\label{eq:adaptive-allowance}
 J_\lambda(X)\le{}&P_{c,\lambda}^{(\le P)}-P_{c,\lambda}
 +q_X\log b_\delta
 +\frac{q_X}{2}\log\left(1+\frac{R_{P(X)}(X,c)}{q_X\delta^2}\right)
 \nonumber\\
 &+T_X+\frac{C_{\rm out}+18}{L_X}
       +\frac{16}{3\pi\kappa_E}P^{3/8}X^{-4}.
\end{align}
Here \(q_X\log b_\delta\le\delta/4\). All comparison costs can
vanish with a suitable regulator; the exact arithmetic residual in
\eqref{eq:exact-relative-residual} is still unpaid. In particular
\eqref{eq:adaptive-provider-error} estimates \(V-S_\lambda\),
not \(V-S_P^{\rm ref}\). Neither the fixed-\(P\) mean law nor
the fixed-\(P\) unboundedness theorem permits substituting this
growing cutoff into its quantifiers.

\subsection{The specified signed tangent still fails for fixed \(P\)}
Retaining a signed cross term is not by itself enough. Fix finite
\(P\), take \(c=1\), and set
\[
 w=V/E_P,\quad a=S_P^{\rm ref}/E_P,\quad h=|a|^2+\delta^2,
 \qquad G_\delta=\tfrac12\log h-\tfrac12+\frac{|w|^2}{2h}.
\]
The inequality \(\log t\le t-1\) gives \(\log|w|\le G_\delta\),
with the exact cross term retained in
\(|w|^2=|a|^2+2\Re(\bar a(w-a))+|w-a|^2\).
Let \(H=4+\delta^2\) and \(Z=\int_{W_X}|w|^2\,d\omega_X\).
Since \(|a|\le2\), \(h\le H\) and \(h/H\ge|a|^2/4\).
Integrating, and then using \(t-1\ge\log t\) at \(t=Z/(q_XH)\),
gives the regulator-independent lower bound
\begin{align}\label{eq:tangent-obstruction}
 \int_{W_X}G_\delta\,d\omega_X
 &\ge\frac{q_X}{2}\log H-\frac{q_X}{2}+\frac Z{2H}
         +\int_{W_X}\log|a|\,d\omega_X-q_X\log2\nonumber\\
 &\ge\frac{q_X}{2}\log(Z/q_X)
         +\int_{W_X}\log|a|\,d\omega_X-q_X\log2.
\end{align}
The reference logarithm is \(O_P(1/L_X)\) by
Lemma~\ref{lem:joint-phase}, including its zeros. For fixed \(P\),
\(Z\ge M(X,1)/B_P^2\). Theorem~\ref{thm:fixed-residual-unbounded}
therefore makes the integrated tangent unbounded at arbitrarily high
centers, independently of \(\delta\). Restoring the bounded reference
prime difference and the nonnegative tail allowances preserves this
failure. The result concerns this fixed-reference tangent, with a
regulator constant on each window. It does not cover a growing
reference, spatially varying tangent heights, or all nonlinear signed
estimates, and gives no lower bound for \(J\).

\section{Short-window residual and Green means}
\label{sec:b10-short-windows}

We use the exact residual \eqref{eq:exact-relative-residual} at \(c=1\),
with physical time \(1/L_X\) held fixed throughout \(W_X\). The prime
cutoff \(P\) is constant on the outer interval \([U,U+\ell]\).
The following theorem combines the growing-cutoff diagonal, the exact
convolution weight and the short-window arm bounds
\cite{GNPRSWeightedEulerLedger,GNPRSLocalizedArm,GNPRSGrowingExponentPair,GNPRSSignedMellin}.
Write
\[
 \mathscr R(U,\ell;P)=\frac1\ell\int_U^{U+\ell}R_P(X,1)\dd X,
 \qquad
 \mu_P=\sum_{\substack{k>1\\(k,D_P)=1}}k^{-3/2}\le2P^{-1/2},
 \qquad D_P=\prod_{p\le P}p .
\]

\begin{theorem}[Short-window mean of the growing-cutoff residual]
\label{thm:b10-averaged}
Put \(\beta=17/81\). For each fixed \(0<\rho\le47/162\) and all
sufficiently large real \(U\), uniformly for integers
\(2\le P\le\lfloor L_U^2\rfloor\) and
\(U^{\beta+\rho}\le\ell\le U\),
\begin{equation}\label{eq:b10-main}
 \left|\mathscr R(U,\ell;P)-\mu_P\right|
 \le C\left(\frac{\sqrt{\mu_P}}{L_U}+L_U^{-2}\right)
     +O_\rho\!\left(U^{-\min(7/6480,\,9\rho/28)}\right).
\end{equation}
Here \(C\) depends only on the fixed inputs. The large-height onset
is finite and unevaluated.
\end{theorem}

\subsection{Exact Mellin projection of the upper part of the arm}

For the short-window proof put \(L=L_U\), \(\mathcal L=1+\log U\),
and choose
\[
 B=N_U=\left\lfloor\sqrt{U/(4\pi)}\right\rfloor,\qquad
 Z=\lfloor U^{1/3}\rfloor,\qquad Y=\lfloor U^{1/10}\rfloor,\qquad
 M=U^{1/10}.
\]
Take a fixed nonnegative smooth \(\psi_0\) of integral one supported
on \([-1/2,1/2]\) and put
\(\chi=\mathbf1_{[-3/2,5/2]}*\psi_0\).
Then \(\chi=1\) on \([-1,2]\), has support in \([-2,3]\) and
has integral four. For \(0\le a<b\le B\) define
\begin{align*}
 w_x=w_x(U)&=x^{-3/4}\exp((\log x)^2/(4L)),\qquad
 \tau(u)=u/2+\pi/(8L),\\
 D_{a,b}(u)&=\sum_{a<n\le b}w_n\,\ee^{\ii\tau(u)\log n},\\
 \mathcal Q_{a,b}&=\ell^{-1}\int_\R
       \chi((u-U)/\ell)|D_{a,b}(u)|^2\dd u,\\
 \mathcal S_{a,b}&=\sum_{a<m<n\le b}w_mw_n
       \ee^{\ii\tau(U)\log(n/m)}
       \widehat\chi\bigl((\ell/2)\log(n/m)\bigr),
 \qquad
 \widehat\chi(\xi)=\int_\R\chi(v)\ee^{\ii v\xi}\dd v.
\end{align*}
Expanding the square and pairing the two orderings gives
\begin{equation}\label{eq:b10-signed}
 \mathcal Q_{a,b}=4\sum_{a<n\le b}w_n^2+2\Re\mathcal S_{a,b}.
\end{equation}

We first show, for \(1\le\ell\le\sqrt U\), that
\begin{equation}\label{eq:b10-high}
 \mathcal Q_{Z,B}+|\Re\mathcal S_{Z,B}|
 \ll\mathcal L^3\left(U^{-1/9}+\frac{U^{37/180}}{\ell}\right).
\end{equation}
Use the cumulative zeta mean-square error of Bourgain and Watt
\cite[Theorem 4]{BourgainWattPerturbedCones2015v1}, published in
\cite{BourgainWattPerturbedCones2018}. For \(T\ge1\) they give
\[
 \int_0^T|\zeta(1/2+\ii t)|^2\dd t
 =T\bigl(\log(T/(2\pi))+2\gamma-1\bigr)
      +O_\epsilon(T^{1515/4816+\epsilon}).
\]
Fix \(\theta=19/60\) and
\(\epsilon=151/144480<\theta-1515/4816\).
For a subinterval of \([1,R]\) subtract the two cumulative
identities, integrating the derivative
\(\log(t/(2\pi))+2\gamma\) of the main term.
The two endpoint errors cost \(O(R^\theta)\). Reflect negative heights
and use boundedness on \([-1,1]\). For every interval \(I\subset[-R,R]\)
of length at most \(c_*\ell\), \(R\ge2\), \(\ell\ge1\), this gives
\begin{equation}\label{eq:b10-local-zeta}
 \int_I|\zeta(1/2+\ii t)|^2\dd t
 \ll_{c_*} \ell\log(2+R)+R^\theta.
\end{equation}
In particular the cumulative error remains in the translated bound.

Let \(q\) be a fixed smooth increasing cutoff with \(q=0\) on
\((-\infty,0]\) and \(q=1\) on \([1/2,\infty)\), and set
\[
 a(x)=\sqrt x\,w_x,\qquad
 b(x)=a(x)q(x-Z-1/4)q(B+3/4-x),\qquad
 \widetilde b(z)=\int_0^\infty b(x)x^z\frac{\dd x}{x}.
\]
The transitions lie in \((Z+1/4,Z+3/4)\) and
\((B+1/4,B+3/4)\). At every integer \(n\),
\(b(n)=\sqrt n\,w_n\mathbf1_{Z<n\le B}\).
Put \(A_*=U^{-1/18}\) and \(g(y)=b(\ee^y)\). On the support,
\(\alpha=\log x/\log U\) ranges from \(1/3+O(U^{-1/3}/\log U)\)
to \(1/2+O(1/\log U)\).
Since \(\log a(U^\alpha)=-\alpha(1-\alpha)\log U/4+O(1)\),
we have \(\sup a\ll A_*\). The logarithmic derivative of \(a(\ee^y)\)
with respect to \(y\) is \(-1/4+y/(2L)\) and changes sign at most once.
The interior total variation is \(O(A_*)\) and the interior
second derivative has integral \(O(A_*\mathcal L)\).
A transition at \(x\asymp N\), \(N=Z,B\), has \(y\)-width
\(O(1/N)\) and first and second \(y\)-derivatives bounded by
\(O(N)\) and \(O(N^2)\). The product rule gives
\[
 \|g\|_1\ll A_*\mathcal L,\qquad
 \|g'\|_1\ll A_*,\qquad
 \|g''\|_1\ll A_*B.
\]
For \(k(v)=\widetilde b(\ii v)/(2\pi)\) integration by parts yields
\[
 |k(v)|\ll A_*\min\{\mathcal L,|v|^{-1},B|v|^{-2}\}
 \quad(v\ne0).
\]
Split the integrals at \(1/\mathcal L,B,U\), with \(B\le\sqrt U\).
The resulting bounds are
\begin{align*}
 \int_\R |k(v)|\dd v&\ll A_*\mathcal L,\\
 \int_\R |k(v)|\log(2+U+|v|)\dd v&\ll A_*\mathcal L^2,\\
 \int_\R |k(v)|(U+|v|)^\theta\dd v&\ll A_*U^\theta\mathcal L.
\end{align*}
For the last two integrals, the \(|v|>U\) pieces cost
\(O(A_*B(1+\log U)/U)\) and \(O(A_*B U^{\theta-1}/(1-\theta))\),
respectively.

With \(s=1/2-\ii\tau(u)\), Mellin inversion first on \(\Re z=1\)
permits the absolutely convergent Dirichlet series for
\(\zeta(s+z)\). Shift to \(\Re z=0\). For each fixed \(U,u\), smooth
compact support gives arbitrary vertical decay of \(\widetilde b\) on
the contour strip, and Euler summation gives polynomial growth of
\(\zeta\) away from its pole. The crossed simple pole has residue
\(\widetilde b(1-s)\). Thus the exact identity is
\begin{equation}\label{eq:b10-mellin}
 D_{Z,B}(u)=\sum_n b(n)n^{-s}
   =\widetilde b(1-s)+\int_\R k(v)\zeta(s+\ii v)\dd v .
\end{equation}
The weight \(\ee^{y/2}a(\ee^y)\) increases on the support and is bounded
by \(C A_*\sqrt B\). After multiplying by the cutoffs its variation is
still \(O(A_*\sqrt B)\), so one integration by parts pays the residue
by \(C A_*\sqrt B/|\tau(u)|\ll A_*U^{-3/4}\) on the \(\chi\)-window.

Apply Cauchy--Schwarz to the last integral in \eqref{eq:b10-mellin}
with measure \(|k(v)|\dd v\), then use Tonelli's theorem with
\(\chi((u-U)/\ell)\dd u/\ell\). For each \(v\), the substitution
\(t=v-\tau(u)\) has absolute Jacobian \(|\dd u/\dd t|=2\).
It sends the support into an interval of length \(5\ell/2\)
in \([-C(U+|v|),C(U+|v|)]\).
Equation~\eqref{eq:b10-local-zeta} and the weighted kernel bounds give
\[
 \mathcal Q_{Z,B}
 \ll A_*^2\left(U^{-3/2}+\mathcal L^3+
                      \frac{U^\theta\mathcal L^2}{\ell}\right).
\]
The diagonal satisfies \(\sum_{Z<n\le B}w_n^2\ll A_*^2\mathcal L\).
Since \(\theta-1/9=37/180\), \eqref{eq:b10-signed} proves
\eqref{eq:b10-high}. Its signed estimate concerns the real quadratic
sum over all pairs, including all lattice resonances.

\subsection{Lower modes and the weighted arm norms}
For the shorter windows \(\ell\le\sqrt U\), use the auxiliary arm
\(A(u)\) at time \(1/L_u\) from \eqref{eq:real-two-arms}, with
\(A_Y(u)\) its restriction to \(n\le Y\). Its weights satisfy
\(w_n(u)\le n^{-5/8}\) for \(n\le N_u\). Define
\begin{align*}
 r_*&=\frac{32}{\pi}\log(2U),&
 \mathcal I&=[U-r_*,U+\ell+r_*],\\
 W(u)&=\int_U^{U+\ell}K(u-X)\mathbf1_{u\in W_X}\dd X,&
 d\sigma(u)&=\frac{2W(u)}{\ell}\dd u .
\end{align*}
The B07 weight \cite{GNPRSLocalizedArm} has \(0\le W\le1/2\),
support in \(\mathcal I\), and
\begin{equation}\label{eq:b10-weight}
  1-\frac4\pi U^{-4}\le m_\sigma:=\int d\sigma\le1,\qquad
  \operatorname{TV}(W)\le1+\frac{\ell}{2U^4}.
\end{equation}
The mass bound integrates the kernel tail
\(\int_{|v|>R_X}K(v)\dd v\le (2/\pi)X^{-4}\).
The untruncated convolution has derivative
\(K(u-U)-K(u-U-\ell)\) of \(L^1\) norm at most one.
Each removed tail has total variation \(4K(R_X)\le X^{-4}/2\),
including the two jumps. Integrating their variations proves the
second bound. For \(r_*\le\ell\), \(\mathcal I\subset[U-\ell,U+2\ell]\),
where the smooth weight \(\chi((u-U)/\ell)\) equals one. Thus it
majorizes \(2W(u)\) in upper bounds for nonnegative norm squares.

On the full smooth support \([U-2\ell,U+3\ell]\), freeze the
weights and the small \(\pi/(8L_u)\) phase correction at \(U\).
For common active indices their summed drift is
\(O(\ell U^{-13/16})\), using \(O(U^{-1})\) bounds for the
logarithmic weight derivative and for the derivative of the
small phase after multiplication by \(\log n\). At the upper
cutoff there are \(O(1+\ell/\sqrt U)\) changed indices of size
\(\asymp\sqrt U\), each of weight \(O(U^{-5/16})\).
With \(Y\le N_{U-2\sqrt U}\), this gives uniformly
\[
 |(A-A_Y)-D_{Y,B}|+|A_Y-D_{0,Y}|\ll U^{-5/16}.
\]
The cutoff jumps therefore cost \(O(U^{-5/8})\) in the norm-square
comparison, which uses \(|f+g|^2\le2|f|^2+2|g|^2\).

In the Fourier expansion of \(\mathcal Q_{Y,Z}\), partition the
pairs \(Y<m<n\le Z\) into \(n\ge2m\) and \(m<n<2m\).
For the latter write \(n=m+h\) and partition \(m\) into
dyadic blocks \(N\le m<2N\). For fixed \(h\) the cutoff
conditions restrict \(m\) to an interval in this block. Put
\(t=(U/2+\pi/(8L))/(2\pi)\asymp U\). The conjugate exponential
phase, in the convention \(\exp(2\pi\ii f)\), is
\[
 f(x)=-t\log(1+h/x)=-t\int_0^h(x+s)^{-1}\dd s .
\]
For \(\varphi(x)=-th/x\), every derivative order \(j\ge1\) satisfies
\[
 \frac{|f^{(j)}(x)-\varphi^{(j)}(x)|}{|\varphi^{(j)}(x)|}
 \le \frac{(j+1)h}{2x}.
\]
This follows by averaging
\(1-(1+s/x)^{-j-1}\le(j+1)s/x\). Take the derivative
parameter \(2\) and scale \(th\) in the exponent-pair class
of Montgomery and Vaughan \cite[Definitions 16.1--16.2 and
Theorems 16.23, 16.25]{MontgomeryVaughanExponentialSums2026}.
The operations \(B,A,A,B\) from \((0,1)\) give \((2/7,4/7)\).
Fix the derivative order and tolerance for this pair, then
choose a fixed \(\delta_0>0\) small enough that \(h\le\delta_0N\)
ensures the required relative derivative bounds. The unweighted
sum over each subinterval is
\[
 \ll (th/N^2)^{2/7}N^{4/7}+N^2/(th)
 \ll (Uh)^{2/7}+N^2/(Uh).
\]
Both terms are used.

The exact weight is decreasing in its index, with
\(|w_x'(U)|\ll w_x(U)/x\). For
\[
 a_h(x)=w_x(U)w_{x+h}(U)
          \widehat\chi\bigl((\ell/2)\log(1+h/x)\bigr)
\]
the supremum plus total variation on each allowed interval is
\(\ll_D w_N(U)^2(1+\ell h/N)^{-D}\) for fixed \(D\).
Indeed the weight derivatives cost \(O(w_N(U)^2/N)\),
and one extra Fourier decay power pays the integral of the
argument derivative \(O(\ell h/N^2)\).
Abel summation and
\(\sum_{h\ge1}h^{2/7}(1+\ell h/N)^{-D}\ll (N/\ell)^{9/7}\)
for \(D\ge4\) give the leading cost per block
\[
 U^{2/7}N^{-3/14}\exp\bigl((\log N)^2/(2L)\bigr)\ell^{-9/7}.
\]
Writing \(\alpha=(\log N)/(\log U)\), the exponent of \(U\)
in its coefficient is \(2/7-3\alpha/14+\alpha^2/2+O(1/\log U)\).
The quadratic is convex. Its endpoint values on \([1/10,1/3]\)
are \(377/1400\) and \(17/63\), with difference \(1/1800\).
Floor and dyadic factors cost \(O(1)\) in the logarithm of the
coefficient. The block count is \(O(\mathcal L)\).
For the \(N^2/(Uh)\) term, the coarser \(N^{-5/4}\) amplitude
and \(\sum_{h\ge1} h^{-1}(1+\ell h/N)^{-D}\ll\log(2+N/\ell)\)
give \(O(U^{-1}Z^{3/4}\mathcal L)\) after dyadic summation.
For \(h>\delta_0 N\) and for noncomparable pairs, the log ratio
is bounded below by a fixed positive constant. The Fourier
factor is \(O(\ell^{-4})\) and the squared absolute sum of the
weights is \(O(U^{3/8})\). The diagonal is at most \(16Y^{-1/4}\).
Altogether,
\begin{equation}\label{eq:b10-low}
 \mathcal Q_{Y,Z}\ll Y^{-1/4}+\mathcal L U^{17/63}\ell^{-9/7}
                       +U^{3/8}\ell^{-4}+U^{-5/8}\mathcal L .
\end{equation}

The exact split \(D_{Y,B}=D_{Y,Z}+D_{Z,B}\) pays all cross
terms through the square triangle bound with factor two.
In the frozen head the diagonal is bounded, and the distinct
log spacings are at least \(1/Y\). Second-power Fourier decay
and \((\sum_{n\le Y}w_n(U))^2\ll Y^{3/4}\) give the
off-diagonal cost \(O(Y^{11/4}/\ell^2)\), which decays on the
short-window range in Theorem~\ref{thm:b10-averaged}.
Together with \eqref{eq:b10-high}, \eqref{eq:b10-low}, and the
freezing comparison this gives, for a sufficiently large fixed \(C\),
\begin{align}
 \mathcal E&=\mathcal L U^{17/63}\ell^{-9/7}
      +\mathcal L^3(U^{-1/9}+U^{37/180}/\ell)
      +U^{3/8}\ell^{-4}+U^{-5/8}\mathcal L,\nonumber\\
 \int |A-A_Y|^2d\sigma&\le A_1:=C(Y^{-1/4}+\mathcal E),\nonumber\\
 \int |A|^2d\sigma,\quad\int |A_Y|^2d\sigma
   &\le A_0:=C(1+\mathcal E).
 \label{eq:b10-norms}
\end{align}

\subsection{The residual estimate}
At auxiliary time \(1/L_u\) write \(Q=E_P^{-1}\), \(D_P=\prod_{p\le P}p\),
and for an arm \(a\) put
\[
 \mathcal Z(a)=aQ-1+\ee^{\ii\phi}(Q/\overline Q)\overline{aQ-1},
 \qquad
 \begin{aligned}
 \phi(u)&=u(L_u-1)/2+\pi/8+7\pi/4,\\
 \nu(u)&=u/2+\pi/(8L_u).
 \end{aligned}
\]
The weighted residual estimate uses the B06 coefficient and clock
bounds \cite{GNPRSWeightedEulerLedger} and B07's exact boundary
calculation. First compare the physical square at time \(1/L_X\),
fixed throughout \(W_X\), to \(|\mathcal Z(A)(u)|^2\).
At fixed \(u\in W_X\) the two auxiliary arms have size \(O(U^{7/32})\)
and time derivative \(O(L^2U^{7/32})\) on the segment from \(1/L_X\)
to \(1/L_u\), by the weights and carrier phase in
\eqref{eq:real-two-arms}. On this segment keep the Euler exponent
\(3/4\). With
\[
 H=\prod_{p\le P}(1-p^{-3/4})^{-1},\qquad
 T_1(P)=\sum_{p\le P}\frac{p^{-3/4}\log p}{1-p^{-3/4}},
\]
the reciprocal product has modulus at most \(H\) and time derivative
\(O(HT_1(P))\). Its ratio with its conjugate has unit modulus and
time derivative \(O(T_1(P))\). Since \(T_1(P)\ll P^{1/4}\log P\le L^2\)
eventually, the relative square has time derivative
\(O(H^2L^2 U^{7/16})\). The displacement
\(O(r_* /(UL^2))\) costs \(O(H^2r_*U^{-9/16})\).
The exact field error in \eqref{eq:real-provider-error}, after
multiplication by \(Q\) and comparison of squares, adds
\(O(H^2U^{-1/32})\). This calculation differentiates only the
arithmetic approximation at fixed \(u\). Fubini now gives
\begin{equation}\label{eq:b10-fubini}
 \left|\mathscr R-\frac12\int|\mathcal Z(A)|^2\,d\sigma\right|
 \le C_{\rm clk}H^2(U^{-1/32}+r_*U^{-9/16}).
\end{equation}
The factor \(1/2\) comes from the definition of \(d\sigma\).

Let \(Q_M,E_M\) denote the expansions of \(Q,E_P\) truncated
by integer index at \(M\). Their indices are squarefree divisors
of \(D_P\) for \(Q_M\) and integers with all prime factors at most
\(P\) for \(E_M\). Their absolute coefficient sums are at most
\(H\), as for the complete expansions. Put \(B_M=A_YQ_M-1\) and
\(m_M=\overline{Q_M}E_M\). With
\[
 T_M=M^{-1/8}\prod_{p\le P}(1-p^{-5/8})^{-1},
\]
both discarded tails have supremum at most \(T_M\). This follows
from \(n^{-3/4}\le M^{-1/8}n^{-5/8}\) for \(n>M\).
Also \(\|m_M-\overline Q/Q\|_\infty\le2HT_M\), and
the coefficient absolute sum of \(m_M\) is at most \(H^2\).

For \(L_0=L_{U-r_*}\), \(K_0=\lfloor M\rfloor Y\) set
\begin{align*}
 a_0&=C_2/(4L_0)+C_0Y^{-1/8}+C_0M^{-1/8},
 &D_0&=2\sqrt{\mu_P}\,a_0+a_0^2,\\
 s_0&=150000 K_0(1+\log K_0)/\ell,
 &b_0&=(\sqrt{\mu_P}+a_0)^2+s_0,\\
 c_0&=8000 H^4Y^{3/4}/(\ell L_0).
\end{align*}
The coefficient vector of \(B_M\) has distance at most \(a_0\)
in \(\ell^2\) from \(k^{-3/4}\mathbf1_{k>1,(k,D_P)=1}\),
by Möbius cancellation. More explicitly, the individual difference
is bounded by \(k^{-5/8}d(k)\) times
\((\log k)^2/(4L_0)+Y^{-1/8}+M^{-1/8}\), where \(d(k)\) counts
divisors. The three terms bound the Gaussian weight correction,
the head truncation and the reciprocal truncation, respectively.
Here \(C_0^2=\sum_k d(k)^2 k^{-5/4}\) and
\(C_2^2=\sum_k d(k)^2(\log k)^4k^{-5/4}\), which converge by
\(d(k)^2\le d_4(k)\) and the absolute Dirichlet series for
\(\zeta(s)^4\), \(s>1\). In particular \(C_0\le25\).
The diagonal error at each \(u\) is at most \(D_0\).

For same-arm phases \(\nu\log(k/l)\), \(k\ne l\), the partial
integrals are at most \(12/|\log(k/l)|\) by monotonicity
and \(1/4\le\nu'\le1/2\). The coefficients have supremum at most
\(g_k=k^{-5/8}d(k)\) and variation at most \(2g_k\) by summing
the decreasing weights over divisors, including any activation
jumps. For the product of two coefficients multiplied by \(W\),
the supremum plus variation is at most \(4g_kg_l\) by
\eqref{eq:b10-weight}. Using
\(|\log(k/l)|^{-1}\le K_0/|k-l|\), \(2g_kg_l\le g_k^2+g_l^2\)
and the harmonic row sum gives the weighted off-diagonal cost
\(s_0\), with the factor \(2/\ell\) from \(d\sigma\). Hence
\[
 \left|\int|B_M|^2d\sigma-\mu_P\right|
 \le D_0+s_0+\mu_P(1-m_\sigma),\qquad
 \int|B_M|^2d\sigma\le b_0.
\]
The reflected phases are \(-\phi+\nu\omega\) with
\(\omega=\log(kle/d)\), \(k,l\le K_0\), \(e,d\le M\).
Their absolute log frequencies are bounded by
\(2\log Y+3\log M\le\tfrac12\log U\).
On \(\mathcal I\), eventually this is at most \(7L_0/8\),
and \(\nu''|\omega|\ll U^{-2}/L\le1/(12U)\) while
\(\phi''=1/(2u)\ge1/(6U)\). The first derivative is thus
monotone and at most \(-L_0/16\), giving partial integrals
bounded by \(48/L_0\).
The pre-convolution absolute coefficient envelope of \(B_M\)
sums to at most \((8/3)HY^{3/8}\), since the coefficient at
1 cancels. Combining the product variation estimate for these
terms with the absolute coefficient sum \(H^2\) of \(m_M\)
bounds \(\big|\int 2\Re(\ee^{-\ii\phi}B_M^2m_M)d\sigma\big|\)
by \(c_0\).

The exact square for \(A_Y\) is
\(2|A_YQ-1|^2+2\Re(\ee^{-\ii\phi}(A_YQ-1)^2\overline Q/Q)\).
To replace the complete arm by \(A_Y\), use
\(\|\mathcal Z(A)-\mathcal Z(A_Y)\|_{L^2(d\sigma)}
 \le2H\sqrt{A_1}\) and
\(\|\mathcal Z(A)\|_{L^2(d\sigma)},\|\mathcal Z(A_Y)\|_{L^2(d\sigma)}
 \le4H\sqrt{A_0}\). For the multiplier replacement use
\(\|(A_YQ-1)-B_M\|_{L^2(d\sigma)}\le T_M\sqrt{A_0}\),
the bound \(b_0\), and \(|\overline Q/Q|=1\).
Combining the square differences, the diagonal and reflected
bounds, and \eqref{eq:b10-fubini}, yields
\begin{align}
 |\mathscr R-\mu_P|\le{}&
 D_0+s_0+\tfrac12c_0+8H^2\sqrt{A_0A_1}
    +4T_M\sqrt{A_0b_0}+2T_M^2A_0+2HT_M b_0\nonumber\\
 &+C_{\rm clk}H^2(U^{-1/32}+r_*U^{-9/16})
    +\frac4\pi\mu_P U^{-4}.
 \label{eq:b10-ledger}
\end{align}

To conclude Theorem~\ref{thm:b10-averaged}, take
\(\beta=17/81\) and \(\ell\ge U^{\beta+\rho}\).
In \eqref{eq:b10-norms} the low term has power \(-9\rho/7\),
and the translated high term has power \(-7/1620-\rho\),
before logarithms. The diagonal tail has power \(-1/40\),
the remaining large-shift cost has power \(-301/648-4\rho\),
and freezing and the phase remainder cost \(O(U^{-5/8}\mathcal L)\).
The head's auxiliary off-diagonal has power \(-469/3240-2\rho\).
In \eqref{eq:b10-ledger}, \(s_0\ll\mathcal L U^{-4/405-\rho}\)
and \(c_0\ll H^4U^{-437/3240-\rho}/L_0\).
The two coefficient tails in \(a_0\) and the power in \(T_M\)
are \(O(U^{-1/80})\). In particular \(A_0,b_0=O(1)\) eventually.
Integral comparison gives \(H\le\exp(CP^{1/4})\) and
\(T_M\le U^{-1/80}\exp(CP^{3/8})\).
The remaining polynomial errors, including the field, clock and
mass terms, fit
\[
 |\mathscr R-\mu_P|
 \le C(\sqrt{\mu_P}/L+L^{-2})
 +C\exp\bigl(C(P^{1/4}+P^{3/8})\bigr)\mathcal L^{C'}
       (U^{-7/3240}+U^{-9\rho/14}).
\]
The positive exponents here precede absorption of the product and
logarithm factors. Uniformly for \(P\le L^2\), the logarithm of
those factors is \(O(L^{3/4}+\log\mathcal L)=o(\log U)\). It
can be bounded by \(\tfrac12\min(7/3240,9\rho/14)\log U\)
after a finite onset. This proves the stated remainder. All
preceding gates hold at finite heights for fixed \(\rho>0\):
\(r_*\le\ell\), \(\mathcal I\subset[U/2,3U]\), \(L_0\ge16\),
\(2\le Y<Z<B\), \(Y\le N_{U-2\sqrt U}\), and the field,
frequency and product bounds above. The common onset is
unevaluated. For \(\sqrt U\le\ell\le U\), the inherited B08
translated residual theorem \cite[Eq.~(B08.1)]{GNPRSGrowingExponentPair}
with its fixed margin \(19/72\) starts at \(\ell=\sqrt U\)
and gives the same leading logarithmic terms with remainder
\(O(U^{-1/144})\). It uses the same block-constant \(P\) and
physical clock.

\subsection{Green consequences at the block cutoff}
\begin{corollary}\label{cor:b10-green}
On the interval range of Theorem~\ref{thm:b10-averaged}, take
\(P_U=\lfloor L_U^2\rfloor\), constant over \([U,U+\ell]\).
Then
\begin{align*}
 \mathscr R(U,\ell;P_U)&=\mu_{P_U}+O_\rho(L_U^{-3/2})
                                    =O_\rho(L_U^{-1}),\\
 \frac1\ell\int_U^{U+\ell}J_{1/L_X}(X)\dd X
   &=O_\rho(L_U^{-1/3}),\\
 \frac{|\{X\in[U,U+\ell]:J_{1/L_X}(X)\ge h\}|}{\ell}
   &=O_{h,\rho}(L_U^{-1})\qquad(h>0\ \text{fixed}).
\end{align*}
\end{corollary}
\begin{proof}
The first assertion uses \(\mu_P\le2P^{-1/2}\).
For the Green transfer, \(L_X=L_U+O(1)\) and \(P_U\le L_X^2\)
throughout the interval. The phase gates in
Lemma~\ref{lem:joint-phase} thus hold uniformly by
\eqref{eq:euler-phase-derivatives}. In \eqref{eq:euler-master},
the omitted primes cost \(O(P_U^{-3/2})=O(L_U^{-3})\).
The reference exterior costs \(O(L_U^{3/4}U^{-4})\) and the
complete-field logarithmic exterior \(O(U^{-3}\log^2U)\).
Since \(q_X\log b_\delta\le\delta/4\) and \(\log(1+x)\le x\),
\[
 0\le J_{1/L_X}(X)\le
     \delta/4+\frac{R_{P_U}(X,1)}{2\delta^2}+C/L_U .
\]
Taking \(\delta=L_U^{-1/3}\) and averaging proves the second
assertion. For fixed \(h>0\) take a fixed \(\delta>0\) with
\(\delta/4\le h/4\), and let \(U\) be large enough for
\(C/L_U\le h/4\). Then \(J_{1/L_X}(X)\ge h\) implies
\(R_{P_U}(X,1)\ge h\delta^2\), so Markov's inequality proves
the last assertion.
\end{proof}

The bound on total Green mass is \(O_\rho(\ell L_U^{-1/3})\),
which grows on the admissible windows. A localized defect may fit inside that
budget. A bound at every center would require further control
of \(J_\lambda(X)\) along the moving-time curve
\(\lambda=1/L_X\).

\section{A square-root obstruction to uniform resolvent stability}\label{sec:conditioning}

The difficulty in \eqref{eq:ratio-error} is not removed merely by
imposing common source and zero-count bounds. The following family
makes the loss quantitative while fixing the translation scale.

Let \(H=H_0\), \(\alpha=(\log2)/2\), and
\[
 S_\alpha(u)=\frac{\Phi(u+\alpha)+\Phi(u-\alpha)}2,\qquad
 \Psi_\eta(u)=\frac{(1+\eta)\Phi(u)+(1-\eta)S_\alpha(u)}2,
\]
where \(0\le\eta\le\eta_0:=\tanh^2(\alpha/4)\).
The transforms are
\begin{equation}\label{eq:control}
 f_\eta(z)=\frac12\int_\R\Psi_\eta(u)\ee^{\ii zu}\dd u
 =H(z)M_\eta(z),\qquad
 M_\eta(z)=\cos^2(\alpha z/2)+\eta\sin^2(\alpha z/2).
\end{equation}
Every source is strictly positive, smooth and even, and
\(f_\eta(0)=H(0)\). The parameter \(\eta\) varies the source; it is
not Newman time. Each source has its own genuine heat orbit, but
\eqref{eq:control} is the time-zero factorization.

\begin{theorem}[Uniform square-root conditioning]\label{thm:conditioning}
The family \(\Psi_\eta\) has common bounds in every fixed exponentially
weighted source seminorm and common full zero-count constants. Its
relative source distance from \(\Psi_0\) is at most \(\eta\), and
each fixed safe-line logarithmic-resolvent jet changes by \(O(\eta)\),
uniformly along that line. Nevertheless,
\begin{equation}\label{eq:delta-scale}
 c\sqrt\eta\le\calJ[f_\eta](X)-\calJ[f_0](X)
                  \le C_\alpha\sqrt\eta
\end{equation}
for every \(X\in\R\), with constants independent of \(\eta,X\).
Thus no nondecreasing modulus \(o(\sqrt t)\), in particular no
uniform Lipschitz modulus, controls the Green difference from these
observations on a class containing this family.
\end{theorem}

\begin{proof}
The exact source difference is
\[
 \Psi_\eta-\Psi_0=\frac\eta2(\Phi-S_\alpha),\qquad
 |\Psi_\eta-\Psi_0|\le\eta\Psi_0.
\]
For fixed \(L,b\ge0\) and integers \(j,k\ge0\), its differentiated
weighted \(L^1\) norm is bounded by
\[
 \frac\eta2\int_\R |u|^j\ee^{Lu^2+b|u|}
            (|\Phi^{(k)}|+|S_\alpha^{(k)}|)\dd u.
\]
This integral is finite by Lemma~\ref{lem:source} and fixed translation.
The bound is independent of horizontal observation height.

Write
\[
 \ell=\frac{2\pi}{\alpha},\qquad
 x_k=\frac{(2k+1)\pi}{\alpha},\qquad
 d(\eta)=\frac2\alpha\atanh\sqrt\eta\le\frac12.
\]
At \(\eta=0\), the multiplier has double real zeros \(x_k\).
At \(\eta>0\), they are the simple zeros \(x_k\pm\ii d(\eta)\).
Any coincidence with an \(H\)-zero adds multiplicities. The fixed
lattice spacing gives a common unit-count bound and hence a common
full count after adding the canonical zeros.

The finite-factor argument of Section~\ref{sec:green}, applied
periodically to \(M_\eta\), gives
\begin{equation}\label{eq:lattice}
 \Delta_\eta(X):=\calJ[f_\eta](X)-\calJ[f_0](X)
       =\sum_{k\in\mathbb Z}G(X-x_k,d(\eta)),\qquad
 \calJ[f_0]=\calJ[H].
\end{equation}
One can justify the periodic limit directly on the cylinder: after
subtracting the Green singularities, the remaining function is
harmonic, has the correct boundary values and is bounded. Uniqueness
fixes its harmonic part. Alternatively \eqref{eq:G-decay} gives
absolute convergence of the periodized kernels. The real doubles of
\(M_0\) have zero interior contribution.

The nearest lattice point is at distance at most \(\ell/2\).
Using \eqref{eq:margin} for that point and
\eqref{eq:G-decay} for the complete lattice gives
\begin{equation}\label{eq:lattice-bounds}
 \frac{\sin a}{\cosh(a\ell/2)}\,d
 \le\Delta_\eta(X)
 \le K_G\coth(a\ell/2)\,d .
\end{equation}
The upper bound uses
\(\sup_X\sum_k\ee^{-a|X-x_k|}=\coth(a\ell/2)\).
Finally
\[
 \frac2\alpha\sqrt\eta\le d(\eta)
 \le\frac{2}{\alpha(1-\eta_0)}\sqrt\eta,
\]
which proves the two-sided Green scale.

It remains to verify the resolvent observations. With
\(R_\eta=f_\eta'/f_\eta\), \(R_*=f_0'/f_0\), and
\(B_\eta=(1+\eta)/(1-\eta)=\cosh(\alpha d)\), exact subtraction yields
\begin{equation}\label{eq:resolvent-difference}
 R_\eta-R_*=
 \frac{\alpha(B_\eta-1)\sin(\alpha z)}
 {(B_\eta+\cos(\alpha z))(1+\cos(\alpha z))}.
\end{equation}
The canonical resolvent cancels, but the multiplier's zeros remain
fully represented in \eqref{eq:lattice}. For \(y>1/2\) and
\(1\le B\le\cosh(\alpha/2)\),
\[
 |B+\cos(\alpha(x+\ii y))|\ge\cosh(\alpha y)-B,\qquad
 |\sin(\alpha(x+\ii y))|\le\cosh(\alpha y).
\]
For the first inequality, set \(t=\cos(\alpha x)\); the squared
modulus is \(t^2+2B\cosh(\alpha y)t+B^2+\sinh^2(\alpha y)\),
minimized at \(t=-1\).

Fix \(b>1\) and \(0<\varrho<b-1\), and define
\[
 C_{b,\varrho}=
 \frac{2\alpha\cosh(\alpha(b+\varrho))}
 {(1-\eta_0)
 [\cosh(\alpha(b-\varrho))-\cosh(\alpha/2)]
 [\cosh(\alpha(b-\varrho))-1]}.
\]
On every radius-\(\varrho\) disk centered at \(x+\ii b\),
\eqref{eq:resolvent-difference} is bounded by \(C_{b,\varrho}\eta\).
The actual resolvents are defined there by the zero strip.
Cauchy's estimate gives
\begin{equation}\label{eq:jet-change}
 \sup_x|\partial_z^j(R_\eta-R_*)(x+\ii b)|
 \le j!\varrho^{-j}C_{b,\varrho}\eta.
\end{equation}
Any fixed sum, maximum, or ordinary finite-dimensional norm of the
relative source distance and finitely many such jet seminorms is
therefore at most \(C_D\eta\), with weights independent of \(\eta\).
A nondecreasing \(\omega(t)=o(\sqrt t)\) would imply
\(c\sqrt\eta\le\omega(C_D\eta)=o(\sqrt\eta)\), a contradiction.
\end{proof}

The baseline \(\Psi_0\) is itself a translated-source mixture, not
theta. The theorem concerns continuity on that enlarged information
class. It is not a discontinuity statement, a universal upper
Hölder theorem, or an obstruction restricted to the literal theta
source. It leaves open \(d^2\)-scale source resolution, contours
approaching the event with their denominator costs paid, and genuinely
arithmetic signed estimates. The explicit factorial in
\eqref{eq:jet-change} also prevents interpreting fixed-order
control as a free growing-order theorem.


\section{Quadratic comparison and the short-mollifier floor}\label{sec:quadratic}

\subsection{The exact mean and the multiplier's zeros}

Fix \(0<\theta<1/2\), put \(S=T/2\), and consider
\begin{equation}\label{eq:A-class}
 A(s)=\sum_{n\le N}a_nn^{-s},\qquad
 N\le S^\theta,\quad a_1=1,\quad a_n\in\C,\quad |a_n|\le1.
\end{equation}
For each \(T\), one polynomial is used on the whole boundary and for
all centers. Set
\[
 F_A(z)=F(z)A(s(z)),\quad c_A(X)=F_A(X+4\ii),\quad
 \epsilon_A(X)^2=\int\left|\frac{F_A}{c_A(X)}-1\right|^2\dd\mu_X .
\]
The denominator is the actual center value, not a surrogate constant.

The fractional-part identity \eqref{eq:fractional} applies here.
It follows by integrating the step-counting formula for zeta initially
on \(\Re s>1\); the displayed integral then continues the pole-cancelled
field to \(\Re s>0\). At \(s=1\) its value is one. If
\(C_A=\sum_{n\le N}|a_n|n^{-1/2}\le2\sqrt N\), then on the closed
strip \(0\le\Im z\le8\),
\(|F_A(x+\ii y)|\le C_A(8+|x|)\).
Poisson representation of its real and imaginary parts therefore gives
\begin{equation}\label{eq:mean}
 \int F_A\dd\mu_X=c_A(X).
\end{equation}
Polynomial growth is killed by the exponentially small vertical-side
harmonic measures.

At \(s_c=5/2-\ii X/2\),
\[
 2-\zeta(5/2)\le|A(s_c)|\le\zeta(5/2),\qquad
 \frac{\zeta(5)}{\zeta(5/2)}\le|\zeta(s_c)|\le\zeta(5/2),\qquad
 \frac35\le|q(s_c)|\le1.
\]
The Euler lower bound uses
\(\prod_p(1+p^{-5/2})^{-1}=\zeta(5)/\zeta(5/2)\).
Since \(\zeta(5/2)\le5/3\), the constants
\begin{equation}\label{eq:center-bounds}
 c_*=\frac35\frac{\zeta(5)}{\zeta(5/2)}(2-\zeta(5/2))>0,\qquad
 U_*=\zeta(5/2)^2
\end{equation}
satisfy \(c_*\le|c_A(X)|\le U_*\).
Expanding the square using \eqref{eq:mean} yields
\begin{equation}\label{eq:variance}
 V_A(X):=\int|F_A|^2\dd\mu_X-|c_A(X)|^2
 =|c_A(X)|^2\epsilon_A(X)^2\ge0,\qquad
 \epsilon_A(X)^2\ge V_A(X)/U_*^2.
\end{equation}

\begin{proposition}[Quadratic upper comparison]\label{prop:quadratic}
With the class \eqref{eq:A-class}, all multiplier boundary logarithms
are absolutely integrable and
\begin{equation}\label{eq:quadratic}
 0\le J_0(X)\le\calJ[F_A](X)
 =J_0(X)+\calJ[A\circ s](X)
 \le\frac12\log(1+\epsilon_A(X)^2).
\end{equation}
Zeros of \(A\), including boundary zeros, are retained.
\end{proposition}
\begin{proof}
For \(0<\delta\le1\), the bounded subharmonic function
\[
 u_\delta(z)=\frac12\log(|A(s(z))|^2+\delta^2)
\]
has upper bound \(M=\log(C_A+1)\).
Subharmonic Poisson majorization implies
\[
 0\le\int(M-u_\delta)\dd\mu_X
 \le M-u_\delta(X+4\ii)\le M-\log|A(s_c)|.
\]
Monotone convergence as \(\delta\downarrow0\) proves integrability
of the negative logarithm; its positive part is bounded. It also gives
\(\calJ[A\circ s]\ge0\). In the exact factorization
\(F_A=DH_0(A\circ s)\), the multiplier \(D\) has zero Green
functional by Section~\ref{sec:prime}. Add the logarithms and use
concavity of \(\log\):
\[
 \calJ[F_A]=\frac12\int\log|F_A/c_A|^2\dd\mu_X
 \le\frac12\log\int|F_A/c_A|^2\dd\mu_X
 =\frac12\log(1+\epsilon_A^2).
\]
No zero of \(A\) was divided away, and the zeta pole was already
cancelled in \(F_A\).
\end{proof}

\subsection{The imported mean-value input and an elementary first moment}

We use the unconditional lower-bound part of
Radziwi{\l}{\l}'s Proposition B, attributed there to Soundararajan
\cite{Radziwill2012MollifyingV3}. For fixed \(\theta<1/2\),
normalization \(a_1=1\), and the stated common coefficient-growth
bounds, it gives
\begin{equation}\label{eq:primary-moment}
 \frac1S\int_S^{2S}
 |1-\zeta(1/2+\ii t)A(1/2+\ii t)|^2\dd t
 \ge\frac1\theta-o_\theta(1).
\end{equation}
This is an imported mean-value theorem, not proved here. We use the
pinned version 3 and its complex-coefficient statement, not its
RH-conditional results.

For \eqref{eq:A-class}, the lower remainder is uniform in \(A\).
Pad shorter polynomials with zeros to \(\lfloor S^\theta\rfloor\).
All coefficient-growth constants are at most one. If a uniform lower
remainder failed, offending arrays along a sequence of heights,
extended by \(A=1\) at the other heights, would form an admissible
family contradicting \eqref{eq:primary-moment}. This does not assert
an additive uniform asymptotic for arbitrarily large quadratic forms.

To convert the error moment to a square mean, an additional first
moment is essential. Put \(B(s)=\zeta(s)A(s)\).
Shift its rectangle from \(\Re s=1/2\) to \(\Re s=2\) between
imaginary heights \(S\) and \(2S\). The zeta pole is outside.
For the integer \(M=\lceil2S\rceil\), the exact truncation
\[
 \zeta(s)=\sum_{n\le M}n^{-s}+\frac{M^{1-s}}{s-1}
         -s\int_M^\infty\{v\}v^{-s-1}\dd v
\]
gives \(|\zeta(s)|=O(\sqrt S)\) on both horizontal edges,
uniformly for \(1/2\le\Re s\le2\). Since \(|A(s)|\le2\sqrt N\),
the two edges cost \(O(\sqrt{SN})\).
On \(\Re s=2\), absolute convergence gives
\[
 B(s)=\sum_{n\ge1}b_nn^{-s},\qquad
 b_n=\sum_{\substack{d\mid n\\d\le N}}a_d,\quad b_1=1,\quad
 |b_n|\le\tau(n).
\]
The nonconstant mean has absolute size at most
\(2S^{-1}\sum_{n\ge2}\tau(n)n^{-2}/\log n=O(S^{-1})\).
Consequently
\begin{equation}\label{eq:first-moment}
 \frac1S\int_S^{2S}B(1/2+\ii t)\dd t
 =1+O_\theta(S^{(\theta-1)/2}).
\end{equation}
Expanding \(|1-B|^2\) in \eqref{eq:primary-moment} now gives
\begin{equation}\label{eq:square-mean}
 \frac1S\int_S^{2S}|B(1/2+\ii t)|^2\dd t
 \ge1+\frac1\theta-o_\theta(1),
\end{equation}
uniformly in the declared coefficient class.

\subsection{Safe-line subtraction and the padded convolution}

For any interval \([u,u+S]\), absolute convergence gives
\begin{equation}\label{eq:safe-mean}
 \frac1S\int_u^{u+S}|B(5/2+\ii t)|^2\dd t
 =\sum_{n\ge1}|b_n|^2n^{-5}+O(S^{-1}),
\end{equation}
uniformly in \(u,A\).
Indeed the absolute off-diagonal error is at most
\[
 \frac4S\sum_{m<n}
 \frac{\tau(m)\tau(n)}{(mn)^{5/2}\log(n/m)}
 \le\frac{16}S\sum_{m\ge1}\frac{H_m}{m^3}<\infty,
 \quad H_m=\sum_{j=1}^m\frac1j.
\]
Here \(\log(n/m)\ge(n-m)/n\), \(\tau(n)\le2\sqrt n\), and
\(\sum_{k\ge1}1/[k(m+k)]=H_m/m\).
The diagonal is bounded by
\begin{equation}\label{eq:diagonal}
 \sum_{n\ge1}\tau(n)^2n^{-5}
 \le1+4\sum_{n\ge2}n^{-4}
 \le1+4\left(\frac1{16}+\int_2^\infty v^{-4}\dd v\right)
 =\frac{17}{12}.
\end{equation}

Let
\begin{equation}\label{eq:padded}
 L_T=a^{-1}\log T,\qquad I_T=[T-L_T,2T+L_T].
\end{equation}
Since \(|q(5/2+\ii t)|\le1\), substitution \(X=2t\), with
coefficient conjugation for the negative imaginary direction, gives
\[
 \frac1T\int_T^{2T}|c_A(X)|^2\dd X\le17/12+O(T^{-1}).
\]
The two extra intervals have total length \(2L_T\), and
\(|c_A|^2\le U_*^2\). Normalizing by the actual length therefore yields
\begin{equation}\label{eq:padded-center}
 \frac1{|I_T|}\int_{I_T}|c_A(X)|^2\dd X\le17/12+o(1).
\end{equation}

\begin{theorem}[Uniform quadratic-residual floor]\label{thm:floor}
For every fixed \(0<\theta<1/2\), the class \eqref{eq:A-class}
and interval \eqref{eq:padded} satisfy
\begin{equation}\label{eq:floor}
 \liminf_{T\to\infty}\inf_A
 \frac1{|I_T|}\int_{I_T}\epsilon_A(X)^2\dd X
 \ge\frac{1/(2\theta)-11/12}{\zeta(5/2)^4}
 >\frac1{12\zeta(5/2)^4}.
\end{equation}
In particular, for all sufficiently large \(T\), uniformly in \(A\),
the average exceeds \(1/[24\zeta(5/2)^4]\).
No numerical onset is asserted.
\end{theorem}
\begin{proof}
Integrate the complete nonnegative boundary square in
\eqref{eq:variance} over \(I_T\), using Tonelli.
Retain only the lower boundary and \(x\in[T,2T]\).
For every such \(x\), the interval
\([x-L_T,x+L_T]\) lies in \(I_T\), and hence
\[
 \int_{I_T}K(x-X)\dd X\ge k_{L_T},\qquad
 k_L=\int_{-L}^LK(u)\dd u=\frac1\pi\arctan(\sinh(aL))
 \ge\frac12-\frac{\ee^{-aL}}{4a}.
\]
Thus \(k_{L_T}\ge1/2-O(T^{-1})\).
On the lower boundary \(|q|=1\).
The conjugate-coefficient polynomial
\(\widetilde A(s)=\sum\overline{a_n}n^{-s}\) is still admissible.
Applying \eqref{eq:square-mean} to it and substituting \(x=2t\)
gives
\[
 \int_T^{2T}|F_A(x)|^2\dd x
 =2\int_S^{2S}|\zeta(1/2+\ii t)\widetilde A(1/2+\ii t)|^2\dd t
 \ge T(1+1/\theta-o_\theta(1)).
\]
Subtract \eqref{eq:padded-center}. Since \(T/|I_T|\to1\),
\[
 \frac1{|I_T|}\int_{I_T}V_A(X)\dd X
 \ge\frac12(1+1/\theta)-\frac{17}{12}-o_\theta(1)
 =\frac1{2\theta}-\frac{11}{12}-o_\theta(1).
\]
All lower errors are common to the coefficient class. Divide by
\(U_*^2=\zeta(5/2)^4\) using \eqref{eq:variance}, then take the
infimum and liminf. The strict assumption \(\theta<1/2\) gives
the final claims.
\end{proof}

\subsection{The real mixed term and its corrected verification}

For clarity we include the elementary verification used in the branch's
assessment \cite{OGCorrection,OGAssessment}. It corrects an intermediate
source attribution and does not replace the imported mean-value
theorem. With \(N=\lfloor S^\theta\rfloor\), define
\[
 w_\ell=\frac{\varphi(\ell)}{\ell^2},\quad
 y_\ell=\sum_{d\le N/\ell}\frac{a_{d\ell}}d,\quad
 G_N=\sum_{\ell\le N}\frac{\mu(\ell)^2}{\varphi(\ell)},\quad
 z_\ell=\frac{\mu(\ell)\ell}{G_N\varphi(\ell)},
\]
\[
 e_\ell=y_\ell-z_\ell,\qquad
 E_N=\sum_{\ell\le N}w_\ell|e_\ell|^2.
\]
Möbius inversion and \(a_1=1\) imply
\begin{equation}\label{eq:orthogonality}
 \sum_{\ell\le N}\frac{\mu(\ell)e_\ell}{\ell}=0.
\end{equation}
The relevant \emph{real} mixed form is
\[
 \mathcal C_N=
 2\Re\sum_{p^k\ell\le N}\frac{\log p}{p^k}w_\ell
       \{z_\ell\overline{e_{p^k\ell}}+z_{p^k\ell}e_\ell\}.
\]
Call the two directed sums \(U_N,V_N\).
Reindexing the first and using
\(\sum_{p^k\mid m}\mu(m/p^k)\log p=-\mu(m)\log m\) gives
\[
 U_N=-\frac1{G_N}\sum_{m\le N}
             \frac{\mu(m)\log m}{m}\overline{e_m}.
\]
In the second, \(\mu(p^k\ell)\ne0\) requires \(k=1\), \(p\nmid\ell\).
Thus the exact expression is
\[
 V_N=-\frac1{G_N}\sum_{\ell\le N}\frac{\mu(\ell)e_\ell}{\ell}
         P_\ell(N/\ell),\qquad
 P_\ell(x)=\sum_{\substack{p\le x\\p\nmid\ell}}\frac{\log p}{p-1}.
\]
The classical prime-sum estimate used in the primary argument gives
\[
 P_\ell(N/\ell)=\log(N/\ell)+r_\ell,\qquad
 |r_\ell|\ll\log\log(3N).
\]
Its uniformity follows by splitting excluded prime divisors at
\(Y=\log(3N)\): the small ones cost \(O(\log Y)\), and the large
ones cost at most \(2\log N/Y\).
The unrestricted sum is \(\log x+O(1)\); replacing
\(1/(p-1)\) by \(1/p\) costs a convergent series.

After \eqref{eq:orthogonality}, the logarithmic main terms of \(U_N\)
and \(V_N\) cancel \emph{only after taking real parts}. Therefore
\[
 \mathcal C_N=-\frac2{G_N}\Re
 \sum_{\ell\le N}\frac{\mu(\ell)r_\ell e_\ell}{\ell}.
\]
Cauchy--Schwarz and the classical \(G_N\sim\log N\) give
\begin{equation}\label{eq:real-mixed}
 |\mathcal C_N|
 \le\frac2{G_N}
 \left(\sum_{\ell\le N}\frac{\mu(\ell)^2r_\ell^2}{\varphi(\ell)}\right)^{1/2}
 \sqrt{E_N}
 \ll\frac{\log\log(3N)}{\sqrt{\log N}}\sqrt{E_N}.
\end{equation}
For a fixed \(\nu>0\), if the last coefficient is \(b_N\), then
\[
 b_N\sqrt{E_N}
 \le\nu\log S\,E_N+\frac{b_N^2}{4\nu\log S}.
\]
The scalar remainder tends to zero uniformly in the coefficients.
The two elementary prime-sum asymptotics are classical inputs already
used in the source argument. Equation~\eqref{eq:real-mixed} is the
direct real calculation recorded by this project; it is not quoted
as the literal printed formula of Lemma 13 in version 3, which displays
\(E_N\) rather than its square root. Nor does it bound the modulus
of the unprojected complex mixed sum.

\subsection{What the floor excludes}

By Corollary~\ref{cor:margin} and Proposition~\ref{prop:quadratic},
a fully acquired inequality
\(\epsilon_A(X)^2<\exp(2m(d,r))-1\) would exclude the specified
nonreal zero. The positive average floor prevents these sufficient
residual certificates from tending uniformly to zero throughout
\(I_T\) using one polynomial in \eqref{eq:A-class}.
It gives no lower bound for \(J_0\), because the upper comparison
\eqref{eq:quadratic} cannot be reversed. It supplies neither failure
at every center nor a positive density of failing centers.
Longer polynomials, \(\theta\ge1/2\), adaptive center choices,
unpadded averages and signed logarithmic estimates are outside
the theorem.


\section{Complete source acquisition for the quadratic comparison}\label{sec:acquisition}

The preceding floor concerns the exact field. To state honestly what
a finite implementation of its sufficient comparison would require,
we retain the complete acquisition estimate from
\cite{OGMollifier,OGCorrection}. This is an analytical error bound;
no new numerical certificate is run in this integration.

\subsection{Reuse of the complete source payment}

Use exactly the reflected source and radial cutoff of
Section~\ref{sec:signed-acquisition}. Equation~\eqref{eq:delta-H}
pays its two independent errors. For the fixed polynomial of
\eqref{eq:A-class}, put \(C_A=\sum_{n\le N}|a_n|n^{-1/2}\).
On the closed strip, \(|A(s(z))|\le C_A\). Hence
\[
 \mathcal B_A(X,R)=C_A\mathcal B(X,R),\qquad
 |F_A-DA\widehat H|\le\mathcal B_A\delta_H
\]
on the compact region of \eqref{eq:multiplier}. All following boundary
tails concern the exact field, so the compact multiplier is not used
at infinity.

\subsection{Both exact boundary tails}

Let \(\mu_{X,R}\) be the restriction of \(\mu_X\) to \(|x-X|\le R\).
Equation~\eqref{eq:fractional} gives
\[
 |F_A(x)|\le C_A(2+|x|),\qquad
 |F_A(x+8\ii)|\le\zeta(9/2)^2.
\]
Writing \(B_X=2+|X|\), one obtains
\begin{align}\label{eq:quadratic-tail}
 \norm{F_A/c_A-1}_{L^2(\mu_X-\mu_{X,R})}^2
 &\le\tau_A(X,R),\nonumber\\
 \tau_A(X,R)
 &=\frac{\ee^{-aR}}{2a}\left\{
 \frac{C_A^2}{c_*^2}
 \left[(B_X+R)^2+\frac{2(B_X+R)}a+\frac2{a^2}\right]
 +\frac{\zeta(9/2)^4}{c_*^2}+2\right\}.
\end{align}
Indeed \(|v-1|^2\le2|v|^2+2\), and
\(K(u)\le\ee^{-a|u|}/8\).
Integrate the lower-boundary quadratic polynomial on both spatial
tails and then the upper-boundary constant on both tails.
The upper square requires the fourth power of \(\zeta(9/2)\).
All tails here belong to the exact field.

\subsection{Center acquisition and denominator propagation}

For an integer \(M_c\ge2\), the safe Euler logarithm has error
\[
 \left|\log\zeta(s_c)-\sum_{p^k\le M_c}\frac{p^{-ks_c}}k\right|
 \le\beta_{M_c}=\frac23M_c^{-3/2}.
\]
Keep the finite polynomial \(A(s_c)\) exactly and set
\[
 \widetilde c=q(s_c)A(s_c)
       \exp\left(\sum_{p^k\le M_c}\frac{p^{-ks_c}}k\right),\qquad
 |c_A-\widetilde c|\le
 \eta:=U_*(\ee^{\beta_{M_c}}-1).
\]
The modulus of the finite-log exponential is at most \(\zeta(5/2)\).
Future arithmetic evaluation error is added to \(\eta\).
No logarithm of \(A\) is used.

Let
\[
 \ell_c=|\widetilde c|-\eta>0,\qquad
 \mu_R=\mu_{X,R}(\partial{\rm strip})
       =\frac2\pi\arctan(\sinh(aR)),\qquad
 A_R=\norm{\widehat F_A/\widetilde c-1}_{L^2(\mu_{X,R})}.
\]
Then
\begin{equation}\label{eq:full-quadratic-acquisition}
 \epsilon_A(X)\le A_R+
 \frac{\eta}{\ell_c}(A_R+\sqrt{\mu_R})
 +\frac{\mathcal B_A\delta_H}{c_*}+\sqrt{\tau_A(X,R)}.
\end{equation}
For the proof, first split the complete \(L^2\) norm into compact
and remote pieces using the triangle inequality. On the compact
piece, add and subtract \(\widehat F_A/c_A\) and
\(\widehat F_A/\widetilde c\). The field error is at most
\(\mathcal B_A\delta_H/c_*\), since the restricted mass is at most one.
The denominator error is
\[
 \frac{|c_A-\widetilde c|}{|c_A|}
     \norm{\widehat F_A/\widetilde c}
 \le\frac{\eta}{\ell_c}(A_R+\sqrt{\mu_R}),
\]
which proves \eqref{eq:full-quadratic-acquisition}.

For each fixed \(X,A\) and a positive acquisition budget \(h\),
first choose \(R\) to make \(\sqrt{\tau_A}<h/3\).
For this \(R\), choose \(N_{\rm src},V\) so that
\(\mathcal B_A\delta_H/c_*<h/3\).
Finally choose the integer center cutoff and arithmetic precision so that
\(\eta<c_*/4\) and
\(2\eta\norm{\widehat F_A}_{L^2(\mu_{X,R})}/c_*^2<h/3\).
The exact denominator difference is then bounded by
\[
 \frac{\eta\norm{\widehat F_A}}{|c_A||\widetilde c|}
 \le\frac{4\eta\norm{\widehat F_A}}{3c_*^2}<h/3.
\]
This sharper expression, rather than requiring the conservative
surrogate in \eqref{eq:full-quadratic-acquisition} to meet the same
allocation, proves the ordered finite-acquisition assertion.
If the norm vanishes, the denominator error vanishes.
Every tuple can be treated this way, including a prescribed sequence
of growing heights or shrinking defect distances. This supplies no
common cutoff, affordable complexity bound, small residual value,
or evaluated exclusion.


\section{Parallel inputs and the scale of a single defect}
\label{sec:parallel}

The operator parent developed several other observables before the Green
branch separated. Their transferable content explains both useful
comparisons and the remaining gap. Localized correlations (OC) and
collision constraints (OA) retain separate manuscripts, registers and
research questions \cite{GNPCorrelations,GNPCollisions}. We do not treat
progress in one observable as an estimate of another.

\subsection{Smooth transverse moments and the Green scale}
For a nonnegative compactly supported \(C^2\) function \(a\), set
\(Q_a=\sum_j a(x_j)y_j^2\), counting every root and multiplicity.
A conjugate pair at \(X\pm\ii d\), where \(a(X)=1\), contributes
\(2d^2\). In reverse Newman time, at simple relevant roots the paired
zero dynamics gives
\begin{equation}\label{eq:smooth-kernel}
 \frac{\dd Q_a}{\dd\tau}=2\sum_{j<k}'K_f(z_j,z_k),\qquad
 K_f=\int_0^1\frac{2a(x_t)e^2-a''(x_t)y_t^2d_x^2}
                         {d_x^2+e^2}\dd t,
\end{equation}
where \(f=a(x)y^2\), \(d_x=x_j-x_k\), \(e=y_j-y_k\) and
\(z_t=z_k+t(z_j-z_k)\). To derive it, differentiate \(f\) along the
root velocities and pair the two interaction rows. The mixed Hessian
terms cancel. In \(|y|\le1\),
\(|K_f|\le\max(2\|a\|_\infty,\|a''\|_\infty)\): no internal gap
appears in this bound. The prime requires complete opposite-pair
summation. The count \eqref{eq:count} makes its far velocity tail
\(O(Z\log U/U)\) on \(|z|\le Z\) after cutoff \(U\ge2Z+2\).
This follows by combining \((z-\zeta)^{-1}+(z+\zeta)^{-1}
=2z/(z^2-\zeta^2)\) before partial summation.

The parent proves, from the canonical unit counts and safe-line
resolvent estimates, that a common \(A\ge1\) gives
\begin{equation}\label{eq:smooth-modulus}
 |Q_{X,R}(\lambda_2)-Q_{X,R}(\lambda_1)|
 \le344A^2\|\chi\|_{(2)}R^{-2}\log^2(2+X)|\lambda_2-\lambda_1|,
\end{equation}
for \(X\ge10\), \(0<R\le1\), \(0\le\lambda_i\le1/2\), where
\(a(x)=\chi((x-X)/R)\), \(0\le\chi\le1\),
\(\chi\in C_c^2([-1,1])\), and
\(\|\chi\|_{(2)}=\sum_{j=0}^2\|\chi^{(j)}\|_\infty\)
\cite{SmoothLocalizationProof,SmoothLocalizationAssessment}.
The proof groups the finitely many roots near the support and compares
the exterior with \(H'/H(x+4\ii)\). The exterior difference has an
absolute inverse-square tail; the unpaired harmonic sum is not bounded
absolutely. Local heat-polynomial splitting isolates the finitely many
multiple events in a fixed compact space-time box. Integrating the same
derivative bound between them and using multiset continuity proves
\eqref{eq:smooth-modulus} through collisions. This argument allows a
derivative jump, but no atom in the first distributional derivative.

In particular an off-axis pair at \(\lambda=0\) yields the conditional
threshold lower bound
\(\Lambda\ge2d^2/[344A^2\|\chi\|_{(2)}\log^2(2+X)]\), with
\(R=1\) and \(\chi(0)=1\). At the threshold every zero is real,
since a nonreal root would persist above it. Apply the modulus between
zero and \(\Lambda\). This lower bound is not a contradiction and is
not an upper bound on the lifetime of every possible defect.

The Green observable has a different local scale. Differentiating its
explicit kernel and substituting \(t=\sinh(au)\) gives, for \(0\le y\le1\),
\[
 \int_\R\partial_yG(u,y)\dd u
 =\int_\R\frac{\cos(ay)}{t^2+\cos^2(ay)}\dd t=\pi.
\]
Since \(G(u,0)=0\),
\begin{equation}\label{eq:Green-mass-depth}
 \int_\R G(u,y)\dd u=\pi y.
\end{equation}
Thus a defect of depth \(d\) has a signal of order \(d/T\) in a
normalized center average over length \(T\). At a newly born double
pair the positive ordinate is of order \(\sqrt\tau\), so the Green
signal has that scale, while the smooth moment has scale \(\tau\).
Collision-stable Lipschitz bounds for the latter do not transfer to
the former.

\subsection{What uniform density and all-order averages do provide}
The parent also obtains, for each fixed \(0<b<1/2\),
\begin{equation}\label{eq:inherited-density}
 N_\lambda(T,v):=\#\{j:T\le x_j\le2T,\ |y_j|>v\}
 \le B_bT\log T\,\ee^{-bv\log T},
\end{equation}
uniformly for \(T\ge8\), \(0\le\lambda\le1/2\), \(0<v<1\)
\cite{SmoothLocalizationProof,SmoothLocalizationAssessment}.
Its imported time-zero input is the density estimate of
Maples--Rodgers, with its constant uniform as the strip boundary
approaches the critical line \cite[Theorem 2.4 and Section 4 (18)]{MaplesRodgers2015}.
Its transport step uses a broad window \(a_T\) and an even convex
\(g\) with \(g(0)=0\). In Newman time the favorable \(a_Tg''\)
term of the chord kernel can be discarded, and conjugate exterior
roots have a favorable vertical sign at horizontal separation at least
two. The remaining cost is
\(D_\chi\|g\|_\infty(\lambda_2-\lambda_1)\log^3T\), independent
of \(g'\) and \(g''\). Take \(g(y)=\cosh(b\log T\,y)-1\), apply
the time-zero tail by layer cake, and absorb
\(T^{b-1}\log^2T\). Markov's inequality gives
\eqref{eq:inherited-density}; small \(v\log T\) is paid by the unit count.

Layer cake now gives the explicit consequence
\[
 \sum_{T\le x_j\le2T}|y_j|^{2p}
 \le B_bT\log T\frac{(2p)!}{(b\log T)^{2p}},\qquad p\ge1.
\]
For local windows the root population is at most \(J\log T\), with
a fixed canonical \(J\). Weighted H\"older, followed by integration
over centers (each root occupies center measure at most \(R\int\chi\)),
therefore yields
\begin{equation}\label{eq:moment-all-orders}
 \frac1T\int_T^{2T}Q_{X,R}^p\dd X
 \le C_bR A_b^p\frac{(2p)!}{(\log T)^p}.
\end{equation}
A fixed enlargement of the dyadic band pays support endpoints. The
constants can be chosen independent of \(p,\lambda,T,R\) in the
declared ranges. This is a useful simultaneous all-order statement;
the factorial cost must remain when choosing a growing order.

None of these bounds removes the exceptional-center problem. For
example \(p=1\) bounds a normalized smooth average at order
\(1/\log T\), whereas one pair contributes only order \(d^2/T\).
Similarly, applying \eqref{eq:inherited-density} to the first absolute
transverse moment gives a Green average allowance of constant order,
too large to resolve \eqref{eq:Green-mass-depth}'s single-defect scale.
The inherited resolution assessment makes this comparison for its
specific density-based allowance \cite{OGResolution}. It is a scoped
failure of that allowance, not a lower bound for the actual Green value
or an impossibility theorem for all statistics.

\subsection{A field-based comparison: Gaussian Laguerre defects}
Another parent observable is
\[
 A_n(\lambda,X,v)=\int_\R w_v(x-X)[(f_n')^2-f_nf_n'']\dd x,
 \quad f_n=\partial_x^nH_\lambda,\quad
 w_v(t)=\frac{\ee^{-t^2/(2v)}}{\sqrt{2\pi v}},\quad v>0.
\]
The complete source gives the exact signed form
\begin{equation}\label{eq:Gaussian-source}
 A_n=\frac18\int_{\R^2}(-pq)^n(p-q)^2W_\lambda(p)W_\lambda(q)
             \ee^{\ii X(p+q)-v(p+q)^2/2}\dd p\dd q.
\end{equation}
Symmetrize \((f_n')^2-f_nf_n''\) under \(p\leftrightarrow q\) to get
\((p-q)^2/8\); Gaussian Fourier integration gives the last factor.
Complete source moments justify Fubini. Positivity of the theta factors
does not remove the signed phase or the factor \((-pq)^n\).

Direct differentiation gives
\begin{equation}\label{eq:Gaussian-hierarchy}
 \partial_\lambda A_n=-\partial_{XX}A_n+2A_{n+1},\qquad
 \partial_v A_n=\tfrac12\partial_{XX}A_n.
\end{equation}
If \(M_j(L)=\int|u|^j\ee^{Lu^2}\Phi(u)\dd u\), then
\(|pq|^n\le n!\epsilon^{-n}\ee^{\epsilon(p^2+q^2)/2}\) proves
\( |A_n|\le n!\epsilon^{-n}M_0(L+\epsilon/2)M_2(L+\epsilon/2)/4\)
for \(\lambda\le L\). Consequently the characteristic
\(v=v_0+2h>0\) has the convergent expansion
\[
 A_n(\lambda_0+h,X,v_0+2h)
 =\sum_{k\ge0}\frac{2^kh^k}{k!}A_{n+k}(\lambda_0,X,v_0).
\]
It can also be proved directly by expanding \(\ee^{-2hpq}\) in
\eqref{eq:Gaussian-source}, with complete moment domination. The parent
records explicit joint order, arithmetic and radial remainders and the
assessment of those fixed statements \cite{GaussianLaguerreProof,LocalizedMetricsAssessment}.

At a real event of finite multiplicity \(m\), let
\(a_m=H_{\lambda_*}^{(m)}(x_*)\ne0\). Differentiating the local heat
polynomial \(m-2\) times leaves
\(g(t,y)=a_m(y^2/2+t)+\text{remainder}\), for
\(g=\partial_x^{m-2}H_{\lambda_*-t}(x_*+y)\).
Uniform real-line bounds on \(g_{yyy},g_{yyyy}\), supplied by the
complete moments, bound the remainder by
\(C(|y|^3/6+t|y|)+Dt^2/2\); the first two spatial derivative errors
are bounded by \(C(y^2/2+t)\) and \(C|y|+Dt\).
Substitution in \((g_y)^2-gg_{yy}\) and Gaussian integration with
variance \(v=\kappa t\) therefore prove
\begin{equation}\label{eq:Gaussian-event}
 A_{m-2}(\lambda_*-t,x_*,\kappa t)
 =a_m^2t(\kappa/2-1)+O_{a_m,C,D,\kappa}(t^{3/2}),\quad0<\kappa<2.
\end{equation}
The explicit remainder in the cited parent theorem pays the full
Gaussian exterior. At each fixed event the expression is negative
for sufficiently small \(t\). A uniform sign-prevention theorem would
still need a source-specific estimate at every relevant order and
shrinking variance, with event amplitude paid. The hierarchy is a
diagnostic, not such a theorem.

For \(n=0\), putting \(Z=\int w_vH_\lambda^2\) and
\(E_w=\int|(H_\lambda\sqrt{w_v})'|^2\) also gives
\(A_0=2E_w-Z/(2v)\), by integration by parts. Thus \(E_w\ge0\)
is weaker than the signed threshold \(E_w\ge Z/(4v)\). This is the
same distinction between regular positive energy and a signed
zero-sensitive observable encountered in Section~\ref{sec:geometry}.

\subsection{The concrete collision input maintained by OA}
The real-pair gap law \eqref{eq:q-exact} has a finite-interval consumer
in OA \cite[Section 4]{GNPCollisions}. On a predeclared box of times,
centers and virtual real node pairs, acquire the deflated factor by
divided differences with a uniform model floor and derivative errors.
These must imply \(|g|\le A\), \(|g_x|\le K\), and an upper bound
\(\beta\) for the averaged exterior curvature in \eqref{eq:q-exact}.
The floor is required throughout that box, before selecting the unknown
trajectory. Values only at the initial roots do not supply it.

Here \(c=(r_++r_-)/2\) is the pair center and \(D\) bounds its half-gap.
Suppose initially \(|c(0)-X|\le R_0<R\) and
\(0<Q_-\le q(0)\le Q_+<4D^2\). If
\begin{equation}\label{eq:pair-survival}
 R_0+2A\Delta<R,\qquad Q_+\ee^{4K\Delta}<4D^2,\qquad
 Q_--8\int_0^\Delta\ee^{4\beta t}\dd t>0,
\end{equation}
the specified pair stays real, distinct and simple for
\(0\le t\le\Delta\). Indeed \(|c'|\le2A\), \(q'\le4Kq\), and
\(q'\ge-8-4\beta q\) give confinement and the positive lower bound
\[
 q(t)\ge\ee^{-4\beta t}
       \left(Q_--8\int_0^t\ee^{4\beta s}\dd s\right).
\]
At a proposed stopping time the root slopes remain at least
\(\sqrt q\) times the factor floor, so the implicit function theorem
continues the pair. This proves the closed-interval assertion, including
the endpoint. A failed sufficient inequality leaves the actual pair
undecided. For the quadratic heat model the gap condition is sharp;
bounded exterior curvature by itself cannot prevent the terminal
\(-8\) closing rate.

The alternative contour interface extracts \(g=U'/U\) as
\((2\pi\ii)^{-1}\int_\Gamma(H'/H)(\zeta)/(\zeta-z)\dd\zeta\): the
cluster divisor's two residues cancel. Ratio errors on the boundary
then pay errors in \(g\) and \(g_x\) by the contour length and buffer
distance. This requires a complete cluster count and boundary margin,
but no lower internal gap. The same analytical tools therefore serve
Green acquisition and collision-force acquisition with different
consumers. A time-zero Green value alone does not provide the
space-time floor required in \eqref{eq:pair-survival}.

\section{Evidence, controls and inherited numerical support}\label{sec:evidence}

\subsection{Control families and the meaning of the results}

The complete Green identity is an exact representation; the stability
and residual floors are quantitative scoped obstructions; and
Corollary~\ref{cor:canonical-log} is an explicit local acquisition
estimate. The source/operator dictionaries are representations. The
Riemann--Siegel logarithmic transfer is an acquisition estimate, and the
finite-pair survival statement is conditional on its stated box data. These roles must not be interchanged. In particular a
lower bound for a sufficient quadratic residual is not a lower
bound for \(J_0\).

The transition-scale outer and phase bounds are acquisition estimates.
The coherent moment, residual mean, density and
Theorem~\ref{thm:green-power-means} are analytical averaged results
with their stated quantifiers. Theorem~\ref{thm:b10-averaged} and
Corollary~\ref{cor:b10-green} quantify shorter translated intervals
at \(c=1\). Their proof recovers the exact frozen high arm from the
ordinary zeta Dirichlet series and uses the exact weighted residual
comparison.
The quadratic, phase-corrected and
fixed-reference tangent failures concern the specified allowances.
The new power-mean consequence is derived from the inherited moment
and density theorems, with a written same-assistant check; it is not
independent arithmetic information at exceptional centers. No numerical
test is evidence for an RH-equivalent reformulation.

The collision controls summarized in \cite{OGControls} give the
following analytical screen. No new control computation is required
for these premise comparisons.

\begin{center}\small
\begin{tabularx}{\textwidth}{@{}>{\raggedright\arraybackslash}p{0.24\textwidth}Y@{}}
\toprule
Control family & Premise not supplied for the canonical application\\
\midrule
Euler-factor and translated-source families &
The extra multiplier changes the literal theta/xi field. Canonical
prime subtraction and Euler-anchor constants cannot simply be assigned
to the altered function. The fixed-prime baseline in
Section~\ref{sec:conditioning} is itself not theta.\\
Weighted-square, finite-prefix and finite-prime-orbit controls &
Matched finite data do not supply the complete normalized theta
coefficients together with full reflection.\\
Reflected super-Gaussian control &
Reflection and positive Gram forms do not supply the integer-square
theta expansion or the unchanged xi/zeta identity.\\
Smooth shape and positive-primitive controls &
Their specified shape, endpoint and asymptotic data do not supply the
complete canonical arithmetic source.\\
Atomic, finite-profile and heat counter-kernels &
Selected anchors and moments do not supply the complete canonical
field, arithmetic arms, correction and coherent moment law. Continuous
extensions retain this limitation.\\
RH-conditional normalized endpoint controls &
Matching completion and residue does not supply ordinary theta/zeta
Dirichlet arithmetic. Their conditional status is retained.\\
Fixed theta amplitude with variable readout &
The literal cosine readout and corresponding physical Newman equation,
used in the complete field input, are absent.\\
Relaxed positive Green profiles &
A relaxed positive measure need not be the exact integer zero measure;
finite profile constraints do not supply the boundary field or its
residual and moment laws.\\
\bottomrule
\end{tabularx}
\end{center}

The generic local theorem can apply to an admissible control with its
own anchor and growth constants. If that control has an off-axis zero
in the relevant domain, accurate acquisition retains its positive
Green signal and cannot pass the corresponding strict exclusion
threshold. Failure of the literal theta identity does not imply
failure of every inequality derived from it. A new canonical sign
estimate would have to identify its quantitative distinguishing
premise.
These comparisons are made at the controls' actual witness domains;
no finite witness is moved into the shrinking-time band. The generic
harmonic, logarithmic, phase, Markov and layer-cake arguments remain
valid for any data satisfying their hypotheses. A control's failure
of canonical arithmetic does not establish failure of the final mean
conclusion; conversely a mean conclusion can coexist with isolated
collisions. No control certificate was rerun for this analytical screen.
In the short-interval argument, the Mellin projection requires the
literal ordinary zeta coefficients to recover the specified high arm.
Lattice resonances of the canonical phases are included in its
real quadratic sum.

\subsection{What the existing computations establish}
The analytical proofs above do not depend on a new numerical run.
The predecessor's records remain supporting evidence with their exact
finite scope:
\begin{itemize}
\item E06's interval-certified feasibility screen failed one sufficient
uniform-midpoint recipe at \(X=30,d=1/2,r=1\), under a million-panel
cap. It did not establish a lower bound on the true integration cost
or rule out other quadratures \cite{OGMidpoint}.
\item B05 r2's contour design uses 264 patches on \([-36,96]\),
degree-64 literal-source polynomials and 67584 product-integration
panels. Conditional on all its contour, count, factor and range gates
and numerical payments, it proves \(J_{\rm up}-J_0(30)<17/1000\).
The number bounds acquisition error, not \(J_0(30)\)
\cite{OGContourDesign}.
\item E08 certified a zero-free direct rectangle
\([29,31]+\ii[1/2,1]\) and three patches, with counts \(1,0,1\), at
128-bit working precision. The two count-one real-symmetric clusters
contain real simple roots. The other 261 patches and the full signed
Green evaluation were not executed. The rectangle was already within
established coverage \cite{OGPilot,OGPilotAssessment}.
\item A separate reuse of OC's complete thin zero-location census gives
\(0\le J_0(1500)<291/65536+2^{-157}<1/200\), at that census's
inherited winding-certification scope. Its input already provides
stronger local zero-location information; it is not acquisition of the
signed boundary logarithm \cite{OGSiblingReuse}.
\end{itemize}
None of these records is promoted here into an all-height result,
an independent comparison of full-method cost, or a new exclusion.

\subsection{Source, review and authorship provenance}
This manuscript was drafted for Edward Baker with substantial OpenAI
Codex LLM assistance in synthesis, elementary connecting derivations,
source checking and production. The actual model identifier and reasoning
effort for this founding chat are unavailable in its capture metadata;
recommended models in the research plan are not session provenance.
Where inherited notes record their actual configurations, those remain
attached to those notes rather than being assigned to this editor.
The 11 November integration and averaged-value derivation likewise
have no recorded model or effort in child chat0004's registry. The
older explicit B03--B05 dispatches were gpt-6-astra/ultra; E04--E06
and their internal assessors were gpt-6.1-sol/ultra. E05's assessor
contributed its twisted-moment route and E06's assessor its tangent
proof. These assessments share proof development and inherited inputs;
their source status is preserved in the cited records.

The 12 November B10 proof had a same-assistant check. A subsequent
internal audit in child chat0008 rederived the Mellin projection, local
mean transfer, low-arm bounds and weighted residual comparison with
separate contributing agents and shared inputs. The recorded
configurations for the B10 proof and the audit are unavailable.
The complete-field provider and Green transfer retain their inherited
review scopes.

The inherited Green identity and prime baseline retain their scoped
internal assessments \cite{LocalizedMetricsAssessment,PrimeBaselineAssessment}.
The conditioning and short-mollifier results retain the correction and
assessment \cite{OGCorrection,OGAssessment}; the boundary-log theorem and
assembly retain E05's acceptance with its additive evidence/control
clarification \cite{OGAssemblyAssessment}. E07 accepted the earlier
deflation and independent-cell comparison; the later contour design
retains its proof-author scope. Integration does not extend a review to
later statements that it did not examine.

The geometric constructions retain supplied-preparation scope. The
Riemann--Siegel passages retain their exact source versions and parameter
domains. Connecting derivations and this manuscript's production check
have same-assistant scope. There has been no independent specialist
human review or formal verification. The accompanying source map and
manifest identify the consumed files; hashes establish file identity,
not mathematical validity.


\section{Outstanding analytical tasks}\label{sec:outstanding}

The transition-scale consumer is selected, and its outer reduction and
adaptive-reference comparison costs are paid. The remaining obstacle is
a value estimate at exceptional centers. The density and averaged-value
conclusions, including the quantitative short-window bounds, leave
that obstacle open. The following tasks distinguish
the missing arithmetic information from further verification and changes
of scope.

\paragraph{1. Acquire new arithmetic information at exceptional centers.}
For the unchanged band \(\lambda=c/L_X\), \(1/2\le c\le1\), and
\(P(X)=\lfloor L_X^2\rfloor\), estimate the complete
\(R_{P(X)}(X,c)\) in \eqref{eq:exact-relative-residual} strongly
enough for \eqref{eq:adaptive-allowance} to lie below \(m_*\).
Alternatively estimate the actual signed nonlinear logarithmic
contribution, retaining its amplitude and phase dependence. A new input
must specify a proof route and a complete error budget; repeating the
residual's definition or the desired signed inequality is not that input.
The field error \(X^{-1/4+o(1)}\) is already paid and does not bound
the arithmetic residual. The short-window Green bound allows total
mass \(O_\rho(\ell L_U^{-1/3})\), which grows on the admissible
windows. A persistence or maximum-to-average estimate would have to
control the actual trajectory \(X\mapsto(1/L_X,X)\), including its
time derivative contribution where differentiable.

\paragraph{2. Preserve the proved limits of the tested methods.}
The exact quadratic-Jensen recipe fails at arbitrarily high centers;
even its diagnostic ideal cosine subtraction fails by at least
\(1/1000\). Every fixed Euler reference has unbounded residuals,
and its specified signed tangent fails uniformly over regulators.
These statements do not decide a growing-reference tangent, spatially
varying tangent heights or all nonlinear signed methods. Dividing by
the reflected reference would meet its real zeros and requires a
separate small-value argument. No such quotient is assumed zero-free.
Keep these tested methods parked unless the proposed change supplies
information outside their proved failure class.

\paragraph{3. Resolve uniformity only in the domain actually needed.}
The density, finite-power and subcritical exponential mean conclusions
apply at \(c=1\), with \(\lambda=1/L_X\) varying between centers.
Theorem~\ref{thm:b10-averaged} uses a cutoff fixed over each translated
interval. Corollary~\ref{cor:b10-green} gives a rate for the Green mean
and the fixed-threshold exceptional proportion with that cutoff.
The varying-cutoff residual, the extension throughout \(c\in[1/2,1]\),
and the shrinking-depth regime need separate uniform estimates.
The fourth exponential mean is bounded and its convergence remains
open. The asymptotic statements begin at a finite unevaluated onset.

\paragraph{4. Retain the separate time-zero arithmetic problem.}
At time zero the parent target remains an upper bound for
\(\int K\log|F|-P+G(X,1)\) below the chosen defect margin.
The increasing-order Riemann--Siegel remainder in
\eqref{eq:RS-log-transfer} controls its logarithmic error but not the
retained arithmetic value. The shifted-zeta identity
\eqref{eq:shifted-zeta-divisor} retains an unknown zero potential
and signed difference; it does not settle this parent question.
The localized-correlation and collision-constraint tracks likewise
retain their separate data requirements and ownership.

\paragraph{5. Check the foundation and make constants effective when useful.}
The imported complete-field estimate, its fixed buffers and Gaussian
exterior, the source normalization, opposite heat clocks, divisor
signs and logarithmic tails retain their original internal-review
status. The B10 averaged theorem has a separate internal audit with
contributing agents and shared inputs. Specialist review remains
outstanding. The finite provider constants and onsets are unevaluated.

\paragraph{6. Use numerics for a specified analytical question.}
The new results are analytical. Existing numerical certificates retain
their finite supporting scope in Section~\ref{sec:evidence}. A future
computation should calibrate a named constant, test a specified proposed
estimate or inspect a proved failure mechanism, with a fixed comparator
and budget. A height scan cannot establish the required every-center
bound. Further execution should follow a new arithmetic estimate or a
concrete proof question; no additional task is issued merely to rename
the remaining residual.

\begingroup\footnotesize\raggedright
\bibliographystyle{plainurl}
\bibliography{green_newman_sources}
\endgroup
\end{document}
